% !TeX root = ../main.tex
%
% Body-only import derived from papers/lop/main.tex.
% Standalone title, author block, abstract, acknowledgements, and bibliography
% commands have been omitted for thesis integration.

\section{Introduction}
One of the central problems in complexity theory is to prove strong lower bounds on the size of boolean circuits computing some explicit function $f: \B^n \to \B$; by a classical result of Shannon, most such functions $f$ require circuits of size $~ 2^n/n$. In the original formulation of this problem, ``explicit'' meant $f \in \mathsf{NP}$, since in this case case proving a lower bound $n^{\omega(1)}$ would separate the classes $\mathsf{P}$ and $\mathsf{NP}$. It was soon realized that this problem remains difficult even for much weaker definitions of ``explicit,'' i.e. if we broaden our search for high circuit complexity functions to a much larger uniform complexity class than $\mathsf{NP}$. 

Obviously there are decidable problems with maximal circuit complexity, since a Turing machine with sufficient resources may search for the hardest function $f_n: \B^n \to \B$ by brute force given $n$, and then compute that function $f_n$ on a given input of length $n$. Kannan \cite{Kannan} was the first to investigate the precise complexity upper bound of this brute-force construction of a hard language: he observed that the class $\mathsf{E}^{\Sigma_2^P}$, which denotes $2^{O(n)}$-time machines with access to an oracle for a $\mathsf{\Sigma_2^P}$ complete language, contains a function of maximum circuit complexity, precisely because this class possesses sufficient resources to diagonalize over all low-complexity functions of a given input length. The question of finding the smallest uniform complexity class containing an exponentially hard boolean function came to more prominence following the seminal works of Nisan, Wigderson and Impagliazzo\cite{NW,IW}, who showed that if one could improve Kannan's upper bound from $\mathsf{E}^{\Sigma_2^P} $ to $\mathsf{E}$, it would imply $\mathsf{BPP}=\mathsf{P}$, i.e. universal derandomization of all polynomial time algorithms.

A recent line of work on the \emph{Range Avoidance} problem (``Avoid'') has reformulated this topic in the language of search problems \cite{Korten,CHR}. Consider the so-called ``truth table generator'' $\mathrm{TT}_n: \B^{\ell} \to \B^{2^n}$ which takes as input the description of a circuit $C: \B^n \to \B$ of size $s$, and outputs the truth table of the function it computes. If $s << 2^n$, then a standard argument shows that we may encode $C$ in a direct way using $< 2^n$ bits, so that $\ell < 2^n$ and the function $\mathrm{TT}_n$ is computable in $2^{O(n)}$ time. Therefore, to produce a function of high circuit complexity, it suffices to find a string $y \in \B^{2^n}$ outside the range of $C$. This general problem, where we are given a circuit with more output bits then input bits and wish to find a string outside its range, is called Range Avoidance. By the pigeonhole principle, it is a \emph{total search problem}, meaning it always has solutions.
Viewing $\mathrm{TT}_n$ as an instance of the Range Avoidance problem we can observe that if Avoid has a polynomial time algorithm then $\mathsf{E}$ requires exponential-size circuits. For larger complexity classes, it was shown in \cite{Korten} (based on an earlier result in \cite{Jerabek}) that solving Avoid and proving circuit lower bounds are actually \emph{equivalent problems}: for exponential-time classes at least as large as $\mathsf{E^{NP}}$, proving an exponential circuit lower bound is equivalent to solving Avoid in the polynomial-time analogue of that class. Beyond constructing hard boolean functions, algorithms for Range Avoidance have a host of further applications in \emph{explicit constructions} of pseudorandom objects \cite{Korten,low-depth-avoid}.

This perspective was crucially used in recent breakthrough works of Chen, Hirahara, Ren and Li \cite{CHR,Li}, who gave the first unconditional improvement to Kannan's classical result\footnote{This is the first improvement on finding the smallest complexity class with an \emph{exponential} circuit lower bound. If the goal is to prove merely superpolynomial lower bounds, lower bounds for smaller classes can be established using Karp-Lipton theorems \cite{karp-lipton}, an approach pioneered in Kannan's original paper \cite{Kannan}. This method can at best show ``sub-half-exponential'' lower bounds, and only for infinitely many input lengths.}. These works showed that the classes $\mathsf{S_2^P}\subseteq \mathsf{ZPE^{NP}} \subseteq\mathsf{\Sigma_2^E} $ all contain a function of circuit complexity $2^n/n$, which is within a $(1 \pm o(1))$ factor of the maximum possible circuit complexity of a boolean function. The results are established by giving a new algorithm for Avoid which runs in the class $\mathsf{FS_2^P}$, the functional variant of the decision class $\mathsf{S_2^P}$. 
%\todo{Toni: should we add a sentence saying this is equivalent to giving small size depth-3 ac0 ckts for Avoid? Also maybe this last sentence would be better placed in next paragraph.}

The authors of \cite{CHR,Li} made progress on a classical lower bound problem by discovering a new kind of algorithm for the total search problem Range Avoidance. This search problem lies in an unusual complexity class which was only first investigated a few years ago by \cite{KKMP}: it is a member of $\mathsf{TF\Sigma_2^P}$, the class of total search problems in the \emph{second level} of the polynomial hierarchy. Since the introduction of $\mathsf{TF\Sigma_2^P}$ in \cite{KKMP}, there has been no follow up work developing a structural classification of the problems therein, despite the considerable attention that has been devoted to Range Avoidance \cite{Korten, CHR, time-space, hard-partial, io,karthik,hardness-lwe,low-depth-avoid} and explicit construction problems more generally. The work of \cite{CHR,Li} indicates that some problems in $ \mathsf{TF\Sigma_2^P}$ admit highly nontrivial algorithms, algorithms which have consequences in seemingly unrelated areas of complexity. Do such algorithms exist for all search problems in $\mathsf{TF\Sigma_2^P}$? What more can we say about the relation of Avoid to other problems in this class? The purpose of this work is to address these two questions in particular, and more broadly to further the systematic study of $\mathsf{TF\Sigma_2^P}$ beyond its introduction in \cite{KKMP}, with an emphasis both on \emph{inclusions} and \emph{black box separations}. Prior to our work, no non-trivial separations between problems in $\mathsf{TF\Sigma_2^P}$ were known.

% \todo{below old}
% The recent success of investigating circuit lower bounds through the lens of total search problems naturally motivates a broader study of the concrete complexity of total search problems of a similar kind \todo{rephrase. Maybe some thing like: Motivated by the recent success in proving circuit lower bounds for easier functions, in this work we initiate a broader study ..}. While the class of total $\mathsf{NP}$ search problems, denoted $\mathsf{TFNP}$, has been investigated thoroughly over the past three decades \todo{throw in some citations}, the Range Avoidance problem seems to be of a higher complexity, belonging to the class $\mathsf{TF\Sigma_2^P}$ of total search problems at the second level of the polynomial hierarchy. The class $\mathsf{TF\Sigma_2^P}$ was only introduced quite recently in \todo{KKMP} (the same paper which introduced the Range Avoidance problem), and so far very little is known about its structure. Do all $\mathsf{TF\Sigma_2^P}$ search problems admit the same nontrivial algorithms as were discovered for Range Avoidance in \todo{cite}? If not, then what structure exists for Range Avoidance which distinguishes it from other search problems in this class? In this work we give answers to both of these problems. 

$\mathsf{TF\Sigma_2^P}$ has a more famous cousin one level down in the polynomial hierarchy, the class $\mathsf{TFNP}$ of total $\mathsf{NP}$ search problems. In contrast to $\mathsf{TF\Sigma_2^P}$, $\mathsf{TFNP}$ has received a thorough investigation over the past three decades and its structure is rather well understood in comparison. 
%\todo{At the risk of making intro even longer, might be worth saying how important it has turned out to be -- connections to cryptography, very strong connections to proof complexity, and also helped to understand/catagorize some important search problems such as Nash.} 
A major distinction we reveal between $\mathsf{TF\Sigma_2^P}$ and $\mathsf{TFNP}$ is that the former has a plethora of \emph{resource-constained} subclasses, many of which can be separated from one another by explicit and natural search problems. By ``resource-constrained subclass'' we mean a class of search problems characterized by the existence of some resource-constrained algorithm which can find a solution. One example we will see is the class $\mathsf{psFZPP^{NP}}$, the class of problems where for every input $x$, there is a canonical solution $y_x$ which is output with high probability by some polynomial time $\mathsf{NP}$-oracle algorithm. Another class is $\mathsf{FP^{\Sigma_2^{NP} \cap \Pi_2^P}}$, consisting of those problems which can be solved in polynomial time with oracle access to a language in $\mathsf{\Sigma_2^P \cap \Pi_2^P}$. We will introduce and study several other such subclasses, and exhibit four natural search problems which exhibit separations (in the decision tree model) between them. This situation is in stark contrast to the known structure of $\mathsf{TFNP}$: while $\mathsf{FNP}$ does contain some intermediate resource-constrained classes, such as $\mathsf{psFZPP}$ and $\mathsf{FP^{NP \cap coNP}}$ (these are roughly analogous to $\mathsf{psFZPP^{NP}}$ and $\mathsf{FP^{\Sigma_2^{NP} \cap \Pi_2^P}}$ mentioned above), it is known that in the decision tree model all of these classes collapse to $\mathsf{FP}$. As a result in the decision tree model, the only resource-based distinctions amongst the standard $\mathsf{TFNP}$ problems is the distinction between $\mathsf{FP}$ and the rest of $\mathsf{TFNP}$.
%\todo{i reworded, since it sounded like FP and its complement are different.}
%i guess you mean distinction between FP and the other nontrivial TFNP classes. ?}

% We now describe the four $\mathsf{TF\Sigma_2^P}$ search problems which we focus on. \todo{explain that we will call ``avoid'' ``weak avoid'' from now on to differentiate it}
\subsection{Overview of Main Results}
A search problem is defined by a relation $R \subseteq \B^* \times \B^*$, where for each ``instance'' $x$ we say $y$ is a ``solution for $x$'' if $(x,y) \in R$; the relevant task is to find a solution given an instance. We say that a search problem is \emph{total} if every instance has a solution. A defining feature of search problems which distinguishes them from decision problems is that a given instance may have many different solutions. 
Indeed, if a search problem has a unique solution on every instance, then it may be equivalently phrased as a decision problem: given $(x,i)$, output the $i^{th}$ bit of the unique solution for $x$.
For many search problems of interest it is not clear how to reduce them to a decision problem of the same complexity. For example given total search problem $R$ we may define the decision problem $\mathrm{LexFirst}_R$, where given an instance $(x,i)$ we must output the $i^{th}$ bit of the lexicographically first solution to $x$. While $\mathrm{LexFirst}_R$ is clearly at least as hard as $R$, in many cases it will be much harder and a reduction does not seem to exist in the opposite direction. 


Say that $\mathcal{A}$ is a deterministic algorithm solving some search problem $R$. Then $\mathcal{A}$ naturally associates to each instance $x$ a \emph{canonical solution} $y_x:= \mathcal{A}(x)$ which it outputs. In particular we can think of $\mathcal{A}$ as \emph{defining a second search problem} $R'$ with $(x,y) \in R'$ iff $\mathcal{A}(x)=y$; if $\mathcal{A}$ solves the original search problem $R$, then $R$ is reducible to $R'$ since $(x,y)\in R' \rightarrow (x,y)\in R$. Now, if the algorithm $\mathcal{A}$ lies in some restricted complexity class, then this places the same complexity upper bound derived search problem $R'$. The point of this is that when we have some nontrivial deterministic algorithm solving a search problem $R$, we may think of it as giving a reduction from $R$ to a search problem with \emph{unique solutions}.

For most of the classical problems in the class $\mathsf{TFNP}$ it is believed that there is \emph{no decision problem} which captures their complexity precisely, a conjecture supported by black box separations \cite{TFNP-separations}. More generally it is believed that for most of the standard problems $R \in \mathsf{TFNP}$, any unique-solution problem $R'$ which we can reduce $R$ to must lie \emph{outside of} $\mathsf{TFNP}$\footnote{This is equivalent to saying that we believe these standard problems, e.g. $\mathsf{PPAD}$ or $\mathsf{PPP}$, not to lie in $\mathsf{FP^{NP \cap coNP}}$.}. The new results in \cite{Li,CHR} reveal that the situation for Range Avoid is different: because their Range Avoidance algorithms are deterministic and are upper bounded inside $\mathsf{TF\Sigma_2^P}$, they imply that Range Avoidance is reducible to a problem in $\mathsf{TF\Sigma_2^P}$ with unique solutions. To make this discussion more formal we need to define the class $\mathsf{TF\Sigma_2^P}$ and its \emph{unique solution} subclass $\mathsf{TFU\Sigma_2^P}$:

% \todo{for typical search problems, it is not clear how to reduce them to a decision problem without making the problem harder; indeed for $\mathsf{TFNP}$ it is believed that many of the standard problems are not equivalent to 
% \emph{any} decision problems, and this is supported by known black box separations [cite first TFNP separation paper]. An important feature of the new Avoid results is that they show Avoid is reducible to a total search problem with \emph{unique solutions}, which means that... we define this formally as follows}. To define these things formally we first present \todo{definition of TFSigma and sv}
\begin{definition}[$\mathsf{TF\Sigma_2^P}$ and $\mathsf{TFU\Sigma_2^P}$]
\label{def:TFUSigma2}
A polynomially-bounded\footnote{The ``polynomially bounded'' condition just means we are restricting attention to search problems whose solutions have polynomially bounded length.} search problem $R$ lies in $\mathsf{TF\Sigma_2^P}$ if it is a total search problem, and there exists a $\mathsf{coNP}$ verifier $\mathrm{V}$ so that $(x,y) \in R$ if and only if $\mathrm{V}(x,y)=1$. We say $R \in \mathsf{TFU\Sigma_2^P}$ if moreover every instance has a unique solution.
\end{definition}
In this terminology Li's result implies the following: Range Avoidance is polynomial-time reducible to a problem in $\mathsf{TFU\Sigma_2^P}$. His result is in fact a significant strengthening of this, but for now we focus on this particular consequence, which seems quite unintuitive on the surface: given an instance $f: \B^n \to \B^{n+1}$ of Range Avoidance, there are at least $2^{n}$ distinct solutions, and it is not clear how to narrow down to any particular solution which is ``more special'' then the others. Our first contribution is to clarify this result in the following way: we introduce a \emph{natural} $\mathsf{TF\Sigma_2^P}$ search problem, whose containment in $\mathsf{TFU\Sigma_2^P}$ is obvious from the definition, and then show that Range Avoidance reduces to this problem: 
\begin{definition}[Linear Ordering Principle (LOP)]\label{def:lop}
Given $\prec: \B^n \times \B^n \to \B$ specified by a boolean circuit, find a witness that $\prec$ does not define a total ordering on $\B^n$, or else find the minimal element in the ordering it defines. A witness that $\prec$ does not define a total order consist of $x,y,z \in \B^n$ such that one of the following holds: (a) $x \prec x$; (b) $x\neq y$, $x \not\prec y $, and $ y \not\prec x$; (c) $x \prec y \prec z$ and $x \not\prec z$.
\end{definition}

This problem was originally defined in the context of Bounded Arithmetic in \cite{Chiari1998-CHIWFI}. While this problem does not literally have unique solutions as stated, it has a property which we show is morally equivalent. Observe that one type of solution in this problem is easier to verify then the other: a witness that $\prec$ fails to define a linear order can be verified in polynomial time, whereas a candidate minimum element can only be verified in $\mathsf{coNP}$. Moreover, if an instance has no easily verifiable solutions, then it has a unique solution. This follows from the fact that every linear order has a unique minimal element. We summarize this by saying that LOP has \emph{essentially unique solutions}; a very easy argument (Section~\ref{sec:prelims}, Lemma~\ref{lem:unique-sol}) shows that any problem with essentially unique solutions is reducible to a problem in $\mathsf{TFU\Sigma_2^P}$. We then prove:

\begin{theorem}
\label{thm:avoid-to-lop-intro}
Range Avoidance is polynomial time reducible to LOP.
\end{theorem}

We believe this result goes a long way in explaining the new upper bounds for Range Avoidance. In particular our reduction isolates the two key steps in Li's algorithm which allow us to single out a special \emph{low complexity canonical solution} for Range Avoidance: the first step prepares a special subset of solutions using a certain tree-like iteration construction, and the second step singles out a fixed canonical solution among these by defining a certain \emph{total ordering} on these special solutions.  Recall our comment under Definition \ref{def:TFUSigma2} that Li's result is in fact stronger than a reduction from Range Avoidance to a problem in $\mathsf{TFU\Sigma_2^P}$: more generally he shows that Range Avoidance lies in the complexity class $\mathsf{FS_2^P}$ (we will review the definition of this class in Section ~\ref{sec:prelims}). Our result subsumes this upper bound as well, since another very direct argument shows that LOP lies in $\mathsf{FS_2^P}$. Thus our result gives the current best upper bound on the Range Avoidance problem, and hence on the Kannan's classical problem of constructing the truth table of a hard boolean function.

At this point we have seen that there are search problems in $\mathsf{TF\Sigma_2^P}$ which seem on the surface to have no distinguished solutions, but which nonetheless can be reduced to problems in $\mathsf{TFU\Sigma_2^P}$ by some highly non-obvious means. This naturally points to the following question: are all problems in $\mathsf{TF\Sigma_2^P}$ reducible to $\mathsf{TFU\Sigma_2^P}$? We give a negative answer in the decision tree model. Our separation is exhibited by the following relative of Range Avoidance which was introduced originally in \cite{KKMP}.
\begin{definition}[Strong Avoid]
Given $f: \B^n \setminus \{0^n\} \to \B^n$, find $y \in \B^n \setminus \rng(f)$.
\end{definition}
From now on we will refer to Range Avoidance as ``Weak Avoid'' to distinguish it from Strong Avoid. Our black box separation will be significantly stronger than just showing that Strong Avoid is not reducible to a problem in $\mathsf{TFU\Sigma_2^P}$; we will show more generally that Strong Avoid cannot be solved by making non-adaptive queries to any language in $\mathsf{\Sigma_2^P}$, which is equivalent to proving size lower bounds for depth-3 $\mathsf{AC}_0$ circuits solving Strong Avoid: 

\begin{theorem}
\label{thm:strongavoid-lb}
%Strong Avoid requires exponential-size depth-3 $\mathsf{AC}^0$ circuits. Thus
In the decision tree model, Strong Avoid is not in $\mathsf{FP_{||}^{\Sigma_2^P}}$. More specifically, let ${C:\B^{(N-1) \log N } \to \B^{N^{\epsilon}}}$ be a depth-3 circuit of size $2^{N^{\epsilon}}$  and let $D: \B^{N^{\epsilon}} \to \B^{ \log N}$ be an arbitrary postprocessing function, where $N=2^n = |\B^n|$. Then provided $\epsilon$ is sufficiently small, $D \circ C$ cannot solve Strong Avoid: there must be some input $f: [N-1] \to [N]$ so that $D(C(f))$ fails to find a $y \notin \rng(f)$.
\end{theorem}

This immediately implies non-reducibility to $\mathsf{TFU\Sigma_2^P}$, since any problem $R \in \mathsf{TFU\Sigma_2^P}$ with unique solutions can be solved with non-adaptive queries to the language $\{(x,i) \mid (x,y) \in R \rightarrow y_i=1\}$ which lies in $\mathsf{\Sigma_2^P}$. The connection of these kinds of separations to depth 3 $\mathsf{AC}^0$ lower bounds was spelled out in a recent paper of Vyas and Williams \cite{vyas-williams} for the case of \emph{Weak} Avoid: their work established that nontrivial upper bounds for Weak Avoid are equivalent to certain depth 3 circuits solving the so-called ``Missing String'' problem: given an explicit list of $2^{n-1}$ $n$-bit strings, output a string not in the list. The input size here is $\approx N=2^n$, and the question is whether a depth 3 circuit exists of size polynomial or quasipolynomial in $N$. Note that the Missing String problem is simply the black-box variant of Weak Avoid. Li's result showed that quasipolynomial size depth-3 circuits for this problem actually \emph{do exist}, solving the original question of Vyas and Williams in the positive. We show that if the problem is modified so that the list of strings has length $N-1$ rather then $\frac{N}{2}$, then depth-3 circuits require exponential size to solve this problem. We note that our lower bound holds against a stronger class of circuits than what was originally considered by Williams and Vyas: in their model the depth 3 circuit is of the form $ C:\B^{(N-1) \log N } \to \B^{ \log N}$ and must output the exact solution to Avoid. Here we allow $C$ to output an arbitrary string in $\B^{N^{\epsilon}}$, which can then be postprocessed arbitrarily to construct a solution to Avoid.

By a simple reduction, we also obtain quasipolynomial depth-3 circuit size lower bounds even for \emph{moderately weak} Avoid instances with domain $[N]$ and codomain $[N + N/\log^{O(1)} N]$. This result gives a complete characterization of the degree of ``Weakness'' necessary to obtain quasipolynomial size depth 3 circuits: if the codomain has size $N + N/\log^{O(1)} N$ then circuits of quasipolynomial size suffice, and if its size is $N + N/\log^{\omega(1)}N$ then they do not; see Lemma~\ref{lem:php-lb-larger-stretch} for details.
%which we explain in \todo{ref section}.

Our last main result exhibits a more fine-grained separation amongst the subclasses of $\mathsf{TF\Sigma_2^P}$. Above we have highlighted one important distinction, between the problems which are reducible to $\mathsf{TFU\Sigma_2^P}$ and those which aren't. However both of the problems we've seen so far which reduce to $\mathsf{TFU\Sigma_2^P}$ also have an additional property: they are solvable by a polynomial time randomized algorithm using a $\mathsf{NP}$ oracle. For Weak Avoid this follows from its definition, while for LOP it follows from its containment in the class $\mathsf{FS_2^P}$. We show that this is not possible for all problems reducible to $\mathsf{TFU\Sigma_2^P}$. Our separation is exhibited by the following natural search problem: 
\begin{definition}[Strong 1-1 Avoid]
Given $f: \B^n \setminus \{0^n\} \to \B^n$, find a pair $x\neq y $ in $\B^n \setminus \{0^n\}$ such that $f(x)=f(y)$, or else $y \in \B^n \setminus \rng(f)$, 
\end{definition}
Observe that Strong 1-1 Avoid enjoys the same property as LOP of having \emph{essentially unique solutions}: it is easy to verify the collision solutions $f(x)=f(y) \land x\neq y$, and any instance with no collision solutions has a unique solution. This follows from the fact that any injective function $f: [N] \to [N+1]$ misses exactly one point in its codomain. Hence, like Weak Avoid and LOP, the problem Strong 1-1 Avoid is reducible to $\mathsf{TFU\Sigma_2^P}$. However we prove the following lower bound:
%placing Strong 1-1 Avoid outside of $\mathsf{FBPP^{NP}}$:

\begin{theorem}
\label{thm:bpp-np-lb}
In the decision tree model, Strong 1-1 Avoid is not solvable in $\mathsf{FBPP^{NP}}$.
%\begin{theorem}\label{thm:bpp-np-lb}
%More generally, there is an absolute constant $\epsilon > 0$ so the following holds: if $T: \B^{(N-1} \lceil \log N \rceil \rightarrow  $ is a $\mathsf{P^{NP}}$ decision tree of depth $N^{\epsilon}$ over the input variables to Strong 1-1 Avoid ($\mathrm{BITS}_{N,N+1}$ and leaves labeled by elements of $[N+1]$, then:
%\[
%\Pr_{f \sim \mathcal{F}_N} [T(f) \notin %\rng(f)] \leq N^{-\Omega(1)}
%\]
\end{theorem}

Aside from revealing further the structure of $\mathsf{TF\Sigma_2^P}$, this result yeilds a new separation for \emph{decision classes} which was not previously known:
\begin{theorem}
\label{thm:sigma2-pi2-separation}
In the decision tree model, $\mathsf{\Sigma_2^P \cap \Pi_2^P} \not\subseteq \mathsf{BPP^{NP}}$. In particular $\mathsf{\Sigma_2^P \cap \Pi_2^P} \not\subseteq \mathsf{S_2^P}$.
\end{theorem}
% The class $\mathsf{UTF\Sigma_2^P}$ occurs under the name $\mathsf{svTF\Sigma_2^P}$ in \todo{ref, maybe explain why this is the same?}.
In the next two subsections we will describe these main results in some more technical detail. A diagram of the structure of $\mathsf{TF\Sigma_2^P}$ and our main results is given in Figure~\ref{fig:classes}; some classes in this diagram will not be defined until Section~ \ref{sec:totalsearch}.
\begin{figure}
    \centering
    \includegraphics[width= \textwidth]{Classes.pdf}
    % \captionsetup{width=1\textwidth}
    \caption{Inclusion diagram of relevant classes and search problems. Solid arrows represent the inclusion of the class at the base of the arrow into the class at its tip. Dotted arrows indicate non-inclusion of the base class into the tip class in the decision tree model. The main linear axis of classes along the top are all included from left to right, indicated by the long solid arrow above them. Our major results are the separations marked with $\spadesuit, \clubsuit$ and the inclusion marked with $\blacklozenge$. $\mathsf{TFU\Sigma_2^P}$ essentially corresponds to $\mathsf{FP^{\Sigma_2^P \cap \Pi_2^P}}$, see Lemma~\ref{lem:unique-intersect}.}
    \label{fig:classes}
\end{figure}
\subsection{$\mathsf{AC^0}$ Lower Bounds and Class Separations}
Our two main separations show that Strong Avoid has no non-trivial upper bound inside of $\mathsf{TF\Sigma_2^P}$ (Theorem~\ref{thm:strongavoid-lb}), and that Strong 1-1 Avoid has no randomized $\mathsf{NP}$-oracle algorithm (Theorem ~\ref{thm:1.5-lb}). The first lower bound is the more involved of the two, and requires proving a novel depth 3 $AC^0$ circuit lower bound for a total search problem, which appears to be the first circuit lower bound of this kind. 

%We now give the formal statement of Theorem ~\ref{thm:strongavoid-lb}:
%is is equivalent to the following circuit lower bound:

%\bigskip

%\begin{theorem} 
%\noindent {\bf Theorem ~\ref{thm:strongavoid-lb}.}
%\label{thm:depth3-avoid-lb-intro}
%{\it Consider the natural association of $\B^{(N-1)\lceil \log N \rceil}$ with functions $f: [N-1] \to [N]$ specified by their truth table, where $\alpha \in \B^{(N-1)\lceil \log N \rceil}$ is interpreted as an ordered list of the values of $f$ on each point in $[N-1]$. Say that $C:\B^{(N-1)\lceil \log N \rceil} \to \B^{\lceil \log N \rceil}$ is a depth 3 unbounded fanin boolean circuit which solves the Strong Avoid problem in the sense that $C(f) \notin \rng(f)$ for all assignments $f$. Then $|C| \geq 2^{N^{\Omega(1)}}$. 

%More generally, if $C:\B^{(N-1)\lceil \log N \rceil} \to \B^{N^{\epsilon}}$ is a depth 3 circuit of size $2^{N^{\epsilon}}$  and $D: \B^{N^{\epsilon}} \to \B^{\lceil \log N \rceil}$ is an arbitrary postprocessing function, then $D \circ C$ cannot solve Strong Avoid provided $\epsilon$ is sufficiently small.}

%\todo{(Toni)This last sentence is a bit hard to understand. Maybe something like: Moreover, the lower bound holds even if $C$ can output $n^{\epsilon}$ output bits which can be postprocessed arbitrarily to solve Strong Avoid.  }\todo{(Oliver) is this better? Also is it ok for me to still hint at the "non-adaptive decision tree" formulation as below or do you think its still too confusing?}
%\end{theorem} 

Theorem~\ref{thm:strongavoid-lb} yields a very fine-grained separation for the Strong Avoid problem. It is easy to construct a depth 4 circuit $C:\B^{(N-1)\lceil \log N \rceil} \to \B^{\lceil \log N \rceil}$ solving Strong Avoid, where moreover the bottom fan-in is only $O(\log N)$ \cite{vyas-williams}. More strongly, it is possible to construct a depth $O(\log N)$ decision tree, which at each step queries a depth 3, $\poly(N)$ size circuit on the input $f$, and at each leaf outputs a correct solution $y \notin \rng(f)$. Our lower bound can be interpretted as saying that if such a decision tree is forced to be non-adaptive, then either the circuits it queries at each step must grow to exponential size, or else the number of queries must grow to $N^{\Omega(1)}$. This also contrasts the situation with Weak Avoid, where as mentioned above,  Li's construction gives depth-3 $\mathsf{AC}_0$ circuits of size $N^{O(\log N)}$ for solving Weak Avoid.

%\begin{theorem}\todo{cite}
%There is a depth 3 circuit $C:\B^{N\lceil \log 2N \rceil} \to \B^{\lceil \log 2N \rceil}$ of size $N^{O(\log N)}$ which solves the Weak Avoid problem: for every $f: [N] \to [2N]$, $C(f) \notin \rng(f)$. Moreover the bottom layer of gates in $C$ have fan-in $O(\log^2 N)$.
%\end{theorem}

%\todo{TONI: consolidate} 
We believe our lower bound is of independent interest in %bounded depth 
circuit complexity. In particular, we give a very precise depth-3 lower bound for a \emph{total search problem}. It is of course possible to construct contrived examples of total search problems which are hard for $\mathsf{AC}^0$ circuits, for example ``given $x \in \B^n$ output $b \in \B$ such that $\mathrm{Parity}(x)=b$;'' in such examples we can derive hardness of the search problem \emph{by reduction from a decision problem}. In contrast, the lower bound we show here \emph{cannot be established by reduction from any decision problem}. This follows from the fact that any decision problem which is reducible to a $\mathsf{TF\Sigma_2^P}$ search problem lies in $\mathsf{\Sigma_2^P \cap \Pi_2^P}$. In the decision tree model, this means that any language which can be solved by a small depth decision tree querying instances of Strong Avoid has both $\Sigma_2^P$ and $\Pi_2^P$ circuits. However, the lower bound we are trying to show for Strong Avoid rules out the existence of any $\mathsf{FP^{\Sigma_2^P \cap \Pi_2^P}}$ algorithm for Strong Avoid. Phrased more succinctly, our lower bound establishes that Strong Avoid is harder then \emph{any} decision problem which can be reduced to it, which by definition means we cannot establish the lower bound itself by reduction from a decision problem. In light of this, to prove Theorem \ref{thm:strongavoid-lb} we must develop new $\mathsf{AC^0}$ lower bound techniques which are specially tailored to total search problems.

% \todo{add a pointer to the section where we compare our switching lemma with those in proof complexity}
%\todo{Toni: should we add somewhere that we also get a fine grained separation with respect to the stretch of Avoid since by a simple reduction we get a superpoly depth-3 ckt lb also for avoid from [N] to [N + N/polylogn]?}

Our second main lower bound (Theorem ~\ref{thm:bpp-np-lb}) places Strong 1-1 Avoid outside of $\mathsf{FBPP^{NP}}$. 
%This can be equivalently viewed as giving strong lower bounds for randomized decision trees for Strong 1-1 Avoid:
%\medskip
%\begin{theorem}
%\noindent{\bf Theorem ~\ref{thm:bpp-np-lb} (Restatement).}
%{\it Associate $\B^{(N-1)\lceil \log N \rceil}$ with functions $f: [N-1] \to [N]$. Let $\mathcal{C}$ be a distribution over functions $C: \B^{(N-1) \lceil \log N \rceil} \to \B^{\lceil \log N \rceil}$, where each $C \in \mathrm{support}(\mathcal{C})$ is a $\mathsf{P^{NP}}$ decision tree of complexity $w$, and such that for all one-to-one assignments $f: [N-1]\to[N]$, 
%\[
%\Pr[\mathcal{C}(f) \notin \rng(f)] \geq %\frac{2}{3}
%\]
%Then $w \geq N^{\Omega(1)}$.}
%\end{theorem}
%\medskip
%Informally, $\mathsf{FBPP^{NP}}$ correspond to randomized decision trees of small height which, instead of querying variables, are allowed to query the value of an arbitrary low-width DNF over the inputs. 
%By Yao's minimax principle, it suffices to prove that any low-depth $\mathsf{P^{NP}}$ decision tree cannot solve Avoid with probability 2/3, with respect to the uniform distribution of 1-1 functions.
%any $\mathsf{P^{NP}}$ for
%We will review the formal definition of $\mathsf{P^{NP}}$ decision trees in Section \ref{sec:prelims} (Definition ~\ref{def:PNP-trees}). 
%Informally, these are decision trees of small height which, instead of querying variables, are allowed to query the value of an arbitrary low-width DNF over the inputs in unit time. 
Recall that Strong 1-1 Avoid has the property of having \emph{essentially unique solutions} and is thus reducible to a problem in $\mathsf{TFU\Sigma_2^P}$. We will see (Lemma ~\ref{lem:unique-intersect}) that this means it is reducible to a \emph{decision problem} in $\mathsf{\Sigma_2^P \cap \Pi_2^P}$. Combining this with the above lower bound we obtain Theorem ~\ref{thm:sigma2-pi2-separation}, which separates the decision tree class 
$\mathsf{\Sigma_2^P \cap \Pi_2^P}$ from $\mathsf{BPP^{NP}}$ and in particular from $\mathsf{S_2^P}$.
%the following novel separation for decision classes:
%\begin{theorem}
%In the decision tree model, $\mathsf{\Sigma_2^P \cap \Pi_2^P} \not\subseteq \mathsf{BPP^{NP}}$. In particular $\mathsf{\Sigma_2^P \cap \Pi_2^P} \not\subseteq \mathsf{S_2^P}$.
%\end{theorem}
This is in contrast to the situation for $\mathsf{NP} \cap \mathsf{coNP}$, which is known to collapse to $\mathsf{P}$ in the decision tree model. This improves a previous result of Fortnow and Yamakami \cite{forntnow-yamakami} who showed that $\mathsf{\Sigma_2^P \cap \Pi_2^P} \not\subseteq \mathsf{P^{NP}}$ in the decision tree model.

\subsubsection{Lower Bound Methods}
\label{sec:intro-proofoverview} 
%\noindent{\bf Lower Bound Methods.} 
%\todo{Oliver: some of this doesn't make sense because we haven't defined the relevant pigeonhole concepts yet}
A main technical ingredient in both lower bounds is a new Switching Lemma (Lemma ~\ref{lem:switching-main})  specialized for Avoid. 
%We use a coding style argument first introduced by Razborov \todo{raz} (see also \todo{beame}, \todo{add references}). 
Switching lemmas have been used for both circuit lower bounds for computing functions, and in proof complexity to prove lower bounds on the size of proofs of hard tautologies (e.g., \cite{beame-sl,razborov-sl,BeameIKPPW92}). Between the two, our argument bears a stronger resemblance to the second, however there are some key conceptual differences. 
The basic idea behind all switching lemmas is to show that under a random restriction $\rho$ (from a suitable distribution) a low-width DNF is likely to be represented by a low-depth decision tree. Since a low-depth decision tree representation for a function $f$ implies that both $f$ as well as its negation can be represented by low-width DNFs, this in turn allows us to collapse an AND of low-width DNFs into a single low-width DNF, thus reducing the circuit depth by 1.

A major difference between various switching lemmas is the choice of  distribution over restrictions, and the way in which the decision tree {\it represents} a DNF. In the original Switching Lemmas used to prove $\mathsf{AC}_0$ lower bounds for parity, the restrictions are simply uniformly random partial restrictions, and the notion of represents is with respect to {\it every} input. That is, the decision tree computes the same function as the DNF.

%In the case of switching lemmas used in proof complexity to prove $\mathsf{AC^0}$-Frege lower bounds for the pigeonhole principle, we think of $[N]$ as a ``nonstandard'' model, and thus we can view the input as specifying a 1-1 function from $[N]$ to $[N]$ (rather than from $[N]$ to $[N+1]$. The chosen distribution over restrictions are partial 1-1 matchings from $[N]$ to $[N]$, and a low-depth ``pigeonhole'' decision tree in this context can make queries to a pigeon or to a hole at each vertex, and every path in the tree corresponds to a partial matching. Then a pigeonhole decision tree represents a DNF if it is truth-functionally equivalent to the DNF with respect to all 1-1 functions from $[N]$ to $[N]$.

% \todo{here is an attempt at rephrasing}
In the case of switching lemmas used in proof complexity to prove $\mathsf{AC^0}$-Frege lower bounds for the $[N] \to [N+1]$ pigeonhole principle, we think of the input as specifying a purported 1-1 function from $[N]$ to $[N+1]$ (which cannot actually 
exist if $N$ is finite).  
%clearly such an assignment cannot actually exist. 
The chosen distribution over restrictions are partial 1-1 matchings from $[N]$ to $[N+1]$, and a low-depth ``PHP decision tree'' in this context can make queries to a pigeon or to a hole at each vertex, and every path in the tree corresponds to a partial matching. Since, in reality, no 1-1 function from $[N]$ to $[N+1]$ exists, these trees do not represent the original DNF in any standard way. However, if the input variables instead corresponded to a total 1-1 assignment from $[N]$ to $[N]$ (which do exist), then we can apply the same PHP Switching Lemma to prove that under a random partial 1-1 restriction, a low-width DNF is likely to convert to a low-depth PHP decision tree, which now represents the DNF in the sense that it agrees with the DNF on all input assignments that correspond to 1-1 mappings from $[N]$ to $[N]$. 

%we can use the PHP switching lemma  we only care about paths in the decision tree which follow partial 1-1 maps from $[N]$ to $[N+1]$. Since, in reality, no 1-1 functions $[N] \to [N+1]$ actually exist, these decision trees do not actually have well defined behavior on all true assignment $f:[N] \to [N+1]$.
% \todo{END attempt}

In the case of Avoid, we have to modify the way of constructing a decision tree associated with a DNF so that the decision tree represents the original DNF in the sense that they are truth-functionally equivalent with respect to all 1-1 input functions from $[N]$ to $[M]$, where now $M$ is strictly larger than $N$. 
%The largest differences  is the way in which the decision tree represents the original DNF.  
%that we require our pigeonhole decision trees to be equivalent to the hole disjunction that they represent, over all 1-1 input functions. In contrast, the pigeonhole decision trees used in proof complexity switching lemmas make queries to a pigeon or a hole at each vertex, and every path in the tree corresponds to a partial matching; because of this, if the number of pigeons and holes is unbalanced, the behavior of such trees is not actually well-defined on every total assignment $f: [N] \to [M]$. 
To achieve this, we modify the notion of pigeonhole decision trees as follows. As in the original PHP Switching Lemma in proof complexity, each node of our decision tree will query either a pigeon or a hole. When a pigeon is queried at a node, we allow edges for all possible holes that it could be mapped to. But when a hole is queried, now we have to allow for the possibility that this hole is unmapped: in addition to allowing edges for each pigeon that could map to this hole, we allow an extra edge corresponding to the case where nothing maps to this hole. With this modification, our pigeonhole decision trees will represent the original DNF with respect to all 1-1 inputs from $[N]$ to $[M]$. 
%have $N+1$ possible outgoing edges, one for each pigeon that could map to this hole, plus an extra edge corresponding to the case where nothing maps to this hole. 

Another crucial distinction is how we use the Switching Lemma to reduce the depth of the circuit by one. In the proof complexity setting for the pigeonhole principle lower bounds, we think of $N$ as infinite, and therefore with respect to 1-1 inputs, a DNF can be written as a low-width {\it matching disjunction}, where each term in the disjunction corresponds to a partial 1-1 function (or matching) from $[N]$ to $[N]$. After applying the PHP Switching Lemma and a union bound, each matching disjunction $f$ (under $\rho$) becomes a low-depth ``matching decision tree'', enabling a reduction in the overall circuit depth by one. To summarize, in the classical PHP Switching Lemma, the underlying depth-2 subcircuits are always low-width matching disjunctions, both before and after each application of the PHP Switching Lemma. 
%Underlying the classical switching lemma used in proof complexity shows that with high probability, a low-width matching restriction converts to a ``matching decision tree'' under a random matching restriction $\rho$. This in turn allows one to represent the original matching disjunction as well as its negation by a low-width matching disjunction. 
%such that  we are able to flip a $\lor\land$ formula of some type to a $\land\lor$ formula of the same type; in 
%particular if we start with a matching disjunction we can flip it to (the compliment of) another matching disjunction. 
In our case, the underlying depth-2 subcircuits are not of the same type before and after applying our Pigeonhole Switching Lemma; a consequence is that our Switching Lemma \emph{cannot be applied twice}. This is not a defect of our method, but rather a necessary feature of any technique here, since our search problem \emph{can be solved} by circuits of one higher depth. More specifically, we show that initially the bottom depth-2 subcircuits of ${\cal C}$ can be expressed as low-width matching disjuncdtions. But after applying the Pigeonhole Switching Lemma, and subsequent depth reduction, the new depth-2 subcircuits become \emph{hole disjunctions} which are a generalization of matching disjunctions, where now each term in the DNF can specify not only a partial 1-1 matching, but also a subset of holes that are unmapped.
%In our case we are only able to flip a matching to the compliment of a \emph{hole disjunction}, and for this reason our lemma cannot be applied twice. This is not a defect of our method, but rather a necessary feature of any technique here, since our search problem \emph{can be solved} by circuits of one higher depth. 
With these appropriate modifications, our proof of the switching lemma is similar to previous proofs.

Equipped with the Pigeonhole Switching Lemma, we can give the high level view of both proofs. 
We start with Theorem ~\ref{thm:strongavoid-lb}.
To prove Theorem ~\ref{thm:strongavoid-lb}, we would like to restrict  attention to the class of 1-1 input functions from $[N]$ to $[N+1]$. However if we truly restrict ourselves to 1-1 functions,
%$f \in {\cal F}_{N,N+1}$, 
then a lower bound is not possible by Lemma~\ref{lem:1-1-lop-unique}: for these inputs, there is a unique solution, and therefore a polynomial-size depth-3 circuits can easily check whether the unique solution has its $i^{th}$ bit equal to 1 or to 0 and hence solve the problem unconditionally.
To circumvent this barrier, we will prove a strengthening of Theorem~\ref{thm:strongavoid-lb}, by giving a lower bound for Avoid on input functions $f: [N] \rightarrow [M]$, where $M$ is larger than $N+1$. By enlarging the range of $f$, we can focus our attention of $f$'s that are 1-1 since now for every input, there are at least $M-N$ distinct solutions.
Note that this implies the lower bound stated above (with $M=N+1$), since there is a direct reduction from Avoid on instances $[N] \to [M]$ to instances $[N] \to [N+1]$: we simply map every element of $[M] \setminus [N+1]$ to the element $N+1$. Observe that this reduction does not preserve injectivity.

Now assume there exists small-size $s$ depth-3 circuit ${\cal C}$ computing Strong Avoid on 1-1 functions from $[N]$ to $[M]$, $M>>N$. 
We can first apply a standard argument (the Width Reduction Lemma ~\ref{lem:width-red}) so that we can assume that the bottom-level fanin of ${\cal C}$ is at most $O(\log s)$. After this step, we can assume that ${\cal C}$ is a size-$s$, depth-3 circuit, where the bottom depth-2 subcircuits are low-width matching disjunctions. 
Next we apply our Pigeonhole Switching Lemma (as discussed above) which will guarantee that there exists a matching restriction $\rho$ such that under $\rho$, all depth-2 matching disjunctions in ${\cal C}$ will convert to low-depth pigeonhole decision trees. This will allow us to reduce the overall circuit depth by 1, and afterwards each output bit of ${\cal C}$ will be computed by a low-width \emph{hole disjunction}.
As discussed above, a hole disjunction is a type of DNF that generalizes matching disjunctions: each term $t$ in the hole disjunction can be viewed as partial information about the input $f$. The partial information consists of two parts: (i) first, $t_1$ specifies a small partial matching, pairing up some pigeons in $[N]$ to some holes in $[N+1]$; (ii) secondly, $t_2$ specifies a small set of holes (disjoint from the holes mentioned in $t_1$ that are not in the range of $f$).
It remains to prove a lower bound for circuits ${\cal C}$ for solving Avoid, where each output bit is specified by a low-width hole disjunction. This is also accomplished using a kind of restriction, but rather then choosing it at random we apply a careful deterministic process involving a novel covering argument. This step is somewhat reminiscent of early proofs of the Switching Lemma (e.g., \cite{FSS,Ajtai83}).


The proof of Theorem ~\ref{thm:bpp-np-lb} again uses the Pigeonhole Switching Lemma, together with a direct argument. 
We want to prove depth lower bounds for $\mathsf{FBPP^{NP}}$, which informally are randomized decision trees of small height which, instead of querying variables, are allowed to query the value of an arbitrary low-width DNF over the inputs. 
By Yao's minimax principle, it suffices to prove that any low-depth $\mathsf{P^{NP}}$ decision tree cannot solve Avoid with probability 2/3, with respect to the uniform distribution of 1-1 functions. We think of this distribution in the following way: first sample a uniform \emph{partial} 1-1 assignment $\rho$, then sample a uniform extension of $\rho$ to a total assignment.
Applying our Pigeonhole Switching Lemma and a union bound, we can argue that with high probability over the first choice of $\rho$, all of the $\mathsf{NP}$ queries in our $\mathsf{P^{NP}}$ decision tree $T$ can be simplified to small depth pigeonhole decision trees, which overall allows $T$ to be replaced by a low depth pigeonhole decision tree. It then remains only to argue that a pigeonhole decision tree of low depth cannot solve 1-1 Strong Avoid with non-trivial probability on a uniform extension $f$ of $\rho$, which can be accomplished with a direct argument.

\subsection{Linear Ordering Principle}
\label{sec:intro-lop} 
We now discuss in more detail our results on the newly defined Linear Ordering Principle problem, abbreviated LOP. Recall that our main result here is Theorem ~\ref{thm:avoid-to-lop-intro}, which says that Weak Avoid is polynomial time reducible to LOP. The proof follows much of the high level structure of Li's result placing Weak Avoid in $\mathsf{FS_2^P}$, with some key modifications. Roughly speaking, Li's proof shows that given an instance of Avoid $f: \B^n \to \B^{n+1}$, we can define a comparison relation $\sqsubset$ on $\B^{\poly(n)}$, so that for some unique distinguished element $\pi^* \in\B^{\poly(n)}$ we have $\pi^* \sqsubset \pi$ for all $\pi \neq \pi^*$, and $\pi^*$ contains a solution to the original Avoid instance. In our case we need to define a similar comparison relation, which in addition \emph{globally} acts as a total order on $\B^{\poly(n)}$. To explain the argument more clearly we split the reduction into two parts. We first introduce an intermediate search problem called Forest Termination and reduce Weak Avoid to this problem, then we reduce Forest Termination to LOP.
We note that our proof, as well as Li's, also bears a strong resemblance to the work of \cite{Paris-Wilkie-Woods} who gave the first proof of the weak pigeonhole principle in the bounded arithmetic theory $T_2$.%\todo{Oliver: I changed this back to $T_2$ because I realized $S_2$ reads just like $\mathsf{S_2^P}$ and I don't want to confuse people... iss that ok?} \todo{Toni: may want to talk up this result more?}


% \begin{theorem}
% ~
% \begin{enumerate}
%     \item Any search problem which has a $\mathsf{P^{NP}}$ Turing reduction to $\mathrm{LOP}$ also has a polynomial time many-one reduction to $\mathrm{LOP}$.
%     \item There is a language $\mathrm{DLOP} \subseteq \B^*$ (``Decisional LOP''), so that $\mathrm{LOP}$ is truth table reducible to $\mathrm{DLOP}$ and $\mathrm{DLOP}$ is many-one reducible to $\mathrm{LOP}$.
% \end{enumerate}
% \end{theorem}
% The advantage is that unlike $\mathsf{FS_2^P}$, which has a purely \emph{semantic} definition and therefore seems unlikely to contain a complete problem, LOP is defined by a very natural total search problem. 
Our subsequent results show some appealing structural properties of the complexity class defined by reducibility to LOP. We start by proving closure under a broad class of reductions:
\begin{theorem}\label{thm:l2p-closure-intro}
Any search problem which has a polynomial time $\mathsf{P^{NP}}$ Turing reduction to LOP also has a polynomial time many-one reduction to LOP.
\end{theorem}
Combining this closure property with the fact that LOP has essentially unique solutions, we are able to conclude that LOP is \emph{equivalent} in complexity to a decision problem. In particular we can define a decisional complexity class $\mathsf{L_2^P}$ for which LOP is the ``complete problem'' (despite being a search problem and not a language). We start by presenting a machine-based definition of $\mathsf{L_2^P}$:
\begin{definition}\label{def:l2p}
A language $L$ is in the complexity class $\mathsf{L_2^P}$ if there is a polynomial time relation $R: (\B^*)^3 \to \B$ and a polynomial $p$, so that for all $x$, $R(x, \cdot, \cdot)$ defines a total order on $\B^{p(n)}$ whose minimal element $a$ has $a_1 = L(x)$. 
\end{definition}

We then have the following equivalent characterizations:
\begin{theorem}\label{thm:l2p-complete-intro}
The following are equivalent for a language $L$:
\begin{enumerate}
    \item $L \in \mathsf{L_2^P}$
    \item $L$ is polynomial time many-one reducible to LOP.
    \item $L$ is $\mathsf{P^{NP}}
$-Turing reducible to LOP.
\end{enumerate}
Conversely, the search problem LOP is polynomial time truth table reducible\footnote{This is essentially the most restrictive reduction possible when reducing a search problem to a language.} to a language in $\mathsf{L_2^P}$.
\end{theorem}
We mentioned in passing before that LOP is easily shown to lie in the class $\mathsf{FS_2^P}$, the functional analogue of the decision class $\mathsf{S_2^P}$ (this will be shown in Lemma ~\ref{lem:lop-in-s2p}. The same reason shows that $\mathsf{L_2^P} \subseteq \mathsf{S_2^P}$. For those unfamiliar with the somewhat unconvential class $\mathsf{S_2^P}$, this is a complexity class introduced independently by Russell-Sundaram and Canetti \cite{russell-sundaram,canetti}. Their goal was to identify the smallest class in the polynomial hierarchy which is sufficient to capture randomized algorithms, in particular $\mathsf{BPP}$ and $\mathsf{MA}$. Beyond this purpose the class rarely appears, and so far no natural problems have been exhibited which lie in $\mathsf{S_2^P}$ and not one of its more traditional subclasses (such as $\mathsf{BPP}$ or $\mathsf{NP}$). In addition it seems that $\mathsf{S_2^P}$ does not have a complete problem, due to its definition involving a promise. We have identified here a subclass $\mathsf{L_2^P}$ of $\mathsf{S_2^P}$, which is characterized exactly by a simple and natural total search problem, and which nonetheless maintains the interesting properties that motivated the original definition of $\mathsf{S_2^P}$:
\begin{theorem}\label{thm:l2p-structural-intro}
\phantom\\
\begin{enumerate}
    \item $\mathsf{P^{NP}} \subseteq \mathsf{L_2^P}$ and $\mathsf{BPP} \subseteq \mathsf{MA} \subseteq \mathsf{L_2^P}$
    \item $\mathsf{L_2^E}$\footnote{$\mathsf{L_2^E}$ is the exponential-time analogue of $\mathsf{L_2^P}$, where we replace ``polynomial time'' with $2^{O(n)}$ time'' in its definition} contains a language of circuit complexity $2^n/n$.
\end{enumerate}
\end{theorem}
In each case, the result stated for $\mathsf{L_2^P}$ was previously known to hold for $\mathsf{S_2^P}$ and is now shown to be inherited by the more natural subclass $\mathsf{L_2^P}$. The only interesting property that is known of $\mathsf{S_2^P}$ which we were unable to prove for $\mathsf{L_2^P}$ is the Karp-Lipton theorem; we discuss this further Sections~\ref{sec:open-probs} and \ref{sec:lop}.


\subsection{Open Problems}\label{sec:open-probs}
We conclude our introduction with a few interesting problems which remain open. The first is rather broad:
\begin{problem}
Show any additional inclusions or black-box separations which are not implied by the arrows in Figure~\ref{fig:classes}.
\end{problem}
We specifically highlight the following:
\begin{problem}\label{prob:lop-to-avoid}
Is there an $\mathsf{FP^{NP}}$ Turing reduction from LOP to Weak Avoid?
\end{problem}
Our interest in this problem is the following observation:
\begin{observation}
Say that Linear Ordering Principle is $\mathsf{FP^{NP}}$ Turing reducible to Weak Avoid. Then there is a particular language $L$ of circuit complexity $\geq 2^n/n$, and a deterministic $\mathsf{NP}$ oracle algorithm $\mathcal{A}$ running in time $2^{O(n)}$, such that given oracle access to any language $L'$ of circuit complexity $2^{\Omega(n)}$, $\mathcal{A}^{L'}$ computes $L$.
\end{observation}
This follows by composing the reduction of Weak Avoid to LOP, which produces a unique solution, with the (purported) second reduction from LOP back to Weak Avoid. Such a consequence would be rather surprising and interesting purely from the perspective of circuit complexity.

Next we highlight the problem of better clarifying the relationship between $\mathsf{L_2^P}$ and $\mathsf{S_2^P}$:
\begin{problem}
Does $\mathsf{L_2^P} = \mathsf{S_2^P}$? Can they be separated in the decision tree model? Does $\mathsf{L_2^P}$ satisfy a Karp-Lipton theorem?
\end{problem}
For this problem, we would say $\mathsf{L_2^P}$ ``satisfies a Karp-Lipton theorem'' if one could unconditionally prove the implication ``$\mathsf{NP} \in \mathsf{P/poly} \rightarrow \mathsf{PH}=\mathsf{L_2^P}$.'' A notable property of $\mathsf{S_2^P}$ is that it is the smallest complexity class for which this statement is known to hold.

The next problem we highlight is in the realm of depth 3 circuits:
\begin{problem}
Does Weak Avoid have depth 3 circuits of polynomial size?
\end{problem}
Recall that Li's upper bound is only quasipolynomial, of size around $N^{\log N}$. This problem seems intimately connected to the long-standing open question in proof complexity of whether the Weak Pigeonhole Principle has polynomial size bounded depth Frege proofs; a \emph{quasipolynomial} upper bound was shown by Paris Wilkie and Woods \cite{Paris-Wilkie-Woods} using a very similar technique to Li's, and \cite{MacielPW02} give a different bounded-depth Frege upper bound of lower depth, but still quasipolynomial size. Despite the strong aesthetic similarities we do not know a formal connection in either direction between these problems.


Lastly, the search problems discussed here have other connections to bounded arithmetic.
In particular, the LOP principle has been investigated in several papers, defined explicitly for the first time in \cite{Chiari1998-CHIWFI} and studied mostly within the context of characterizing the  strength of Jerebek's bounded arithmetic theory of approximate counting relative to weaker theories, and also as a new avenue for approaching the longstanding problem of separating Buss' $T_2$ hierarchy by sentences of fixed complexity. 
%(The latter problem is the proof theory analog of separating levels of the polynomial hierarchy by CNF/DNF functions.)
%($APC_1$) probabilistic polynomial-time reasoning
Buss, Kolodziejczyk and Thapen \cite{BussKT14} observe that the LOP principle is provable in both
$T^2_2$ and in $APC_2$, and ask whether or not  
LOP is provable in the weaker theory $T^1_2 ~ + ~ sWPHP$, 
where $sWPHP$ is the surjective weak pigeonhole principle, and corresponds to the search problem Weak Avoid.
Atserias and Thapen \cite{AtseriasT13} resolve this question,
proving that in the relativized setting,
$sWPHP$ does not prove the LOP principle over $T^1_2$.
In fact they prove a stronger result, that 
$sWPHP$ cannot prove the HOP principle over $T^1_2$,
where HOP is a
$\Sigma_1^b$ version of LOP. 
It seems possible that the techniques here could be used to give a negative answer to Problem~\ref{prob:lop-to-avoid}.
%that a converse to Theorem~\ref{thm:avoid-to-lop-intro} (reducing LOP to Weak Avoid) is unlikely to hold. 
A relatively unexplored area that is likely to be fruitful is to discover more relationships between natural search problems lying in the second level of the polynomial hierarchy (and higher) and corresponding systems of bounded arithmetic. 


%\todo{SAY SOMETHING NOT TOO DUMB.}



%%%%%%%%%%%%%
%%%%%%%%%%%%%
%%%%%%%%%%%%%

% \section{[OLD INTRO]}
% \subsection{Search Problems In Question}
% We now define the four $\mathsf{TF\Sigma_2^P}$ search problems which are the main focus of this work. It should be noted that from now on we will refer to the Range Avoidance problem, referenced in the introduction, as ``Weak Avoid.'' This is to distinguish it from the more general problem ``Strong Avoid.'' Both of these problems were originally defined in \todo{KKMP} under different names, however in subsequent work the name ``Avoid'' has become standard for the former problem.
% \begin{definition}[Weak Avoid]
% Given $f: \B^n \to \B^{n+1}$, find $y \in \B^{n+1} \setminus \rng(f)$
% \end{definition}

% \begin{definition}[Strong Avoid]
% Given $f: \B^n \setminus \{0^n\} \to \B^n$, find $y \in \B^n \setminus \rng(f)$.
% \end{definition}

% \begin{definition}[Strong 1-1 Avoid]
% Given $f: \B^n \setminus \{0^n\} \to \B^n$, find a pair $x\neq y $ in $\B^n \setminus \{0^n\}$ such that $f(x)=f(y)$, or else $y \in \B^n \setminus \rng(f)$, 
% \end{definition}
% \begin{definition}[Linear Ordering Principle]
% Given $\prec: \B^n \times \B^n \to \B$, find a witness that $\prec$ does not define a total ordering on $\B^n$, or else find the minimal element in the ordering it defines.
% \end{definition}

% In each case, the input function, e.g. $f: \B^n \to \B^{n+1}$, is given in the form of a boolean circuit with the appropriate number of input/output wires. We start by observing that for three of these problems, nontrivial upper bounds on their complexity follow directly from their definition.
% \begin{observation}\todo{ref}
% Weak Avoid is solvable in $\mathsf{FZPP^{NP}}$
% \end{observation}
% This is a consequence of the fact that a random string $y \sim \B^{n+1}$ will be a solution with probability $\geq 1/2$, and can be checked using an $\mathsf{NP}$ oracle. 
% \begin{observation}\label{obs:s-cap-pi-innitial}
% Linear Ordering Principle and Strong 1-1 Avoid are both solvable in $\mathsf{FP^{\Sigma_2^P \cap \Pi_2^P}}$.
% \end{observation}
% We will not prove this formally here, but it follows from the fact that both search problems have the following property: there are two types of solutions, the first being checkable in polynomial time and thus of a lower verification complexity, and if an instance does not have a low-complexity solution then it has a unique solution. In the case of Strong 1-1 Avoid, this is because a one-to-one map $[2^n] \to [2^n +1]$ will miss exactly one point in the codomain; for Linear Ordering Principle, this is because a linear ordering has a unique minimal element. Digging more into the structure of Linear Ordering Principle (and using the definition of $\mathsf{FS_2^P}$ which we have yet to present), we can also observe easily:
% \begin{observation}
% Linear Ordering Principle is solvable in $\mathsf{FS_2^P}$.
% \end{observation}

% \subsection{Main Results}
% In this work we establish new inclusions and separations amongst the subclasses of $\mathsf{TF\Sigma_2^P}$. Our results are summarized visually in Figure~\ref{fig:classes}, and we give a more in-depth description here.
% \begin{figure}
%     \centering
%     \includegraphics[width= \textwidth]{Classes.pdf}
%     % \captionsetup{width=1\textwidth}
%     \caption{Inclusion diagram of relevant classes and search problems. Solid arrows represent the inclusion of the class at the base of the arrow into the class at its tip. Dotted arrows indicate non-inclusion of the base class into the tip class in the decision tree model. The main linear axis of classes along the top are all included from left to right, indicated by the long solid arrow above them. Our major results are the separations marked with $\spadesuit, \clubsuit$ and the inclusion marked with $\blacklozenge$. \todo{Toni: for the classes in the top long line, are they also all separated in decision tree worls?}}
%     \label{fig:classes}
% \end{figure}
% \subsubsection{Weak Avoid and the Linear Ordering Principle}
% Our first set of results pertains to the newly defined Linear Ordering Principle problem, which we abbreviate as LOP. Amongst the novel results in \todo{the two papers}, the strongest is the upper bound of Li, which shows that Weak Avoid is contained in the class $\mathsf{FS_2^P}$. We show here that LOP represents a natural subclass of $\mathsf{FS_2^P}$ which is still powerful enough to capture the Weak Avoid problem. The main result is:
% \begin{theorem}
% There is a polynomial time many-one reduction from Weak Avoid to LOP.
% \end{theorem}
% This inclusion subsumes the new upper bounds for Weak Avoid given in \todo{ref}, since the inclusion $\mathrm{LOP} \in \mathsf{FS_2^P}$ follows directly from its definition.

% The proof follows much of the high level structure of \todo{Li's} result placing Weak Avoid in $\mathsf{FS_2^P}$. Roughly speaking, Li's proof shows that given an instance of Avoid $f: \B^n \to \B^{n+1}$, we can define a comparison relation $\sqsubset$ on $\B^{\poly(n)}$, so that for some unique distinguished element $\pi^* \in\B^{\poly(n)}$ we have $\pi^* \sqsubset \pi$ for all $\pi \neq \pi^*$, and $\pi^*$ contains a solution to the original Avoid instance. In our case we need to define a similar comparison relation, which in addition \emph{globally} acts as a total order on $\B^{\poly(n)}$; this requires some significant additional work. We note in addition that our proof, as well as Li's, bears a strong similarity to the work of \todo{PWW} which gave the first proof of the weak pigeonhole principle in the bounded arithmetic theory $S_2$. \todo{Toni: may want to talk up this result more?}

% \begin{theorem}
% ~
% \begin{enumerate}
%     \item Any search problem which has a $\mathsf{P^{NP}}$ Turing reduction to $\mathrm{LOP}$ also has a polynomial time many-one reduction to $\mathrm{LOP}$.
%     \item There is a language $\mathrm{DLOP} \subseteq \B^*$ (``Decisional LOP''), so that $\mathrm{LOP}$ is truth table reducible to $\mathrm{DLOP}$ and $\mathrm{DLOP}$ is many-one reducible to $\mathrm{LOP}$.
% \end{enumerate}
% \end{theorem}
% The advantage is that unlike $\mathsf{FS_2^P}$, which has a purely \emph{semantic} definition and therefore seems unlikely to contain a complete problem, LOP is defined by a very natural total search problem. 
% We next show that LOP is equivalent in complexity to a decision problem, and satisfies some strong closure properties under a wide class of reductions. In particular we define a \emph{decisional} complexity class $\mathsf{L_2^P}$ as follows:
% \begin{definition}
% A language $L$ is in the complexity class $\mathsf{L_2^P}$ if there is a polynomial time relation $R: (\B^*)^3 \to \B$ and a polynomial $p$, so that for all $x$, $R(x, \cdot, \cdot)$ defines a total order on $\B^{p(n)}$ whose minimal element $a$ has $a_1 = L(x)$. 
% \end{definition}

% We then have the following structural results:
% \begin{theorem}
% \phantom\\
% \begin{enumerate}
%     \item $L \in \mathsf{L_2^P}$ if and only if $L$ is $\mathsf{P^{NP}}
% $-Turing reducible to LOP. In particular $\mathsf{L_2^P}$ is closed under $\mathsf{P^{NP}}$ Turing reductions.
%     \item $\mathsf{L_2^P}$ has a complete language under many-one polynomial time reductions.
% \end{enumerate}
% \todo{Toni: Would it be more clear to define $L_2^P$ to be the class of decision problems that are polytime reducible to LOP. Or do you think it is important to specify decision version of LOP ?   Also for statement of Thm 2 maybe restate: if L is $P^{NP}$ 
% reducible to LOP, then L is polytime reducible to LOP. Thus $L_2^P$ is closed under $P^{NP}$ reductions.  (Oliver): I thought about both ways; my motivation of doing it this way is that people familiar with the definition of $\mathsf{S_2^P}$ will see this as a natural subclass of $\mathsf{S_2^P}$; then we state as a theorem that it has the equivalent characterization you mention.}
% \end{theorem}
% Combining this with the previous theorem reducing Weak Avoid to LOP we get:
% \begin{theorem}
% \phantom\\
% \begin{enumerate}
%     \item $\mathsf{L_2^P} \subseteq \mathsf{S_2^P}$
%     \item $\mathsf{P^{NP}} \subseteq \mathsf{L_2^P}$ and $\mathsf{BPP} \subseteq \mathsf{MA} \subseteq \mathsf{L_2^P}$
%     \item $\mathsf{L_2^E}$ contains a language of circuit complexity $2^n/n$.
% \end{enumerate}
% \end{theorem}
% % We believe this sequence of results greatly adds to our understanding of the new Weak Avoid upper bounds given in \todo{ref and ref}. 
% The main takeaway of these results is the following. $\mathsf{FS_2^P}$ and its decisional counter part $\mathsf{S_2^P}$ are somewhat unconventional complexity classes, and since their definition involves a promise they are believed not to contain complete problems. We have identified a natural subclass $\mathsf{L_2^P}$ of $\mathsf{S_2^P}$, which is characterized by a simple and natural search problem, thus endowing it with a complete language, but which otherwise maintains most of the power of the larger class $\mathsf{S_2^P}$. \todo{(Oliver): I'm still not sure I'm selling this well... do you have any other ideas for how to make this sound cooler?}
% \subsubsection{Circuit Lower Bounds for Pigeonhole Principles \todo{Lower Bounds and Separations?}}
% Our two main results in this section show that Strong Avoid has no non-trivial upper bound within the identified subclasses of $\mathsf{TF\Sigma_2^P}$, and that for Strong 1-1 Avoid, the $\mathsf{FP^{\Sigma_2^P \cap \Pi_2^P}}$ upper bound in Observation~\todo{ref} is the best possible:
% \begin{theorem}
% In the decision tree model, Strong Avoid is not solvable in $\mathsf{FP}^{\Sigma_2^P}_{||}$.
% \end{theorem}
% \begin{theorem}
% In the decision tree model, Strong 1-1 Avoid is not solvable in $\mathsf{FBPP^{NP}}$.
% \end{theorem}
% \todo{Toni: Perhaps reorder so thm 4 states both versions and we discuss that completely first, and then we state theorem 5 both versions and discuss it.}  
% The first lower bound is the more involved of the two, and requires proving a novel depth 3 $AC^0$ circuit lower bound for a total search problem, which appears to be the first circuit lower bound of this kind. In particular Theorem\todo{ref} is equivalent to the following circuit lower bound:
% \begin{theorem}
% \label{thm:depth3-avoid-lb-intro}
% Consider the natural association of $\B^{n\lceil \log n \rceil}$ with functions $f: [n-1] \to [n]$ specified by their truth table. Say that $C:\B^{n\lceil \log n \rceil} \to \B^{\lceil \log n \rceil}$ is a depth 3 unbounded fanin boolean circuit which solves the Strong Avoid problem in the sense that $C(f) \notin \rng(f)$ for all assignments $f$. Then $|C| \geq 2^{n^{\Omega(1)}}$. 

% More generally, if $C:\B^{n\lceil \log n \rceil} \to \B^{n^{\epsilon}}$ is a depth 3 circuit of size $2^{n^{\epsilon}}$  and $D: \B^{n^{\epsilon}} \to \B^{\lceil \log n \rceil}$ is an arbitrary postprocessing function, then $D \circ C$ cannot solve Strong Avoid provided $\epsilon$ is sufficiently small.

% %\todo{(Toni)This last sentence is a bit hard to understand. Maybe something like: Moreover, the lower bound holds even if $C$ can output $n^{\epsilon}$ output bits which can be postprocessed arbitrarily to solve Strong Avoid.  }\todo{(Oliver) is this better? Also is it ok for me to still hint at the "non-adaptive decision tree" formulation as below or do you think its still too confusing?}
% \end{theorem} 
% This latter result yields a very fine-grained separation for the Strong Avoid problem. It is easy to construct a depth 4 circuit $C:\B^{n\lceil \log n \rceil} \to \B^{\lceil \log n \rceil}$ solving Strong Avoid, where moreover the bottom fan-in is only $O(\log n)$ \todo{reference williams}. More strongly, it is possible to construct a depth $O(\log n)$ decision tree, which at each step queries a depth 3, $\poly(n)$ size circuit on the input $f$, and at each leaf outputs a correct solution $y \notin \rng(f)$. Our lower bound can be interpretted as saying that if such a decision tree is forced to be non-adaptive, then either the circuits it queries at each step must grow to exponential size, or else the number of queries must grow to $n^{\Omega(1)}$.

% This also contrasts the situation with Weak Avoid, where \todo{the new result} shows:
% \begin{theorem}\todo{cite}
% There is a depth 3 circuit $C:\B^{n\lceil \log 2n \rceil} \to \B^{\lceil \log 2n \rceil}$ of size $n^{O(\log n)}$ which solves the Weak Avoid problem: for every $f: [n] \to [2n]$, $C(f) \notin \rng(f)$. Moreover the bottom layer of gates in $C$ have fan-in $O(\log^2 n)$.
% \end{theorem}

% We believe our lower bound is of independent interest in %bounded depth 
% circuit complexity. In particular, we give a very fined grained depth-3 lower bound for a \emph{total search problem}. It is of course possible to construct contrived examples of total search problems which are hard for $\mathsf{AC}^0$ circuits, for example ``given $x \in \B^n$ output $b \in \B$ such that $\mathrm{Parity}(x)=b$;'' in such examples we can derive hardness of the search problem \emph{by reduction from a decision problem}. In contrast, the lower bound we show here \emph{cannot be established by reduction from any decision or function problem}. This follows from the fact that any decision problem (resp. any function problem) which is reducible to a $\mathsf{TF\Sigma_2^P}$ search problem lies in $\mathsf{\Sigma_2^P \cap \Pi_2^P}$ (resp. $\mathsf{FP^{\Sigma_2^P \cap \Pi_2^P}}$). In the decision tree model, this means that any language or function which can be solved by a small depth decision tree querying instances of Strong Avoid has both $\Sigma_2^P$ and $\Pi_2^P$ circuits. However, the lower bound we are trying to show for Strong Avoid rules out the existence of any $\mathsf{FP^{\Sigma_2^P \cap \Pi_2^P}}$ algorithm for Strong Avoid. Phrased more succinctly, our lower bound establishes that Strong Avoid is harder then \emph{any} decision problem which can be reduced to it, which by definition means we cannot establish the lower bound itself by reduction from decision problem. In light of this, to establish the lower bound in Theorem \todo{ref thm} we must develop new $\mathsf{AC^0}$ lower bound techniques which are specially tailored to total search problems.

% \todo{add a pointer to the section where we compare our switching lemma with those in proof complexity}
% \todo{Toni: should we add somewhere that we also get a fine grained separation with respect to the stretch of Avoid since by a simple reduction we get a superpoly depth-3 ckt lb also for avoid from [N] to [N + N/polylogn]?}

% Our second main separation for Strong 1-1 Avoid is equivalent to the following circuit lower bound: 
% \begin{theorem}
% Associate $\B^{n\lceil \log n \rceil}$ with functions $f: [n-1] \to [n]$ as above. Let $\mathcal{C}$ be a distribution over functions $C: \B^{n \lceil \log n \rceil} \to \B^{\lceil \log n \rceil}$, where each $C \in \mathrm{support}(\mathcal{C})$ is $\mathsf{P^{NP}}$ decision tree of complexity $w$, and such that for all one-to-one assignments $f: [n-1]\to[n]$, 
% \[
% \Pr[\mathcal{C}(f) \notin \rng(f)] \geq \frac{2}{3}
% \]
% Then $w \geq n^{\Omega(1)}$.
% \end{theorem}

% We will review the formal definition of $\mathsf{P^{NP}}$ decision trees in \todo{section}. Informally, these are decision trees of small height which, instead of querying variables, are allowed to query the value of an arbitrary low-width DNF over the inputs in unit time. Recall that Strong 1-1 Avoid is easily shown to lie in the class $\mathsf{FP^{\Sigma_2^P \cap \Pi_2^P}}$ (Observation~\ref{obs:s-cap-pi-innitial} above, which will be proven formally in Lemma~\ref{lem:1-1-lop-unique} later). Combining this with the above lower bound we obtain as a corollary the following separation for decision classes:
% \todo{(Toni): shouldnt this be $P^{\Sigma_2 \cap \Pi_2}$ below? (Oliver):$P^{\Sigma_2 \cap \Pi_2} = \Sigma_2 \cap \Pi_2$ so it's the same thing }
% \begin{theorem}
% In the decision tree model, $\mathsf{\Sigma_2^P \cap \Pi_2^P} \not\subseteq \mathsf{BPP^{NP}}$. In particular $\mathsf{\Sigma_2^P \cap \Pi_2^P} \not\subseteq \mathsf{S_2^P}$.
% \end{theorem}
% This is in contrast to the situation for $\mathsf{NP} \cap \mathsf{coNP}$, which is known to collapse to $\mathsf{P}$ in the decision tree model. This improves a previous result of \todo{fortnow and someone} who showed that $\mathsf{\Sigma_2^P \cap \Pi_2^P} \not\subseteq \mathsf{P^{NP}}$ in the decision tree model.
% % \subsection{Techniques}
% % \todo{(Oliver): Should we include this? Or is it too redundant with our "Proof Overview" in section 3 and we can just postpone everything to then?}
% % The main ingredient in our circuit lower bounds is a variant of the switching lemma. Roughly speaking, this lemma says the following. Say we are given a small-width DNF $D$ on $\approx N \log N$ boolean variables, where as above we identify each variable assignment with a function $f: [N] \to [N+1]$ in the natural way. As is common in random restriction arguments, we may choose a certain partial function $\rho: [N] \to [N+1]$, and restrict attention to only those assignments $f$ which extend $\rho$. Our switching lemma shows that if $\rho$ is sampled from an appropriate distribution, then after the restriction we may replace $D$ with a small height ``pigeonhole decision tree'' $T$ which has the same behavior on all assignments extending $\rho$. However, this decision tree $T$ is \emph{not} a standard decision tree querying the underlying variables defining $f$. Instead it is allowed to make both pigeon queries (what is $f(x)$?) as well as \emph{hole queries} (does $y$ have a preimage under $f$, and if so what is it?). This corresponds to a special kind of $\mathsf{P^{NP}}$ decision tree over the underlying variables. We then need a separate argument, not involving a random restriction, to show that the simplified circuit cannot solve the search problem. A different argument is needed for Theorems \todo{ref and ref}. 

% % In the more involved case of the depth 3 circuit lower bound, we start out with a depth 3 circuit $\phi = \land_i D_i$ where $D_i$ are DNFs. After the random restriction, we are able to represent each $D_i$ by a small-height pigeonhole decision tree $T$, which in turn allows us to represent $\lnot D_i$ by a small-width \emph{hole disjunction}, which is a special type of depth 3 circuit which in some sense resembles a DNF. Overall this gives us a representation of $\lnot \phi = \lor_i \lnot D_i$ as a small-width hole disjunction. Although this translation does not reduce the depth of $\phi$ in the traditional sense (it is still depth 3), the special structure of the simplified circuit allows for a direct argument to be applied.

% % Our switching lemma bears \todo{stronger resembelence to proof complexity ones then circuit complexity ones. However highlight following distinction: we need the decision tree to represent the original DNF semantically as a function, rather then in some syntactic sense as a formula. It is also distinct from switching lemmas which have appeared in circuit complexity, in that the decision tree we end up with is not a standard decision tree on the underlying variables, but rather a specialized kind of $\mathsf{P^{NP}}$ decision tree.} A consequence of this latter fact is that we cannot hope to apply the switching lemma repeatedly. This \todo{makes sense} in light of the fact that Strong Avoid \emph{does have} polynomial size circuits of depth 4.
% \subsection{Open Problems}
% We conclude our introduction with a list of natural problems that remain open:
% \begin{enumerate}
%     \item Is Linear Ordering Principle reducible to Weak Avoid?
%     \item Is Linear Ordering Principle reducible to Strong Avoid?
%     \item Is Weak Avoid reducible to Strong 1-1 Avoid?
%     \item Can we separate $\mathsf{S_2^P}$ from $\mathsf{L_2^P}$?
%     \item Does Weak Avoid have depth 3 circuits of polynomial size?
% \end{enumerate}
% In each case the problem remains interesting whether reductions are taken to be polynomial time or $\mathsf{FP^{NP}}$. We highlight the first problem in particular, whose positive answer would have the following interesting consquence:
% \begin{observation}
% Say that Linear Ordering Principle is $\mathsf{FP^{NP}}$ Turing reducible to Weak Avoid. Then there is a particular language $L$ of circuit complexity $\geq 2^n/n$, and a deterministic $\mathsf{NP}$ oracle algorithm $\mathcal{A}$ running in time $2^{O(n)}$, such that given oracle access to any language $L'$ of circuit complexity $2^{\Omega(n)}$, $\mathcal{A}^{L'}$ computes $L$.
% \end{observation}

\section{Preliminaries}
\label{sec:prelims}
\subsection{Search Problems, Complexity Classes, and Basic Inclusions}
\label{sec:totalsearch}
We define here a variety of subclasses inside of $\mathsf{TF\Sigma_2^P}$, classified according to the computational resources necessary to solve a search problem. We then prove some of the more basic results relating these classes to eachother and to the four main search problems of interest in this work.

\begin{definition}[$\mathsf{FP^{NP}}$, $\mathsf{FP^{\Sigma_2^P \cap \Pi_2^P}}$, $\mathsf{FP_{||}^{\Sigma_2^P}}$]
Let $R$ be a search problem and $\mathcal{C}$ a class of decision problems. We say $R \in \mathsf{FP}^\mathcal{C}$ if there is a language $L \in \mathcal{C}$ and a polynomial time algorithm making queries to $L$ which ``solves $R$:'' given $x$ it outputs $y$ such that $(x,y) \in R$. We say $R \in \mathsf{FP}_{||}^{\mathcal{C}}$ if there is such an algorithm, which moreover makes its queries \emph{nonadaptively}: given an input $x$ it computes in polynomial time a list of queries $z_1,\ldots,z_m$, uses its $L$-oracle to test in unit time the membership of each $z_i$ in $L$, and then uses the oracle responses to output an answer in polynomial time.
\end{definition}

We include for reference the following class considered in \cite{Li,CHR}:
\begin{definition}[$\mathsf{svF\Sigma_2^P}$]
$\mathsf{svF\Sigma_2^P}$ is the class of search problems having``singled-valued'' $\mathsf{F\Sigma_2^P}$ algorithms. We say $R$ lies in this class if there is a choice of canonical solution $\{(x,y_x)\in R \mid x\}$ for each input $x$, a second relation $R' \in \mathsf{TF\Sigma_2^P}$, and a polynomial time function $f$ so that whenever $(x,z) \in R'$, $f(z)=y_x$.
\end{definition}
In Lemma~\ref{lem:inclusions} we relate this class to the others we have defined here. Beyond this we will not need to reference this class further: every time we prove an upper bound for a search problem it will be in a class lower then $\mathsf{svF\Sigma_2^P}$, and our separation for Strong Avoid will hold even against the larger class $\mathsf{FP_{||}^{\Sigma_2^P}} \supseteq \mathsf{svF\Sigma_2^P}$. We next review the randomized classes:

\begin{definition}[$\mathsf{(ps)FZPP^{NP}}$, $\mathsf{(ps)FBPP^{NP}}$]
A relation $R$ is in $\mathsf{FBPP^{NP}}$ if there exists a randomized polynomial time algorithm with access to a $\mathrm{SAT}$ oracle which, given $x$, outputs $y$ such that $(x,y) \in R$ with probability $\geq 2/3$. If the algorithm always outputs a valid answer or $\bot$ and answers $\bot$ with probability $< \frac{1}{3}$ this places $R$ in the subclass $\mathsf{FZPP^{NP}}$.

For a search problem $R$ where each $x$ may have many solutions $y$, it is possible that a randomized algorithm outputs different correct answers on the same input $x$ as a function of its random coin tosses. If for each $x$ there exists a canonical $y_x$ with $(x,y_x) \in R$ and some randomized algorithm computes $x \mapsto y_x$ with high probability, we say that algorithm is pseudodeterministic; we use the prefix $\mathsf{ps}-$ to denote the pseudodeterministic analogue of a randomized class.
\end{definition}
The last standard complexity class we examine is the functional analogue of $\mathsf{S_2^P}$, defined as follows:
\begin{definition}[$\mathsf{FS_2^P}$]
A search problem $R$ is in $\mathsf{FS_2^P}$ if there exists polynomial time algorithm $\mathrm{V}$ taking three inputs, so that for all $x$ there exists $y_x$ with $(x,y_x) \in R$ so that:
\begin{enumerate}
    \item There exists $\pi_1$ such that for all $\pi_2$, $\mathrm{V}(x,\pi_1,\pi_2)=y_x$.
    \item There exists $\pi_2$ such that for all $\pi_1$, $\mathrm{V}(x,\pi_1,\pi_2)=y_x$.
\end{enumerate}
\end{definition}
% All of the above search classes have decision analogues, where we restrict attention to total search problems of the form $x \mapsto \mathbbm{1}\{x \in L\}$ for some language $L$ and apply the same definition. In such cases we simply omit the prefix $\mathsf{F}$ from the class name.

\paragraph{Search Problems vs. Function Problems:}
As discussed in the introduction, some kinds of algorithms for search problems have the property that they associate to each input a \emph{fixed solution} which the algorithm produces on that input. With the exception of the non-pseudodeterministic randomized classes $\mathsf{FBPP^{NP}}$ and $\mathsf{FBPP^{NP}}$, each of the classes we have just defined describes algorithms of this sort. These classes are arranged nicely in the following hierarchy, as indicated in Figure~\ref{fig:classes}:

\begin{lemma}\label{lem:inclusions}
$\mathsf{FP} \subseteq \mathsf{FP^{NP}} \subseteq \mathsf{FS_2^P} \subseteq \mathsf{psFZPP^{NP}} \subseteq \mathsf{FP^{\Sigma_2^P \cap \Pi_2^P}} \subseteq \mathsf{svF\Sigma_2^P} \subseteq \mathsf{FP_{||}^{\Sigma_2^P}}$
\end{lemma}
\begin{proof}
All inclusions follow directly from the definition, with the exceptions of $\mathsf{FP^{NP}} \subseteq \mathsf{FS_2^P}$ which is due to Russell-Sundaram \cite{russell-sundaram} and $\mathsf{FS_2^P} \subseteq \mathsf{psZPP^{NP}}$ which is due to Cai \cite{cai}.
\end{proof}

Recall the class $\mathsf{TFU\Sigma_2^P}$, referenced heavily in the introduction, consisting of those $\mathsf{TF\Sigma_2^P}$ search problems with unique solutions. This class is directly associated to $\mathsf{FP^{\Sigma_2^P \cap \Pi_2^P}} $ in this hierarchy:
\begin{lemma}\label{lem:unique-intersect}
The following are equivalent for any $R \in \mathsf{TF\Sigma_2^P}$:
\begin{enumerate}
    \item $R$ is polynomial time reducible to a problem in $\mathsf{TFU\Sigma_2^P}$.
    \item $R \in \mathsf{FP^{\Sigma_2^P \cap \Pi_2^P}}$.
\end{enumerate}
\end{lemma}
\begin{proof}
Say $R$ is polynomial time reducible to $R' \in \mathsf{TFU\Sigma_2^P}$. So there are polynomial time functions so that for all $x$, if $(f(x),y)\in R'$ then $(x,g(x,y))\in R$. Consider the language $L$ defined as follows: $(x,i) \in L$ iff for the unique $y$ with $(f(x),y)\in R'$, we have $y_i=1$. Then $L \in \mathsf{\Sigma_2^P \cap \Pi_2^P}$ and $g$ yeilds a reduction from $R$ to $L$.

In the other direction note that any language $L \in \mathsf{\Sigma_2^P \cap \Pi_2^P}$ defines a search problem in $\mathsf{TFU\Sigma_2^P}$: given $x$ find $b \in \B$ so that $L(x)=b$. The result follows directly from this fact. 
\end{proof}

Recall from the introduction that the problems LOP and Strong 1-1 Avoid do not quite have unique solutions, but come very close. We define this property of having \emph{essentially unique solutions} as follows:
\begin{definition}[Essentially Unique Solutions]
We say that a total search problem $R \in \mathsf{TF\Sigma_2^P}$ has ``essentially unique solutions'' if there are verifiers $\mathrm{V_1}, \mathrm{V_2}$ such that:
\begin{enumerate}
    \item $\mathrm{V}_1$ is testable in polynomial time, while $\mathrm{V_2}$ is testable in $\mathsf{coNP}$.
    \item For all $x$, either there exists $y$ so that $\mathrm{V}_1(x,y)=1$ and $(x,y) \in R$, or else there exists a unique $y$ such that $\mathrm{V}_2(x,y)=1$ and $(x,y) \in R$.
\end{enumerate}
\end{definition}
We then have:
\begin{lemma}\label{lem:unique-sol}
If $R$ has essentially unique solutions then it is polynomial time reducible to a search problem $R' \in \mathsf{TFU\Sigma_2^P}$ which actually has unique solutions. By Lemma~\ref{lem:unique-intersect} this is equivalent to the statement $R \in \mathsf{FP^{\Sigma_2^P\cap \Pi_2^P}}$.
\end{lemma}
\begin{proof}
Let $\mathrm{V}_1,\mathrm{V}_2$ witness that $R$ has essentially unique solutions. Consider the search problem $R'$: given $x$, output either the lexicographically first $y$ such that $\mathrm{V}_1(x,y)=1$, or else the unique $y$ such that $\mathrm{V}_2(x,y)=1$. By the definition of $\mathrm{V}_1,\mathrm{V}_2$ we see that $R'$ is a total search problem with unique solutions. Clearly $R$ is polynomial time reducible to $R'$. We need to show that $R' \in \mathsf{F\Sigma_2^P}$. For a fixed $x$, say that there exists $y$ with $\mathsf{V}_1(x,y)=1$, and let $y_0$ be the lexicographically first such $y$. Then we can confirm that $(x,y_0) \in R'$ in $\mathsf{coNP}$ by confirming $\mathsf{V}_1(x,y_0)=1$ and that for all $y'<y_0$ we have $\mathsf{V}_1(x,y')=0$. On the other hand say that there is no $y$ with $\mathsf{V}_1(x,y)=1$, and let $y^*$ be the unique element with $\mathsf{V}_2(x,y^*)=1$. Then we can confirm that $(x,y^*) \in R'$ in $\mathsf{coNP}$ by checking that for all $y$ we have $\mathsf{V}_1(x,y)=0$, and using $\mathsf{V}_2$ to confirm $\mathsf{V}_2(x,y^*)=1$ (recall that $\mathsf{V}_2$ is in $\mathsf{coNP}$ by definition).
% Consider the language $L$ defined as follows: $(x,i) \in L$ iff for all $y$, $\mathrm{V}_1(x,y)=0$ and $\mathrm{V}_2(x,y)=1 \rightarrow y_i = 1$. We first claim that $L \in \mathsf{\Sigma_2^P \cap \Pi_2^P}$. To certify that $(x,i) \in L$ we supply $y$ such that $\mathrm{V}_2(x,y)=1$ and $y_i = 1$, and to certify $(x,i) \notin L$ we supply either a $y$ such that $\mathrm{V}_1(x,y)=1$ or else a $y$ such that $\mathrm{V}_2(x,y)=1$ and $y_i = 0$. In either case we can verify the certificate in $\mathsf{coNP}$; in the case of a positive certificate we can use the $\mathsf{coNP}$ verifier to confirm that there are no $\mathrm{V}_1$ solutions for $x$. Now to solve the search problem we can first use an $\mathsf{NP} \subseteq \mathsf{\Sigma_2^P \cap \Pi_2^P}$ oracle to find a solution verified by $\mathrm{V}_1$ if one exists. If found we are done, and if not then we can query the language $L$, which lies in $\mathsf{\Sigma_2^P \cap \Pi_2^P}$, to determine all bits of the unique $y$ which is verified by $\mathrm{V}_2$.
\end{proof}
We now state formally the claim made in the introduction that LOP and Strong 1-1 Avoid have essentially unique solutions:
\begin{lemma}\label{lem:1-1-lop-unique}
Linear Ordering Principle and Strong 1-1 Avoid have essentially unique solutions; hence both problems lie in $\mathsf{FP^{\Sigma_2^P \cap \Pi_2^P}}$.
\end{lemma}
\begin{proof}
Any injective function $f: [N] \to [N+1]$ leaves a unique point in $[N+1]$ out of its range. Every total linear order on a finite set has a unique minimal element.
\end{proof}
We next prove formally that LOP is contained in the class $\mathsf{FS_2^P}$, which follows quite directly from the definitions:
\begin{lemma}\label{lem:lop-in-s2p}
Linear Ordering Principle is in $\mathsf{FS_2^P}$.
\end{lemma}
\begin{proof}
Let $\prec: \B^n \times \B^n \to \B$ be an instance of Linear Ordering Principle. We construct the $\mathsf{FS_2^P}$ solver $\mathrm{V}$ for $\prec$ as follows. Let $X = \B^{3n+1}$ be partitioned so that the first $2^n$ elements $A$ are identified with $\B^n$, and the remaining $>2^{3n}$ elements $B$ are identified with potential witnesses that $\prec$ fails to define a linear order. Given $\pi_1,\pi_2 \in X$, $\mathrm{V}(\prec, \pi_1,\pi_2)$ behaves as follows:
\begin{enumerate}
    \item If $\pi_i \in B$ codes a witness that $\prec$ does not define a linear order for some $i \in \{1,2\}$, we output $\pi_i$; if both do then we output the lexicographically first between them.
    \item Say $\pi_i \in B$ and $\pi_{i'} \in A$ for $\{i,i'\}=\{1,2\}$, and $\pi_i$ is not a witness that $\prec$ fails to define a linear order. In this case we output $\pi_{i'}$.
    \item If both $\pi_1,\pi_2 \in A$, we think of them as representing elements $a_1,a_2$ of $\B^n$ and compare them according to $\prec$. If $a_1 \prec a_2$ we output $\pi_1$, otherwise we output $\pi_2$.
\end{enumerate}
First say that $\prec$ is not a total order. Among all witnesses to this let $\pi \in B$ be the lexicographically first. Then for all $\pi' \in X$, we have $\mathrm{V}(\prec,\pi,\pi')=\mathrm{V}(\prec,\pi',\pi)=\pi$. On the other hand say $\prec$ is a total order and let $\pi \in A$ correspond to its unique minimal element. Then again for all $\pi' \in X$, we have $\mathrm{V}(\prec,\pi,\pi')=\mathrm{V}(\prec,\pi',\pi)=\pi$. Thus in all cases $\mathrm{V}$ gives an $\mathsf{FS_2^P}$ algorithm solving the Linear Ordering Principle problem on input $\prec$.
\end{proof}
Finally, amongst the 4 search problems studied here we have the two following obvious inclusions which we haven't mentioned yet:
\begin{observation}
Weak Avoid and Strong 1-1 Avoid are polynomial time reducible to Strong Avoid.
\end{observation}
\subsection{Oracle Separations and the Decision Tree Model}
\label{prelim:querymodels}

All of the upper bounds and inclusions we show in this paper are unconditional and hold relative to every oracle. Since showing any unconditional separations amongst the classes we have identified would imply $\mathsf{P} \neq \mathsf{NP}$, we can only hope to establish separations in a restricted model. As is standard, our restricted model will correspond to the ``decision tree model'' of complexity classes, which can be framed either in terms of oracles and Turing machines, or more directly in terms of decision trees and bounded depth circuits, c.f. \cite{cook}. We choose here the latter terminology. 

All of our search problems are defined in terms of a function or relation specified by a boolean circuit. Take for example the Weak Avoid problem, whose instance is a boolean circuit computing $f: \B^n \to \B^{n+1}$. Clearly this problem remains total if $f$ were not represented by a small circuit, but instead was an arbitrary function $f: [2^n] \to [2^{n+1}]$. In the decision tree model, the relevant search problem has as its input an arbitrary function $f: [N] \to [2N]$ where $N = 2^n$. For example, the decision tree variant of the simplest complexity class $\mathsf{FP}$ then corresponds to algorithms which can access $f$ only by querying its on $\poly(n) = \poly(\log N)$ different inputs $x \in [N]$. If we imagine $f: [N] \to [2N]$ is specified by an assignment $\alpha: \B^{[2^n] \times [n+1]}$, where $\alpha_{x,i} = f(x)_i$, then this is equivalent to allowing the algorithm to query $\poly(n)=\poly(\log N)$ variables of the assignment $\alpha$. 

Decision tree analogues of other complexity classes can be defined similarly. It is a standard result that a decision tree separation implies a a separation of the associated Turing machine classes relative to an oracle; more specifically it is equivalent to a separation relative to a so-called ``generic oracle'' \cite{BeameCEIP98}.

Most of the definitions in our decision tree models will be standard, e.g. bounded depth formulae and DNFs.
One decision tree model we study whose definition and notation is less standardized is a $\mathsf{FP^{NP}}$ decision tree:
\begin{definition}[$\mathsf{FP^{NP}}$ Decision Trees]
\label{def:PNP-trees}
Let $f: \B^n \to A$ be a function where $A$ is some set. A $\mathsf{P^{NP}}$ decision tree $T$ computing $f$ is defined by a binary tree, with each internal node labeled by a DNF formula on the variables $\{x_i \mid i \leq n\}$ and each leaf labeled by a value $y \in A$. On an input $x \in \B^n$, we traverse $T$ starting at the root. At each internal node associated to a DNF $D$, we test if $D(x)=1$; if so we proceed to the right child of the current node, otherwise we proceed to the left. When we reach a leaf we output the value associated to it; $T$ computes $f$ if the value reached is $f(x)$. We say that $T$ has complexity $\leq r$ if its depth as a tree is at most $r$, and each DNF associated to its nodes has width at most $r$.
\end{definition}

We next define the decision tree variant of $\mathsf{FP^{\Sigma_2^P}_{||}}$ which is the subject of our main lower bound Theorem~\ref{thm:strongavoid-lb}:
\begin{definition}\label{def:sigma-parallel}
Let $R \subseteq \B^{N} \times [M]$ black-box (relativized) search problem. We say that $R$ is in $\mathsf{FP^{\Sigma_2^P}_{||}}$ if there exist $\mathsf{\Sigma_2^P}$ circuits $\Phi_1,\ldots,\Phi_k$ with $k \leq \poly\log(N)$ and an arbitrary post-processing function $S: \B^k \to [M]$ so that for all inputs $f \in \B^N$, $(f, S(\Phi_1(f),\ldots,\Phi_k(f))) \in R$. A $\mathsf{\Sigma_2^P}$ circuit is a $2^{\poly\log N}$-size depth 3 circuit with bottom fanout $\leq \poly\log N$.
\end{definition}

In our case the input will be some $f \in \B^{N\log N}$ representing a function $f: [N] \to [N+1]$, and the relevant search problem is to output some $y \in [N+1]$ outside its range. The fact that this captures the relativized version of $\mathsf{FP^{\Sigma_2^P}_{||}}$ is based on the well known equivalence between relativized levels of the polynomial hiearchy and quasipolynomial size bounded depth circuits \cite{FSS}; in particular the $\mathsf{\Sigma_2^P}$ circuits $\Phi_1,\ldots,\Phi_k$ above correspond to a sequence of non-adaptive $\mathsf{\Sigma_2^P}$ oracle queries.

\paragraph{Note on bottom fanin:} The depth 3 circuit model which captures $\mathsf{FP^{\Sigma_2^P}_{||}}$ has the additional restriction that the bottom level of each depth 3 circuit has fanin $\poly\log N$. Our lower bound will apply also to the stronger model in which the depth 3 circuits are only constrained in their size and not their bottom fanin.

Finally, we make note of one separation indicated in Figure~\ref{fig:classes} which was shown prior to our work:
\begin{lemma*}[Wilson and Vyas-Williams\cite{wilson}\cite{vyas-williams}]
In the decision tree model, Weak Avoid is not in $\mathsf{FP^{NP}}$.
\end{lemma*}
% \subsection{Connection Range Avoidance Algorithms and Circuit Lower Bounds}
% \todo{I think maybe it would be good to lay out the formal proof that, say, a $\mathsf{psZPP^{NP}}$ algorithm for Weak Avoid implies the circuit lower bound $\mathsf{ZPE^{NP}} \notin \mathsf{size}[2^n/n]$, for the reader who's unfamiliar with how this connection works.}
% As mentioned in the introduction, it is known that improved algorithms for Weak Avoid imply circuit lower bounds for uniform complexity classes, and this connection was used to give the new circuit lower bounds in \todo{refs}. We include here a proof of this 

\section{Lower Bounds for Pigeonhole Principles}
\subsection{Pigeonhole Principle Basics}
\label{sec:php-basics}


We will be concerned here with search problems where the input is a function $f: [N] \to [M]$ with $M > N$, and the goal is to find an empty pigeonhole of $f$. 
\begin{definition}
We use $[M]^{[N]}$ to refer to the set of all functions $f: [N] \to [M]$.
% and to $\mathcal{F}_{N,M}$ to refer to the set of all 1-1 (injective) functions $[N] \to [M]$. We use the shorthand $\mathcal{F}_{N} = \mathcal{F}_{N,N+1}$.
We define a set of propositional variables, referred to as ``bit variables'', given by 
\[
\mathrm{BITS}_{N,M}:=\{f_{x,i} \mid x \in [N], i \in [\lceil \log M \rceil]\}
\] We associate each truth assignment $\alpha: \mathrm{BITS}_{N,M} \to \B$ with the function $f_{\alpha}: [N] \to [M]$, where where the value $\alpha$ assigns to $f_{x,i}$ indicates the $i^{th}$ bit of $f_{\alpha}(x)$ in binary; if $M$ is not a power of two then we think of all strings in $\B^{\lceil \log M \rceil}$ of binary value exceeding $M$ as being redundant representations of the element $M$. Similarly we associate every function $f: [N] \to [M]$ with the assignment $\alpha_f: \mathrm{BITS}_{N,M} \to \B$ using the same correspondence. 

For a conjunction of literals (term) $\tau$ over the varaibles $\mathrm{BITS}_{N,M}$, we define its \emph{pigeon-width} to be the size of the set $\{x \in [N]\mid  f_{x,i} \text{ or } \lnot f_{x,i} \text{ occurs in } \tau \text{ for some } i \}$.
\end{definition}
From now on we will often not refer explicitly to assignments $\alpha: \mathrm{BITS}_{N,M} \to \B$, only to functions $f: [N] \to [M]$. 
The relevance of this definition is that, when defining circuits/computational devices whose input is a function $f: [M] \to [N]$, we must specify how the device is able to access/read the input $f$: typically this will correspond to the ability to read the bits $\mathrm{BITS}_{N,M}$.

For the vast majority of this section we will restrict our attention to functions $f: [N] \to [M]$ which are 1-1 (injective): consequently when we speak about the evaluation of a formula which takes $f$ as an input, we will typically only care about the behavior of that formula on this special class of inputs. We use the following notation to express this:
\begin{definition}
Let $\mathcal{F}_{N,M} \subset [M]^{[N]}$ denote the set of all 1-1 functions $[N] \to [M]$. For two predicates $F,G:[M]^{[N]}\to \B$, we say that $F$ is equivalent to $G$ (with respect to all 1-1 functions), denoted by  $F \equiv G$, if $F(f)=G(f)$ for all $f \in \mathcal{F}_{N,M}$.
\end{definition}

We now introduce notation for describing partial information about an input $f$. A natural such unit of partial information is a partial assignment $\rho : \mathrm{BITS}_{N,M} \to \{0,1,\star\}$, however for our purposes it will be useful to define additional types of partial information as well:

\begin{definition}[Matchings and Hole Restrictions]
A partial matching $\pi$ is a partially defined 1-1 function from $[N]$ to $[M]$. We use $\dom(\pi), \rng(\pi)$ to refer to the domain and range of $\pi$, and $\nodes(\pi)= \dom(\pi) \sqcup \rng(\pi)$. More generally a ``hole restriction'' $\tau = \langle \pi, E\rangle$ consists of a partial matching $\pi$ and a set $E \subseteq [M]$ satisfying $\rng(\pi) \cap E = \emptyset$.  Let $\nodes(\tau) = \nodes(\pi) \cup E$. We use $|\tau|$ to refer to the value $|\pi| + |E|$. We think of the hole restriction $\langle \pi, E\rangle$ as describing the following partial information about a function $f: [N] \to [M]$: ``$f(x) = \pi(x)$ for all $x \in \dom(\pi)$, and $y \notin \rng(f)$ for all $y \in E$.''  For a total assignment $f \in {\cal F}_{N,M}$, we say $f$ is compatible with $\langle \pi, E\rangle$, written $f \Vert \langle \pi, E\rangle$, if the above statement holds. For two hole restrictions $\tau_1=\langle \pi_1,E_1\rangle, \tau_2=\langle \pi_2,E_2\rangle$, we say they are consistent, written $\tau_1 \Vert \tau_2$, if there is a total assignment $f$ which is compatible with both. We say $\tau_1$ extends $\tau_2$, written $\tau_1\supseteq \tau_2$, if $\pi_1 \supseteq \pi_2$ and $E_1 \supseteq E_2$. We will think of a partial matching $\pi$ as a special kind of hole restriction of the form $\langle \pi, \emptyset\rangle$, and a total 1-1 assignment $f$ as a special hole restriction of the form $\langle f, \corng(f)\rangle$, where $\corng(f)= [M] \setminus \rng(f)$. In this way we will use the above terminology to define relations between matchings, hole restrictions and total 1-1 assignments, e.g. $f \Vert \pi$ and $\pi \subseteq \tau$. 

\end{definition}

The main subject of this section is circuits $C: \B^{|\mathrm{BITS}_{N,M}|} \to \B^{\lceil \log M \rceil}$ which solve Avoid, in the sense that $C(f) \notin \rng(f)$ for all $f$, where $C(f)$ denotes feeding $C$ the representation of $f$ as an assignment to its bit variables. We thus need some notation for basic kinds of circuits computing a function of $f$, which will be hole and matching disjunctions:
\begin{definition}[Matching and Hole Disjunctions]
A hole disjunction $\phi = \lor_i \tau_i$ over $[M]^{[N]}$ is defined by a collection of hole restrictions $[N] \to [M]$, which we refer to as the ``terms'' of $\phi$. We say $\phi(f)=1$, or ``$f$ satisfies $\phi$,'' if $f \Vert \tau_i$ for some term $\tau_i$ of $\phi$. In this way each hole disjunction is associated with a boolean function $\mathcal{F}_{N,M} \to \B$. If $|\tau_i|\leq w$ for all $i$ we say that $\phi$ has width $w$. We call $\phi$ a matching disjunction if all $\tau_i$ are partial matchings.

% For a hole restriction $\tau = \langle \pi, E \rangle$ and a partial matching $\sigma$, we denote by $\tau \upharpoonright \sigma$ the \todo{simplified}. Similarly for a hole disjunction $\phi$ we use $\phi \upharpoonright \sigma$ \todo{simplified disjunction}.
\end{definition}
Recall that we are primarily concerned here with the values a formula takes on 1-1 assignments $f \in \mathcal{F}_{N,M}$. In such a setting we can simplify any circuit so that the bottom two logical layers are a set of matching disjunctions:
\begin{observation}\label{obs:matching-conversion}
Let $D$ be a DNF formula on the bit variables $\mathrm{BITS}_{N,M}$. Then there is a matching disjunction $\phi$ so that $D \equiv \phi$. Moreover if $D$ has pigeon-width $w$ and size $s$, then $\phi$ will have width at most $w$ and size at most $sM^w$.
\end{observation}

\begin{proof}
Say that $D = \lor_i \land_j \ell_{i,j}$ is a DNF where each $\ell_{i,j}$ is a literal on $\mathrm{BITS}_{N,M}$. For each term $t_i = \land_j \ell_{i,j}$ we may replace it by the matching disjunction:
\[
\theta_i = \bigvee_{\pi \in \mathrm{match}(t_i)} \pi
\]
where $\mathrm{match}(t_i)$ is the set of minimal partial matchings $\pi$ so that for each literal $\ell_{i,j} = (\lnot)^{\xi} f_{x,b}$ in $t_i$ (with $\xi \in \B$), there is an edge $(x,y) \in\pi$ with $y_b = \lnot \xi$. Clearly we have that $\theta_i(f) = t_i(f)$ for all 1-1 assignments $f$, and thus we may replace $D$ with the matching disjunction $\lor_i \theta_i$ without affecting its behavior on 1-1 assignments. Recall that the pigeon-width of a term over $\mathrm{BITS}_{N,M}$ is the total number of pigeons $x \in [N]$ mentioned in the term; clearly the width of $\theta_i$ is at most the pigeon width of $t_i$ since it only mentions the $\leq w$ pigeons relevant to $t_i$; for this same reason we may bound its size (number of terms) by $M^w$, the number of mappings from these $\leq w$ pigeons onto the $M$ available holes.
\end{proof}

\begin{definition}
    \label{def-restricting-a-dnf}
    
Let $\phi = \tau_1 \lor \ldots \lor \tau_s$ be a hole disjunction with $\tau_i = \langle \pi_i, E_i\rangle$ and let $\kappa = \langle \sigma,U\rangle$ be a hole restriction.
We define $\phi$ restricted by $\tau$, denoted by $\phi \upharpoonright\kappa$, as follows:
(i) First, any term $\tau_i \in \phi$ that is inconsistent with $\kappa$ is set to 0 by $\kappa$ and thus these terms disappear from $\phi \upharpoonright \kappa$ (we also say they are ``killed'' by $\kappa$). If all terms are set to 0, then $\phi\upharpoonright\kappa  =0$. 
(ii) Otherwise for any term $\tau_i$ consistent with $\kappa$, we replace $\tau_i$ with $\tau_i \upharpoonright \kappa := \langle \pi_i \setminus \sigma, E_i \setminus U\rangle$. If $\tau_i \upharpoonright \kappa = \langle \emptyset, \emptyset \rangle$ this means $\kappa$ already satisfies this term; if this happens for any of the $\tau_i$ then we set $\phi \upharpoonright \kappa=1$. 
\end{definition}

The last basic computational model acting on inputs $f\in \mathcal{F}_{N,M}$ we consider is a ``pigeonhole decision tree'':

\begin{definition}
A pigeonhole decision tree $T$ over $\mathcal{F}_{N,M}$ is defined by a rooted tree with fanout $\leq M$, with leaves labeled by values from some finite set $\mathcal{Z}$. Each internal node $v \in T$ is labelled by either a ``pigeon query'' or a ``hole query''.  Pigeon nodes (those that make a pigeon query) are labeled by a query $q(v) \in [N]$ and have $\leq M$ outgoing edges each labeled by a distinct element $(q(v), y) \in [N] \times [M]$. 
Hole nodes are labeled by a query $q(v) \in [M]$ and have $\leq N+1$ outgoing edges, which are either labelled by a  distinct element $(x,q(v)) \in [N] \times [M]$ or by $q(v) \in [M]$. (The label $q(v)$ corresponds to nothing mapping to pigeon $q(v)$.)  
%\times [M]) \sqcup [M]$. 
In this way, we can associate each node $v \in T$ to some $\langle \pi, E\rangle$ with $\pi \subseteq [N] \times [M]$, $E \subseteq [M]$, consisting of the labels of all edges on the path from the root to $v$ in $T$. We then require that this pair $\langle \pi, E\rangle$ associated to $v$ is a hole disjunction.
% Furthermore, $T$ is required to satisfy the 1-1 property: every path $p$ in $T$ must be consistent with some 1-1 function in ${\cal F}_{N,M}$. That is, associated with every path in $T$ is a partial restriction -- the set of edges fixed so far along this path -- and we require this restriction to be a hole restriction. 

$T$ is said to be a complete pigeonhole decision tree if for every $f \in {\cal F}_{N,M}$ there is a (unique) root to leaf path in $T$ that is consistent with $f$. In the case that $T$ is complete we may associate $T$ with a function $T: \mathcal{F}_{N,M} \to \mathcal{Z}$, where $T(f)$ is the value at the leaf of $T$ consistent with $f$.
% \todo{(Oliver): Can I add here: when $T$ is complete we can associate it to a function $\mathcal{F}_{N,M} \to \mathcal{Z}$?}
% \todo{Toni: yes good idea. maybe we can just make $Z={0,1}$ and then can also say when complete it fully specifies a function from $cal_F$ to 0,1}
The depth of $T$ is the length of the longest root to leaf path. We will sometimes define pigeonhole decision trees without having specific associated leaf values in mind; in this case we refer to the tree as ``unlabeled.''

%When a path is traversed from the root to some node in $T$, when we pass from $v$ to its child 
%The leaf nodes are then labeled by a value in $\mathcal{Z}$. 


%Given a pigeonhole decision tree $T$, each 1-1 assignment $f$ uniquely defines a root to leaf path in $T$ as follows:
%\begin{enumerate}
   % \item At node $v$, if $q(v) \in [N]$ traverse the edge labeled by $f(q(v))$.
   % \item At node $v$, if $q(v) \in [M]$ and $q(v) \in \rng(f)$, then traverse the edge labeled by the (unique) $u \in [N]$ with $f(u)=q(v)$; otherwise traverse the edge labeled by $\bot$.
%\end{enumerate}
%In this way we can identify $T$ with a function $T: \mathcal{F}_{N,M} \to \mathcal{Z}$ given by
%\[
%T(f) \mapsto \text{ value at leaf reached by } f \text{ in }T
%\]
% \todo{(Oliver) We are not using the leaf values anywhere here, and are not explaining how $T$ defines a function $\mathcal{F}_{N,M} \to \mathcal{Z}$ when it is complete. That's why I put the comment above}
\end{definition}

\noindent {\bf Note:} The queries in a pigeonhole decision tree \emph{do not} correspond directly to queries of the underlying bit variables $\mathrm{BITS}_{N,M}$. While it is possible to query $\log M$ of the variables $\mathrm{BITS}_{N,M}$ to determine $f(x)$ for some $x \in [N]$, if we want to determine the preimage of $y \in [M]$ under $f$ (or determine it has no preimage), we would need to query $\approx N \log M$ variables of $\mathrm{BITS}_{N,M}$ to make this determination directly. Instead, a pigeonhole decision tree corresponds more directly to a special kind of $\mathsf{P^{NP}}$ decision tree on the bit variables.



%Next we will define a procedure for representing a matching disjunction by a {\it canonical} complete pigeonhole decision tree. We start with some preliminary definitions.

%\begin{definition}[Complete Pigeonhole Trees on a subset $A$] 
%For a partial matching $\pi$ and a hole restriction $\tau$, we say that $\tau$ covers $\pi$ if $\nodes(\tau) \supseteq \nodes(\pi)$. If $\tau = \langle \sigma, E\rangle$, we say that $\tau$ is a minimal cover of $\pi$ if it fails to be a cover after removing any edge from $\sigma$ or any element from $E$.

%For every tuple $(A,N,M)$ where $A \subseteq [N] \sqcup [M]$, 
%we define a canonical complete pigeonhole decision tree, denoted denoted by $\mathcal{C}(A, N,M)$, as follows. This tree is unlabelled, and is defined with respect to some fixed ordering of the elements of $A$. So lets assume the ordering $a_1,\ldots,a_r$ of $A$. 
%We build $\mathcal{C}(A,N,M)$ in $r$ steps:
%\begin{enumerate}
   % \item In step 1, if $a_1 \in [N]$, we query pigeon $a_1$ at the root, with all possible outedges: $(a_i,y)$, $y \in [M]$. Otherwise if $a_1 \in [M]$, we query hole $a_1$ with all possible outedges: $(x,a_i), x \in [N] \sqcup \{\bot\}$. 
    %\item In step $i>1$, for each leaf node $l$ in $T$, let $\tau_l = <\sigma_l,E_l>$ denote the hole restriction associated with the path to $l$.  If $a_i$ is a pigeon or hole in $nodes(\tau_l)$ (that is, if $a_i$ has already been "determined" by the hole restriction thus far), then we do nothing.  Otherwise,  we query $a_i$ at the leaf $l$, with all possible outedges that are consistent with $\tau_l$. 
   % \end{enumerate}
%\end{definition}
    
%Observe that the leaves of $\mathcal{C}(A,N,M)$ are in one-to-one correspondence with the set of all hole restrictions that minimally cover $A$, and thus ${\cal C}(A,N,M)$ is a complete pigeonhole decision tree.

%\begin{definition}[Canonical Pigeonhole Trees]
%Let $\phi = \pi_1 \lor \ldots \lor \pi_{s}$ be a matching disjunction over $[N] \sqcup [M]$  where we have fixed this ordering on the terms. 
%The canonical pigeonhole tree for $\phi$, 
%$\mathcal{T}({\phi},N,M)$, is defined recursively as follows:

%\begin{enumerate}
   % \item If $s=0$, then $\phi$ is a constant function. If $\phi=1$ then $\mathcal{T}({\phi},N,M)$ consists of a single node labelled by 1, and similarly if $\phi$ is 0, then the decision tree is a single node labelled by 0.
    %\item Otherwise, let $\pi_1$ be the first term in $\phi$ according to the ordering of terms in $\phi$. First, we create the complete pigeonhole tree ${\cal C}(nodes(\pi_1),N,M)$. For each leaf $l$ in the tree constructed so far, let $\tau_l = <\sigma_l,E_l>$ be the hole restriction associated with the path to $l$. Then for each leaf $l$: 
    %\begin{enumerate}
       % \item Let $P \subseteq [N]$ be the subset of pigeons that have not been matched  by $\pi_l$, and let $H \subseteq [M]$ be the subset of holes not matched by $\pi_l$ and not in $E_l$. $P$ and $H$ define the pigeons and holes that are still undetermined after applying $\tau_l$.  If the size of $P$ is $N'$ and the size of $H$ is $M'$, then we can rename the vertices so that we identify $P$ with $[N']$ and $H$ with $[M']$.
        %\item Recursively construct ${\cal T}(\phi\upharpoonright \tau_l ,N',M')$. and attach a copy of this subtree to the leaf $l$.
    %\end{enumerate}
   % \end{enumerate}
%\end{definition}

% Note that if $T$ is a complete pigeonhole tree for a matching disjunction $\phi$ \todo{(oliver): haven't defined "pigeonhole tree for a matching disjunction''. toni: yes we can remove this now with addition above}, then for every 1-1 function $f \in {\cal F}_{N,M}$, there is a unique path  in $T$ that is consistent with $f$, and the value assigned to the leaf of this path is equal to $\phi(f)$.  Furthermore, if $T$ has depth $d$, then both $\phi$ and $\neg \phi$ can be expressed as width $w$ hole disjunctions:
We next observe that if $T$ is a complete depth $d$ pigeonhole decision tree with binary leaf values, then we may represent both $T$ and $\lnot T$ by a width $d$ hole disjunction:
\begin{observation}\label{obs:dt-to-disj}
Let $F: \mathcal{F}_{N,M} \to \B$ be some predicate. If $F \equiv T$ for some complete pigeonhole decision tree $T$ of depth $d$, then there exist hole disjunctions $\phi_1,\phi_2$ of width $\leq d$ so that $F \equiv \phi_1$ and $\lnot F \equiv \phi_2$.
\end{observation}

%We initialize with $T_{\phi}$ initially set to be a single 
%\begin{enumerate}
    %\item  If $\phi =1$ then $\mathcal{T}_{\phi}$ consists of a single node labelled by 1. Similarly if $\phi =0$, then $\mathcal{T}_{\phi}$ consists of a single node labelled by 0. 
   % \item Otherwise, we construct $\mathcal{C}_{\nodes(\pi_1)}$.
    %recall that each leaf is associated to a minimal cover $\tau$ of $\pi_1$. 
    %At the unique such leaf associated to $\tau = \pi_1$ we label the leaf by the value 1. At any leaf $\tau$ such that $\tau \nparallel \pi_1$ we label it with the value $0$. Otherwise we construct recursively $\mathcal{T}_{\phi \upharpoonright \tau}$ and attach a copy of this subtree to the current leaf associated with $\tau$.
     %Then for each leaf of $l$ of ${\mathcal C}_{\nodes(\pi_1)}$, we apply the above  procedure recursively on the restricted matching disjunction $\phi|_{\rho_l}$, where $\rho_l$ is the partial matching associated with leaf $l$.} 
    % \end{enumerate}
%\end{definition}

%We will deal at various points with a hole restriction or partial matching $\rho$ which we have fixed, and wish to consider only those assignments which extend this partial information. In this context we will need to talk about the value of various functions/formula on $\mathcal{F}_{N,M}$ when our attention is restricted to those assignments compatible with $\rho$:
%\begin{definition}[Formulae Relative to a Hole Restriction]
%A pair of general functions $F,G: \mathcal{F}_{N,M} \to \B$ and some hole restriction $\rho$, we say that $F \equiv_{\rho} G$ if $F(f) = G(f)$ for all $f \in \mathcal{F}_{N,M}$ consistent with $\rho$.

%In the case that $F$ is computed by a hole disjunction $\phi$ and $\rho$ is some hole restriction, we would like to simplify $\phi$ so that it contains only the relevant information necessary to decide the value $\phi(f)$ on 1-1 assignments $f$ extending $\rho$. For maps $\tau= \langle \pi, E\rangle$ and $\rho = \langle \sigma, U \rangle$, we denote by $\tau \upharpoonright \rho:= \langle \tau \setminus \rho , E \setminus U\rangle$ in the case that $\tau, \rho$ are compatible, or else $\tau \upharpoonright \rho:= \bot$ if they are incompatible. 
%For a hole disjunction $\phi = \lor_i \tau_i$, we define $\phi \upharpoonright \rho$ the disjunction $\lor_i (\tau_i \upharpoonright \rho)$, where each term $\tau_i$ with $\tau_i \upharpoonright \rho = \bot$ is removed, and if any term has $\tau_i \upharpoonright \rho =\langle \emptyset, \emptyset\rangle$ then we set $\phi \upharpoonright \rho = 1$ (in this case the disjunction is already satisfied by the hole restriction).
%\end{definition}

We will need the following variant of the Switching Lemma, which says that for any  low-width matching disjunction $\phi$, is we sample a random partial matching $\rho$, then with high probability $\phi\upharpoonright\rho$ will have a low-depth, complete pigeonhole decision tree. We defer the proof of the Switching Lemma to Section ~\ref{sec:php-sl}.
%for $\phi \upharpoonright  \rho$ will have low depth. 
%we can find a low depth pigeonhole decision tree $T$ so that $T \equiv_{\rho} \phi$; in other words after applying the restriction prescribed by $\rho$, we can replace $\phi$ with a low depth pigeonhole decision tree. 
%In the case of our depth 3 lower bound, the use of this will be due to the following lemma, which lets us replace such a decision tree \todo{by dnf}.

\begin{definition}[Distribution of Partial Restrictions] 
\label{}
Let $\mathcal{M}^{N,M}_K$ be the set of all partial matchings $[N] \to [M]$ with exactly $N - K$ edges. We use $\rho \sim \mathcal{M}^{N,M}_K$ to denote a sample from the uniform distribution on this set. When $N,M$ are clear from context we write $\mathcal{M}_K$ as shorthand for $\mathcal{M}^{N,M}_K$.  
% \todo{(Oliver) I suppose we have to call it $\mathcal{M}^{N,M}_K$ here since the definition depends on all 3 parameters. Toni: Yup but maybe we can omit M,N when clear from context?}
\end{definition} 

\begin{lemma}[Pigeonhole Switching Lemma]
\label{lem:switching-main} 
{\it Let} $M,N,d \in \mathbb{N}$. 
%$L = M - N$. 
{\it Let} $\phi$ {\it be a width}-$w$ {\it matching disjunction over} $[M]^{[N]}$. 
%Say that $\rho$ is sampled uniformly from the set of partial matchings $[N] \to [M]$ using $N-K$ edges. 
{\it If} $M-N \leq K \leq \frac{N}{4}$ {\it and} $N,K,d,w$ {\it sufficiently large, then}
\[
\Pr_{\rho \sim \mathcal{M}_K} [\phi\upharpoonright \rho \text{ has no depth } \leq d \text{ pigeonhole decision tree }] \leq \exp\Bigl(d (\log wK^5 -\log N^{1/2} + O(1))\Bigr)\]
\end{lemma}

%\todo{(Oliver:) Still not sure if this should be here or the intro... some of these comments won't make sense until we've defined the relevant concepts like pigeonhole dts} 
%The proof of this lemma will be given in the Appendix. 
%Like other Switching Lemmas, we use a coding style argument first introduced by Razborov \todo{raz} (see also \todo{beame}, \todo{add references}). Switching lemmas have been used for both circuit lower bounds for computing functions, and proof complexity lower bounds for deriving tautologies. Between the two, our argument bears a stronger resemblance to the second, however there are some key conceptual differences. The largest difference is that we require our pigeonhole decision trees to be equivalent to the hole disjunction that they represent, over all 1-1 input functions. In contrast, the pigeonhole decision trees used in proof complexity switching lemmas make queries to a pigeon or a hole at each vertex, and every path in the tree corresponds to a partial matching; because of this, if the number of pigeons and holes is unbalanced, the behavior of such trees is not actually well-defined on every total assignment $f: [N] \to [M]$. Here we have to allow for the possibility that a queried hole is unmapped: in our case when a hole is queried, we have $N+1$ possible outgoing edges, one for each pigeon that could map to this hole, plus an extra edge corresponding to the case where nothing maps to this hole. Another crucial distinction is how we can use the Switching Lemma to flip a DNF to a CNF. In the classical proof complexity variants, we are able to flip a $\lor\land$ formula of some type to a $\land\lor$ formula of the same type; in particular if we start with a matching disjunction we can flip it to (the compliment of) another matching disjunction. In our case we are only able to flip a matching to the compliment of a \emph{hole disjunction}, and for this reason our lemma cannot be applied twice. This is not a defect of our method, but rather a necessary feature of any technique here, since our search problem \emph{can be solved} by circuits of one higher depth. 
%With these appropriate modifications, our proof of the switching lemma is similar to previous proofs.

\subsection{Depth 3 Lower Bound for Strong Avoid}
\label{sec:depth3-lb}

In this section we prove the following circuit lower bound, which is a restatement of Theorem ~\ref{thm:strongavoid-lb}.
%\todo{This next thm should be a restatement of thm in intro, with same theorem number}

\medskip

%\begin{theorem}\label{thm:depth-3-lb}
\noindent {\bf Theorem~\ref{thm:strongavoid-lb} (Restated)}. {\it There is some absolute constant $\epsilon> 0$ so that the following holds. Let $\Phi_1,\ldots,\Phi_k$ be depth 3, size $\leq s$ unbounded fanout formulas over the bit variables $\mathrm{BITS}_{N,N+1}$. Provided $k \leq N^{\epsilon}$ and $s \leq 2^{N^{\epsilon}}$, there exists a string $z \in \B^k$ so that for all $y \in [N+1]$, there exists an assignment $f: [N] \rightarrow [N+1]$ such that:
\begin{enumerate}
    \item $\Phi_i(f) = z_i$ for all $i$
    \item $y \in \mathrm{range}(f)$
\end{enumerate}
In particular it is not possible to determine an empty pigeonhole of $f$ by reading the values of $\Phi_1(f),\ldots,\Phi_k(f)$.}
%\end{theorem}

\medskip

\noindent {\bf Overview of Proof.}
As mentioned in Section ~\ref{sec:intro-proofoverview}, we want to prove Theorem ~\ref{thm:strongavoid-lb} by restricting attention to the class of inputs that correspond to 1-1 functions, and in order to do this we will need to prove a strengthening of 
%However if we truly restrict ourselves to 1-1 functions $f \in {\cal F}_{N,N+1}$, then a lower bound is not possible by Lemma~\ref{lem:1-1-lop-unique}: for these inputs, there is a unique solution, and therefore polynomial-size depth-3 circuits can easily check whether the unique solution has its $i^{th}$ bit equal to 1 or to 0 and thus solve the problem unconditionally.
%To circumvent this barrier, we will prove a strengthening of 
Theorem~\ref{thm:strongavoid-lb}, by proving the lower bound for Avoid on input functions $f: [N] \rightarrow [M]$, where $M$ is larger than $N+1$. 
%By enlarging the range of $f$, we can focus our attention of $f$'s that are 1-1 since now for every input, there are at least $M-N$ distinct solutions.
As discussed in Section ~\ref{sec:intro-proofoverview}, this implies the lower bound stated above (with $M=N+1$).
%since there is a direct reduction from Avoid on instances $[N] \to [M]$ to instances $[N] \to [N+1]$: we simply map every element of $[M] \setminus [N+1]$ to the element $N+1$. Observe that this reduction does not preserve injectivity.
 
Thus we assume for sake of contradiction that there exist depth-3 AND-of-OR-of-AND circuits $\Phi_1, \ldots,\Phi_k$ of total size at most $s$ that solves Avoid on all one-to-one instances $f: [N_0] \rightarrow [M_0]$, where $M_0 = N_0 + N_0^{\epsilon}$, for some $0<\epsilon<1$.

% Since we will prove the lower bound with respect to all 1-1 inputs, by Observation~\ref{obs:matching-conversion} we can assume without loss of generality that the bottom level of ANDs are partial matchings. Therefore, each circuit $\Phi_i$ is without loss of generality an AND of matching disjunctions.
\begin{enumerate}
 % \item The first step is to argue that there exists a matching restriction $\rho$ such that after applying $\rho$, all depth-2 subcircuits (which are matching disjunctions) have width at most $w \approx \log s$. This is accomplished in the Width Reduction Lemma (\ref{lem:width-red}) by a standard Chernoff argument and a union bound. Thus after applying Lemma ~\ref{lem:width-red},  we are left with depth-3 circuits $\Phi_1,\ldots,\Phi_k$ which are ANDs of matching disjunctions of width at most $w = c \log s$ for some constant $c>0$. %N_0^{\delta}$ for some $0 < \delta < 1$. 
 % These circuits still solve Avoid but now with respect to 1-1 functions on the reduced domain and range, $[N_1]$, $[M_1]$,  where $N_1 =\Omega(N_0)$, and $M_1 \approx N_1 + N_1^{\epsilon}$.

 \item The first step is to argue that there exists a matching restriction $\rho$ such that after applying $\rho$, all depth-2 subcircuits (which are DNFs) have pigeon-width at most $w \approx \log s$. This is accomplished in the Width Reduction Lemma (\ref{lem:width-red}) by a standard Chernoff argument and a union bound. Composing this with Observation~\ref{obs:matching-conversion}, we can then convert all of the depth two subcircuits to \emph{matching disjunctions} of the same width $w$, without affecting their behavior on 1-1 assignments. Overall, we are left with depth-3 circuits $\Phi_1,\ldots,\Phi_k$ which are ANDs of matching disjunctions of width at most $w = c \log s$ for some constant $c>0$, and there size is at most $s M^{c \log s}$.
 These circuits still solve Avoid but now with respect to 1-1 functions on the reduced domain and range, $[N_1]$, $[M_1]$,  where $N_1 =\Omega(N_0)$, and $M_1 \approx N_1 + N_1^{\epsilon}$.
 \item At this point each $\Phi_i$ is an AND of low-width matching disjunctions, i.e. we have $\Phi_i = \land_j \phi_{i,j}$ where the $\phi_{i,j}$ are width $w$ matching disjunctions. The next step is to apply our Pigeonhole Switching Lemma (\ref{lem:switching-main}) which will tell us that there exists a matching restriction $\rho$ such that for all $i,j$,  $\phi_{i,j}\upharpoonright\rho$ simplifies to a depth $d$ pigeonhole decision tree. We will be able to choose the parameters in our switching lemma so that after the second restriction, there are $N_2$ remaining pigeons and $M_2= 2N_2$ remaining holes, with $d^2k^2 < N_2$. By Observation~\ref{obs:dt-to-disj}, this allows us to rewrite $\lnot \phi_{i,j}$ as a width=$d$ hole disjunction, and thus $\lnot \Phi_i$ as width $d$ hole disjunction. The simplified circuits $\lnot \Phi_1,\ldots, \lnot \Phi_k$ still solve Avoid, but now with respect to 1-1 functions on the reduced domain and range leftover after this restriction, i.e. on instances $[N_2] \to [2N_2]$.
 
 \item The final step is to prove that if $\phi_1,\ldots,\phi_k$ are width-$d$ hole disjunctions that solve Avoid with respect to all 1-1 functions from $[N_2]$ to $[2N_2]$, then we must have $d^2k^2 > N_2$.
 This is accomplished by Lemma ~\ref{thm:1.5-lb}. The proof is a novel argument based on coverings. Recall that our goal is to give some sequence of values $z_1, \ldots, z_k \in \B$, so that for all $y_0 \in [2N_2]$ we can find a 1-1 instance $f: [N_2] \to [2N_2]$ so that $\phi_1(f)=z_1, \ldots,\phi_k(f)=z_k$, but $y_0$ is not an Avoid solution for $f$. This indicates that the values $\phi_1(f), \ldots,\phi_k(f)$ are insufficient to determine an Avoid solution. To construct $z_1,\ldots,z_k$, we start by repeatedly searching for some $\phi_i$ whose underlying terms have a \emph{small hitting set}, which is a small set of pigeons and holes so that every term of $\phi_i$ mentions at least one of them. If some $\phi_i$ is found, we apply a partial restriction $\rho$ to the hitting set variables, which reduces the width of $\phi_i$ by 1.
 Since there are $k$ $\phi_i$'s, and each has width at most $d$, after $dk$ iterations, each $\phi_i$ either has been set to a constant, or is promised to have no small hitting set. We set $z_i = 1$ or $z_i = 0$ for all the $\phi_i$ which have been forced to a constant by our restriction $\rho$ in this process, and set $z_i=1$ for all the unkilled $\phi_i$. Then to finish the proof, we need to show that for any $y_0$ we can find an instance $f \Vert \rho$ so that $y_0 \in \rng(f)$, and $\phi_i(f)=1$ for every $\phi_i$ which was unkilled in the first step. Let $y_0 \in [2N_2]$ be given. Starting with the partial assignment $\rho$, we extend it so some unmapped pigeon $x_0$ goes to $y_0$, and then we iterate over each unkilled $\phi_i$ and try extend $\rho$ to satisfy one of its terms greedily. The correctness will follow from the fact that none of the unkilled $\phi_i$ have a small hitting set.
 \end{enumerate}
 
 We now formalize the above, which relies on the three Lemmas mentioned in the proof overview. 
 %lemmas which are justified subsequently. 
 Two of these Lemmas (Lemmas ~\ref{lem:width-red} and ~\ref{thm:1.5-lb}) will then be proven, while the more involved Lemma \ref{lem:switching-main} is postponed to Appendix~\ref{sec:php-sl}. 
 
\begin{proof}[Proof of Theorem ~\ref{thm:strongavoid-lb}]
Fix some $\epsilon, \delta$ to be specified later. Let $N_0$ be sufficiently large, and let $M_0 = N_0 + N_0^{\epsilon}$.
%$M > N$, for some value of $M$ to be specified later. 
%\todo{L is used below to be number of unset pigeons after switching lemma. Maybe just leave out previous sentence.} \todo{(Oliver): I am using the same parameter $L$ for both the gap $M-N$ and the number of unset pigeons, that's why. }
Assume towards a contradiction that $\Phi_1,\ldots,\Phi_k$ are depth-3, size $s = 2^{N_0^{\epsilon/2}}$ circuits so that for all assignments $f: [N_0] \to [M_0]$, we may determine an empty pigeonhole for $f$ by reading the values $\Phi_1(f),\ldots,\Phi_k(f)$, and that $k \leq N_0^{\epsilon/5}$. Without loss of generality we may assume:
\[
\Phi_i = \bigwedge_{j} D_{i,j}
\]
where $D_{i,j}$ is a DNF over the variables $\mathrm{BITS}_{N_0,M_0}$. Obviously if $\Phi_1(f),\ldots,\Phi_k(f)$ could determine an empty pigeonhole on all assignments $f$ then it would do so on all 1-1 assignments; we will derive a contradiction already from this fact.

% Therefore by Observation~\ref{obs:matching-conversion} we may assume
% \[
% \Phi_i = \bigwedge_{j} \phi_{i,j}
% \]
% where $\phi_{i,j}$ is a matching disjunction, at the cost of increasing the overall size of $\Phi$ from $s$ to $M_0 s$, while preserving the behavior of the $\Phi_i$ on 1-1 assignments.

Applying the width reduction lemma (Lemma~\ref{lem:width-red})
with $w = c \log s$ for a sufficiently large constant $c>1$, and a union bound over $i,j$, we can find a partial matching $\rho_0$ 
leaving $N_1$ pigeons unmapped, where $N_1 =  \delta N_0 $, 
%with $N - \hat{N}$ edges (i.e. leaving $\hat{N}$ pigeons unmapped), so that $\hat{N} = \Omega(N)$ 
such that for all $i,j$, $D_{i,j} \upharpoonright \rho_0$ has pigeon-width $w$. We then convert each $D_{i,j}$ to matching disjunction $\phi_{i,j}$ using Observation~\ref{obs:matching-conversion}; the width of the matching disjunctions is at most $w$, and their number is at most $s':= s \cdot M_0^{c \log s} \leq 2^{N_0^\epsilon}$.
%Now, say we choose $L = {\hat{N}}^{\epsilon}$ for some $\epsilon > 0$, and have that $k, \log s \leq L^{\frac{1}{5}}$. 
At this point each circuit $\Phi_i$ has been reduced to a conjunction of width $w$ matching disjunctions of size $s'$. The reduced circuits still solve Avoid with respect to all 1-1 functions $f$ with domain  $N_1= \delta N_0 $ and range $M_1 = N_1 + N_0^{\epsilon}$.
%Then for $\epsilon = bla$ is a sufficiently small constant,
%(not depending on $N$) 

Next we apply our Switching Lemma (Lemma~\ref{lem:switching-main}) with the following choice of parameters:
%domain 
%$$\epsilon = 1/10$$
$$N:= N_1$$
$$M:=M_1 = N_1 + N_0^{\epsilon}$$
$$K:=N_0^{\epsilon}$$
$$w = c \log s = c N_0^{\epsilon}$$
$$d:=K^{1/5} = N_0^{\epsilon/5} $$



Choosing $\epsilon$ sufficiently small, it follows that for all $i,j$:
$$Pr_{\rho \sim {\cal M}_K}[{\cal T}_{\phi_{i,j}\upharpoonright\rho}~ has ~depth~\geq d] < \frac{1}{ks'}$$ 
Thus by a union bound over $i,j$, we conclude that
there exists a partial matching $\rho$ 
%\rho:=\rho_0\rho_1$ 
so that for each $i,j$ there is a pigeonhole decision tree $T_{i,j}$ of depth $d$
%:=\log S = N_1^{1/4}$.
%L^{\frac{1}{5}}$ 
such that $\phi_{i,j} \equiv_{\rho} T_{i,j}$.

Applying Observation~\ref{obs:dt-to-disj} we can find hole disjunctions $\psi_{i,j}$ of width $d$ so that $\lnot \phi_{i,j} \equiv_{\rho} \psi_{i,j}$, and therefore $\lnot \Phi_i \equiv_{\rho} \bigvee_j \psi_{i,j}$ which is a width $d$ hole disjunction. By the conditions on the original $\Phi_1,\ldots,\Phi_k$, for all 1-1 assignments $f$ extending $\rho$ we must be able to determine an empty pigeonhole of $f$ by reading the values $\lnot \Phi_1(f),\ldots,\lnot \Phi_k(f)$. 

After applying $\rho$, we are left with a reduced domain and range $[N_2]$ and $[M_2]$, where $N_2 = K = N_0^{\epsilon}$, and $M_2 = M_1 - (N_1 - K) = 2K = 2N_2$.
%Clearly the set of 1-1 assignments extending $\rho$ are in direct correspondence with the set of 1-1 maps $[N_2] \to [2N_2]$. %\todo{Explain/fix parameters here -- how to set them  so that we are left with map from [L] to [2L].} 
Therefore on the smaller input size $[2N_2]^{[N_2]}$, we have found a sequence of $k$ width-$d$ hole disjunctions such that on all 1-1 assignments $g: [N_2] \to [2 N_2]$, we can determine an empty pigeonhole of $g$ by reading the values of these hole disjunctions applied to $g$. At this point we reach a contradiction with Theorem~\ref{thm:1.5-lb}, since by our choice of parameters we have $4k^2d^2 < N_2$ is satisfied (recall that $k \leq N_0^{\epsilon/5}$ by assumption).%which yields a contradiction to our choice of $k,d,L$.
% Recall that $w=c\log s$ for a sufficiently large constant $c$, which by our setting of $w=N_1^{\epsilon}$ implies that with $\epsilon=1/10$, and the circuit size $s$ is at least size $exp(\Omega(N^{1/10})$. 
 

\end{proof}

Equipped with the Pigeonhole Switching Lemma (Lemma~\ref{lem:switching-main}, proof in Section ~\ref{sec:php-sl}), it is left to prove Lemmas ~\ref{lem:width-red} and ~\ref{thm:1.5-lb}.   We will first prove Lemma 
~\ref{thm:1.5-lb} followed by a proof of Lemma ~\ref{lem:width-red}. 
%(Width Reduction Lemma) and Lemma ~\ref{thm:1.5-lb} (the lower bound for solving Avoid by non-adaptively querying hole disjunctions). We start with the latter:

% \begin{definition}[Generalized Proposition]
% With $f$ as above we say that $\tau(f)$ is a ``generalized proposition'' if either:
% \begin{enumerate}
%     \item $\tau(f) = \mathbbm{1}\{f(x)=y\}$ for some fixed $x\in N, y \in M$
%     \item $\tau(f) = \mathbbm{1}\{y \notin \mathrm{range}(f)\}$ for some fixed $y \in M$
% \end{enumerate}

% We say that $T$ is a generalized term if $T = \land_{i \leq w} \tau_i$ where $\tau_i$ are generalized propositions; in this case the width of the term $T$ is $w$. Finally we say that $\psi$ is a generalized hole disjunction if $\psi = \lor_i T_i$ where $T_i$ are generalized terms; in the case that all $T_i$ have width $w$ we say that $\psi$ overall has width $w$.
% \end{definition}
% \todo{switch $s$ to $k$ here}
\begin{lemma}\label{thm:1.5-lb}
Let $\phi_1,\ldots,\phi_k$ be hole disjunctions of width $w$ over $[2N]^{[N]}$. Then provided $N > 4k^2w^2$, there exists a string $z \in \B^k$ so that for all $y \in [2N]$, there exists a 1-1 function $f: [N] \rightarrow [2N]$ such that:
\begin{enumerate}
    \item $\phi_i(f) = z_i$ for all $i$
    \item $y \in \mathrm{range}(f)$
\end{enumerate}
In particular it is not possible to determine an empty pigeonhole of $f$ from $\phi_1(f),\ldots,\phi_k(f)$ for all 1-1 assignments $f$.
\end{lemma}

To prove this we need the following definition:
\begin{definition}[Hitting Sets]
If $\phi$ is a hole disjunction and $A \subseteq [N] \sqcup [M]$, we say that $A$ is a hitting set for $\phi$ if $A \cap \nodes(\tau)  \neq \emptyset$ for all terms $\tau \in \phi$. 
\end{definition}
We can observe the following property of hitting sets:
\begin{observation}\label{obs:width-killing}
Let $\pi: [N] \to [M]$ be a partial matching, $\phi$ a nonempty hole disjunction, and $A\subseteq [N] \sqcup [M]$ a non-empty hitting set for $\phi\upharpoonright\pi$.  Say that $\pi'$ is an extension of $\pi$ such that $A \subseteq \nodes(\pi')$. Then for each term $\tau \in \phi\upharpoonright\pi$, $\tau$ has strictly smaller width in $\phi\upharpoonright\pi'$ then it originally did in $\phi\upharpoonright\pi$. 
\end{observation}
We can now prove the main claim:
\begin{proof}
To prove this we will construct partial matching $\pi: [N] \to [2N]$, a set $I \subseteq [k]$, and values $\{z_i \in \B \mid i \in I\}$ so that:
\begin{enumerate}
    \item For all total 1-1 assignments $f$ extending $\pi$, $\phi_i(f) = z_i$ for all $i \in I$.
    \item For all $y_0 \in [2N]$, there exists a total 1-1 assignment $f$ extending $\pi$ so that $\phi_i(f)=1$ for all $i \notin I$, and $y_0 \in \mathrm{range}(f)$.
\end{enumerate}
We will initialize $\pi = \emptyset$ and expand it in stages by the following procedure:
\begin{enumerate}
    \item If there exists a nonempty $A \subseteq [N] \sqcup [2N], |A| \leq 2kw + 2$ and some $i\in [k]$ so that $A$ is a hitting set for $\phi_i \upharpoonright \pi$, then extend $\pi$ to a 1-1 map $\pi'$ so that $A \subseteq \nodes(\pi')$. If none exist then halt the procedure and output the current $\pi$.
    \item Set $\pi := \pi'$ and return to the previous step.
\end{enumerate}
We claim that we will always be able to extend $\pi$ in the appropriate way until no further hitting sets can be found, and that at the end we will have $|\pi| \leq kw(2kw+2)=2k^2w^2 +2kw$. To see this, observe that if $\phi$ is a hole disjunction of width $w$ and $\nodes(\pi)$ is a hitting set for $\phi$, then every term in $\phi$ will either be killed in $\phi\upharpoonright\pi$ or else have its width decreased by 1 by Observation~\ref{obs:width-killing}. Therefore the above procedure can only repeat $kw$ times before all the $\phi_i$ have been killed. In each step at most $2kw+2$ edges need to be added to $\pi$ in order to cover $A$, so the total size of $\pi$ at any step is at most $kw(2kw+2)=2k^2w^2 +2kw$, and therefore there are always enough available pigeons/holes to extend $\pi$ in the appropriate way at each step since $N > k^2w^2$.

Now we can choose $I \subseteq [k]$, which will consist of those indices $i$ so that $\phi_i\upharpoonright\pi$ is forced to a constant (i.e. one of its terms is already satisfied by $\pi$ in which case it is forced to one, or else all of its terms are killed and it is forced to 0). For $i \in I$ we denote by $z_i \in \B$ the value it is forced to. It remains to show that for any given $y_0 \in [2N]$, we can find a total 1-1 assignment $f$ extending $\pi$ so that $y_0 \in \mathrm{range}(f)$ and $\phi_i(f)=1$ for all $i \notin I$. To do this we will construct a hole restriction $\langle \rho, U\rangle $, where $\rho$ is another partial matching disjoint consistent with $\pi$, and $U \subseteq [2N]$ satisfying:
\begin{enumerate}
    \item $y_0 \in \mathrm{range}(\pi \cup \rho)$
    \item $|\rho| \leq kw$, $|U| \leq kw$
    \item $U \cap \mathrm{range}(\pi \cup \rho) = \emptyset$
    \item For each $i \in [k] \setminus I$ there is a term $\langle \sigma, E\rangle \in \phi_i$, so that for all $x \in \dom(\sigma)$ we have $\pi \cup \rho(x) = \sigma(x)$, and $E \subseteq U$.
\end{enumerate}

If we can accomplish this then the proof is complete. We simply extend $\pi \cup \rho$ to a total 1-1 assignment $f$ which leaves $U$ out of its range. By the assumption $N > 4k^2w^2$ such a total assignment exists; here we are using the fact that the total number of holes $2N$ is larger then the size of $U$ plus the number of holes filled thus far by $\pi \cup \rho$. By construction we then have $\phi_i(f)=1$ for all $i \notin I$ and $y \in \mathrm{range}(f)$.

Observe that by construction of $\pi$, for each $i \notin I$ we have that $\phi_i\upharpoonright\pi$ has no hitting set of size $\leq 2kw+2$. We will construct $\rho, U$ in stages. Initially we check if $y_0 \in \mathrm{range}(\pi)$, if not we initialize $\rho = \{(x_0,y_0)\}$ where $x_0$ is an arbitrary element unmapped in $\pi$, otherwise we initialize $\rho = \emptyset$. 
In addition we initialize $U = \emptyset$. Now we go through each $i \in [k] \setminus I$ in order and do the following. We search for some term $\tau = \langle \sigma, E\rangle \in \phi_i \upharpoonright (\pi \cup \rho)$ so that $\nodes(\tau) \cap \nodes(\langle \rho, U\rangle) = \emptyset$. If found, then we add the nodes in $E$ to $U$ and the edges in $\sigma$ to $\rho$. We claim that it is always possible to find such a term while maintaining that $\pi \cup \rho$ is 1-1 and $U \cap \mathrm{range}(\pi \cup \rho) = \emptyset$. In particular say that we have gotten to some step $i \in [k] \setminus I$ where this is not possible. Let $A = \nodes(\langle \rho, U\rangle)$. Observe that by construction $|A|\leq 2kw+2$. Recall that by construction of $\pi$ we have that $A$ cannot be a hitting set for $\phi_i\upharpoonright\pi$. By definition of a hitting set this implies the existence of the required term $\tau \in \phi_i$.
\end{proof}

Next, we prove the Width Reduction Lemma:

\begin{lemma}[Width Reduction]\label{lem:width-red}
Let $\epsilon > 0$ be a sufficiently small constant. Let $\phi$ be a DNF over $\mathrm{BITS}_{N,M}$ with $s$ terms. Say a partial matching $\rho$ is sampled as follows: sample $A \subseteq [N]$ by including each element independently with probability $1-\epsilon$; now choose a uniform 1-1 function $A \to [M]$.
\[
\Pr_{\rho}[D\upharpoonright \rho \text{ has pigeon width } \geq w] < \exp(\log s - \Omega(w))
\]
\end{lemma}
\begin{proof}
Let $\phi = \lambda_1,\ldots,\lambda_s$. We will ignore each term whose pigeon width is already $<w$; its pigeon width cannot increase under the restriction. Let $\lambda \in \phi$ be a term of width $\geq w$. We will bound the probability that $\rho$ does not kill the term $\lambda$ and then take a union bound over the $s$ terms. Without loss of generality we may assume $\lambda$ is the literal $f_{x,1}$; if it were of the form $f_{x,i}$ or $\lnot f_{x,i}$ for some $i > 1$ the argument is completely symmetric. Let $B \subset M$ consist of those $y\in M$ whose first bit (in binary) is 1. 
\begin{gather*}
\Pr_{\rho}[\lambda \text{ survives }\rho] = \sum\limits_{U \subseteq \dom(\lambda)} \Pr[\dom(\lambda) \cap \dom(\rho) = U] \cdot \bigl(\frac{|B|}{M}\bigr)^{|U|} \\
\leq \sum_{t=1}^w\Pr[|\dom(\lambda) \cap \dom(\rho)| = t] \cdot 2^{-t}\\
\leq 2^{-w/2}+\Pr[|\dom(\lambda) \cap \dom(\rho)| < \frac{w}{2}] \\
\leq 2^{-w/2} + \Pr_{\xi_1,\ldots,\xi_w \sim \mathrm{Bern}(1-\epsilon)} [\sum_i \xi_i \leq \frac{w}{2}] \leq \exp(- \Omega(w))
\end{gather*}
Where the last inequality follows from Chernoff's bound and $\mathrm{Bern}(1-\epsilon)$ denotes the Bernoulli distribution with expected value $1-\epsilon$. Thus by a union bound 
\[
\Pr_{\rho}[\phi\upharpoonright \rho \text{ has width } \geq w] \leq \exp(\log s - \Omega(w))
\]
\end{proof}


\paragraph{Lower Bounds for Moderately Weak Avoid:}
Above we showed that depth-3 circuits cannot solve the Strong Avoid problem unless their size is exponential. In contrast, Li's construction gives \emph{quasipolynomial} size depth 3 circuits which solve the Weak Avoid problem, i.e. Avoid on instances $[N] \to [2N]$ when the codomain is at least twice as large as the domain. By a standard iteration construction (the same argument which lets us reduce $[N] \to [2N]$ Avoid to $[N] \to [N^2]$ Avoid as in \cite{Korten}), for every constant $d$ we can obtain quasipolynomial size depth 3 circuits solving Avoid on instances $[N] \to [N+N/\log^dN]$. We show here that this is essentially optimal: for every constant $c$ there exists another constant $d$ so that $[N] \to [N+N/\log^dN]$ Avoid cannot be solved by depth 3 circuits of size $2^{\log^cN}$. The proof is a simple reduction to Theorem~\ref{thm:strongavoid-lb} which is originally due to Razborov \cite{raz-talk}.
\begin{lemma}\label{lem:php-lb-larger-stretch}
For every $c \in \mathbb{N}$ there exists $d \in \mathbb{N}$ so that the following holds. Let $\Phi_1,\ldots,\Phi_k$ be depth 3 circuits of size $s$ such that for every assignment $f:[N] \to [N+N/\log^dN]$, an Avoid solution can be determined from the values $\Phi_1(f),\ldots,\Phi_k(f)$. Then either $k \geq \log^cN$ or $s \geq 2^{\log^c N}$.
\end{lemma}
\begin{proof}
Let $c \in \mathbb{N}$ be fixed. We will choose $M = \log^dN$ for an appropriate choice of $d$ which will be specified later and which will depend only on $c$. We then reduce the Avoid problem on instances $[M] \to [M+1]$ to instances $[N] \to [N+N/\log^dN]$ as follows: given $g: [M] \to [M+1]$, let $g^{\otimes \ell}: [M\ell] \to [(M+1)\ell]$ be the map which sends each block $[1,M], \ldots, [M(\ell-1), M\ell]$ to its corresponding block in $[1,M+1], \ldots [(M+1)(\ell-1),(M+1)\ell]$ according to $g$. Setting $\ell = \frac{N}{M}$ we see that $g$ is a map $[N] \to [N + \ell]$, where $\ell = N/M = N/\log^d N$. 

Obviously given $g$ we can generate $g^{\otimes \ell}$ without any computational overhead by simply substituting variables. In addition, any Avoid solution to $g^{\otimes \ell}$ uniquely determines an Avoid solution to $g$, obtained by simply forgetting the block number and outputting the position within the relevant block. Thus if the circuits $\Phi_1(f),\ldots,\Phi_k(f)$ could solve Avoid on inputs $[N] \to [N+N/\log^dN]$ then they can also solve Avoid on inputs $[M] \to [M+1]$. By Theorem~\ref{thm:strongavoid-lb} there is some absolute constant $\epsilon$, such that for this to be possible we must have $k \geq M^{\epsilon}$ and $s \geq 2^{M^{\epsilon}}$. Setting $d > \frac{c}{\epsilon}$ the theorem follows.
\end{proof}

% \begin{lemma}
% \label{lem:php-lb-larger-stretch}

% Suppose there exists size $s$, depth-$3$ unbounded fanin circuits $C_1\ldots,C_k$ solving
% Avoid on functions $g$ from $[N]$ to $[M]$, where $M=N + N/B$, $B \in [N]$, and $B$ divides $N$.
% Then there exists size $s$, depth-$3$ unbounded fanin circuits $C_1',\ldots,C_k$ solving Strong Avoid on functions $f$ from $[B]$ to $[B+1]$.
% \end{lemma}

% \begin{proof}
% Assume that circuits $C_1,\ldots,C_k$ solve Avoid on inputs $f: [N] \rightarrow [M]$, where $M=N/B$. Let $f$ be an input to Strong Avoid, where $f$ maps $[B]$ to $[B+1]$.
% We create an instance, $g: [N] \rightarrow [N + N/B]$ by dividing $[N]$ into $N/B$ blocks each of size $B$ and dividing $[N+N/B]$ into $N/B$ blocks each of size $B+1$. Then we map the pigeons in the $i^{th}$ block to holes in the $i^{th}$ block using $f$. 
% That is, for $i \in \{0,\ldots,N/B-1\}$ and $j \in \{1,\ldots,B\}$,
% let $g(iB+j) = i(B+1) + f(j)$.
% It is easy to see that a solution to Avoid on input $g$ implies a solution to Strong Avoid on input $f$. 
% \end{proof}

% Applying Theorem ADD, with $B=polylog(N)$, we obtain the following 
% quasipolynomial lower bound for Avoid where the ``stretch'' of $f$ is $M-N = N/poly\logN$.
% We note that this is not far off from the upper bounds given by ADD and ADD 
% who construct quasipolynomial-size depth-3 circuits for Avoid with the same stretch.

% \begin{corollary}
% Let ${\cal C}$ be a depth-3 circuit solving Avoid on functions from
% domain $[N]$ to range $[N+N/polylog(N)]$. 
% Then ${\cal C}$ must have quasipolynomial size.
% \end{corollary}


\subsection{\texorpdfstring{$\mathsf{BPP^{NP}}$}{BPP\^NP} Lower Bound for 1-1 Strong Avoid}
\label{sec:BPP-NP-lb}
We restate the main result to be proved:

\medskip
\noindent{\bf Theorem~\ref{thm:bpp-np-lb}}
%\begin{theorem*}[The0rem~\ref{thm:bpp-np-lb}]
{\it In the decision tree model, Strong 1-1 Avoid is not solvable in $\mathsf{FBPP^{NP}}$.}
%\end{theorem*}
% \todo{Toni: restate original theorem stating BPP-NP lower bound and discuss reformuation below.} 

\medskip
Recalling the definition of $\mathsf{BPP^{NP}}$ (in the decision tree model), if 1-1 Strong Avoid has $\mathsf{BPP^{NP}}$ decision trees of complexity $w$, then there is a distribution $\mathcal{T}$ over $\mathsf{P^{NP}}$ decision trees of complexity $r$, so that for all $f: [N] \to [N+1]$ $\mathcal{T}$ outputs a 1-1 Strong Avoid Solution with probability $\geq 2/3$. By Yao's minimax principle, this implies the existence of a fixed tree $T \in \mathrm{support}(\mathcal{T})$, so that $T(f)$ succeeds in finding an empty pigeonhole of $f$ with probability $\geq 2/3$ when $f$ is sampled uniformly from the space of 1-1 functions $\mathcal{F}_{N,N+1}$. It thus suffices to rule out the existence of such a deterministic tree $T$. We will use here the notation $\mathcal{F}_N := \mathcal{F}_{N,N+1}$, and $f \sim \mathcal{F}_N$ to denote a uniform sample from this set. Thus Theorem ~\ref{thm:bpp-np-lb} follows from the following Theorem:

\begin{theorem}\label{thm:bpp-np-lb-dt}
There is an absolute constant $\epsilon > 0$ so the following holds: if $T$ is a $\mathsf{P^{NP}}$ decision tree of complexity $N^{\epsilon}$ over the bit variables $\mathrm{BITS}_{N,N+1}$ and leaves labeled by elements of $[N+1]$, then:
\[
\Pr_{f \sim \mathcal{F}_N} [T(f) \notin \rng(f)] \leq N^{-\Omega(1)}
\]

\end{theorem}

The high level idea of the proof is as follows. Observe that a random $f \sim \mathcal{F}_{N,M}$ can be sampled by first choosing a uniform partial matching $\rho$ of a certain size, and then choosing a uniform completion of $\rho$ to a total assignment. Applying our Switching Lemma, we can argue that with high probability, after we sample $\rho$ we can replace all of the $\mathsf{NP}$ queries in our $\mathsf{P^{NP}}$ decision tree $T$ with small depth pigeonhole decision trees, which overall allows $T$ to be replaced by a low depth pigeonhole decision tree. It then remains only to argue that a pigeonhole decision tree of low depth cannot solve 1-1 Strong Avoid with non-trivial probability on a uniform extension $f$ of $\rho$, which can be accomplished with a direct argument.
\begin{proof}[Proof of Theorem ~\ref{thm:bpp-np-lb-dt}]
By definition of a $\mathsf{P^{NP}}$ decision tree, $T$ is a binary tree of depth $r = N^{\epsilon}$ with each internal node branching on the value of some DNF of width $r$ over the bit variables $\mathrm{BITS}_{N,N+1}$. Let $D_1,\ldots,D_s$ be the set of all DNF associated to the nodes of $T$; so $s \leq 2^r$. As in the proof of Theorem~\ref{thm:strongavoid-lb} we may apply Observation~\ref{obs:matching-conversion} and replace all $D_i$ by width $r\log N$ matching disjunctions $\phi_i$ since we only care about their behavior on assignments in $\mathcal{F}_N$. Fix some parameter $K$ to be specified later. We will consider sampling $f \sim \mathcal{F}_N$ in the following way: first choose a uniform partial matching $\rho$ with $N-K$ edges, then a uniform extension $f \supseteq \rho$ of $\rho$ to a total 1-1 assignment. We define $\mathcal{B}$ as the event that $\phi_i \upharpoonright \rho$ does not have a pigeonhole decision tree of depth $\leq r$ for some $i \leq s$. Observe in the case $\rho \notin \mathcal{B}$, we may construct a pigeonhole decision tree $\mathcal{T}$ of depth $r^2$ so that $\mathcal{T} \equiv_{\rho} T$; we simply simulate the original computation of $T$, and each time a DNF $D_i$ was originally queried, we instead simulate the depth $\leq r$ pigeonhole decision which represents $\phi_i$. Let $\mathbb{P}_d(J)$ be the maximum, over all pigeonhole decision trees $Q$ of depth $d$ on $g: [J] \to [J+1]$, of $\Pr_{g \sim \mathcal{F}_J}[Q(g) \notin \rng(g)]$. Then we have: 

\begin{gather*}
\Pr_{f \sim \mathcal{F}_N} [T(f) \notin \rng(f)] = \Pr_{\rho,~f \supseteq \rho } [T(f) \notin \rng(f)] \\ \leq \Pr_{\rho}[\rho \in \mathcal{B}] + \max_{\rho \notin \mathcal{B}}\Pr_{f \supseteq \rho}[T(f) \notin \rng(f)]\\
\leq \Pr_{\rho}[\rho \in \mathcal{B}] + \mathbb{P}_{r^2}(K)\\
\leq 2^r \max_i \Pr_{\rho}[\phi_i\upharpoonright \rho \text{ has no pigeonhole decision tree of depth } \leq r] + \mathbb{P}_{r^2}(K)
\end{gather*}
We will prove in the next lemma that $\mathbb{P}_{r^2}(K) \leq \frac{r^2}{K+1}$. Setting $\epsilon$ sufficiently small and $K = r^3$, we can apply Lemma~\ref{lem:switching-main} and complete the proof.
\end{proof}

It remains to prove the bound on $\mathbb{P}_{r^2}(K)$. At the cost of increasing the depth of the tree by 1, we assume that a pigeonhole decision tree always queries a hole before outputting it as a solution.

\begin{lemma}\label{lem:pdt-lb}
$\mathbb{P}_d(N) = \frac{d}{N+1}$
\end{lemma}
\begin{proof}
Let $T$ be a tree witnessing $\mathbb{P}_d(N)$. 

Say that $T$ first queries a hole $y \in N$, so that 
\[
T(f) = \begin{cases}
T_1(f) \text{\;\;if } f(1) = y\\
\cdots \\
T_{N}(f) \text{ if } f(N) = y\\
y \text{\;\;\;\;\;\;\;\; if } y \notin \rng(f)
\end{cases}
\]
then we have:
\begin{gather*}
\mathbb{P}_d(N) = \Pr_{f}[T(f) \notin \rng(f)] \\
= \Pr_f[y \notin \rng(f)] + \sum\limits_{x}\Pr_f[f(x)=y] \cdot \Pr_f[T_x(f) \notin \rng(f) \mid f(x)=y] \\
\leq \frac{1}{N+1} + \frac{N}{N+1} \mathbb{P}_{d-1}(N-1)
\end{gather*}

Here we are using that fact that a uniform $f$ conditioned on $f(x)=y$ is equivalent, up to relabeling, to a uniform member of $\mathcal{F}_{N-1}$, and the labeling of pigeons and holes has no effect on the value of $\mathbb{P}_{d-1}(N-1)$. By the same reasoning if $T$ first queries a pigeon $x$ then we have the simpler bound:

\begin{gather*}
\Pr_{f}[T(f) \notin \rng(f)] = \sum\limits_{y} \Pr_f[f(x)=y] \cdot \Pr_f[T_y(f) \notin \rng(f) \mid f(x)=y]
\leq  \mathbb{P}_{d-1}(N-1)
\end{gather*}

so overall
\[
\mathbb{P}_d(N) \leq \max\{\mathbb{P}_{d-1}(N-1),\frac{1}{N+1} + \frac{N}{N+1} \mathbb{P}_{d-1}(N-1)\}
\]
Therefore by induction on $d$ we conclude that the optimal pigeonhole decision tree achieving $\mathbb{P}_d(N)$ only queries holes, and without loss of generality queries them in order $1,\ldots, d$ (in the base case we use the assumption that $T$ must query a hole before using it as an answer), then outputs the index of the first empty hole that was queried (if any). Therefore we have:
\[
\mathbb{P}_d(N) = \Pr_f[\{1,\ldots,d\} \nsubseteq \rng(f)] = \frac{d}{N+1}
\]
\end{proof}

%\begin{corollary}

Next we show that Theorem ~\ref{thm:sigma2-pi2-separation}, separating the decision class $\mathsf{\Sigma_2^P \cap \Pi_2^P}$ from $\mathsf{BPP^{NP}}$, follows as a corollary of Theorem \ref{thm:bpp-np-lb}.

%In the decision tree model, $\mathsf{\Sigma_2^P} \cap \mathsf{\Pi_2^P} \not\subseteq \mathsf{BPP^{NP}}$.
% \end{corollary}
%\begin{proof}

% \medskip

\begin{proof}[Proof of Theorem~\ref{thm:sigma2-pi2-separation}]
% Recall that Theorem ~\ref{thm:sigma2-pi2-separation} states that  in the decision tree model, 
% $\mathsf{\Sigma_2^P} \cap \mathsf{\Pi_2^P} \not\subseteq \mathsf{BPP^{NP}}$. 
Lemma ~\ref{lem:1-1-lop-unique} says that Strong 1-1 Avoid is solvable in $\mathsf{FP^{\Sigma_2^P \cap \Pi_2^P}}$. This result holds relative to every oracle (and thus in the decision tree model). Hence relative to every oracle, if $\mathsf{\Sigma_2^P} \cap \mathsf{\Pi_2^P} \subseteq \mathsf{BPP^{NP}}$ then Strong 1-1 Avoid collapses into $\mathsf{BPP^{NP}}$. The result then follows from Theorem~\ref{thm:bpp-np-lb}.
\end{proof}
%\end{proof}



% We now need the following simpler lemma:
% \begin{lemma}
% Let $\epsilon > 0$ be a sufficiently small constant. Let $\phi$ be a matching disjunction with $s$ terms over $[M]^{[N]}$. Say a partial matching $\rho$ is sampled as follows: sample $A \subseteq [N]$ by including each element independently with probability $1-\epsilon$; now choose a uniform 1-1 function $A \to [M]$.
% \[
% \Pr_{\rho}[\phi\upharpoonright \rho \text{ has width } \geq w] < \exp(\log s - \Omega(w))
% \]
% \end{lemma}
% \begin{proof}
% Let $\phi = \lambda_1,\ldots,\lambda_s$. We will ignore each term whose width is already $<w$; its width cannot increase under the restriction. Let $\lambda \in \phi$ be a term of width $\geq w$. We bound the probability that $\rho$ does not kill the term $\lambda$ as:
% \begin{gather*}
% \Pr_{\rho}[\lambda \text{ survives }\rho] = \sum\limits_{U \subseteq \dom(\lambda)} \Pr[\dom(\lambda) \cap \dom(\rho) = U] \frac{1}{|U|! \binom{M}{|U|}} \\
% \leq \sum_{t=1}^w\Pr[|\dom(\lambda) \cap \dom(\rho)| = t] \frac{1}{t! \binom{M}{t}}\\
% \leq \frac{1}{(w/2)! \binom{M}{w/2}} +\Pr[|\dom(\lambda) \cap \dom(\rho)| < \frac{w}{2}] \\
% \leq \frac{1}{(w/2)! \binom{M}{w/2}} + \Pr_{\xi_1,\ldots,\xi_w \sim \mathrm{Bern}(1-\epsilon)} [\sum_i \xi_i \leq \frac{w}{2}] \leq \exp(- \Omega(w))
% \end{gather*}
% Thus by a union bound 
% \[
% \Pr_{\rho}[\phi\upharpoonright \rho \text{ has width } \geq w] \leq \exp(\log s - \Omega(w))
% \]
% \end{proof}
% Taking another union bound we get:
% \begin{corollary}
% There are absolute constants $\epsilon > 0$, $C \in \mathbb{N}$ so that the following holds. Say $\phi_1,\ldots,\phi_s$ are a collection of matching disjunctions, each of size at most $s$, over $[M]^{[N]}$. Then there exists a partial assignment $\rho$ of size $N-K$, $K\geq \epsilon N$ so that for all $\phi_i$, $\phi_i \upharpoonright \rho$ has width at most $C\log S$.
% \end{corollary}
% Finally, composing our two random restriction arguments we have:
% \begin{theorem}\label{thm:full-switching}
% There is some absolute constant $\epsilon>0$ so that the following holds. Let $\phi_1,\ldots,\phi_s$ be a collection of matching disjunctions of size $\leq s$ over $f: [N] \to [M]$ and let $L = M-N$. Provided $K, d \leq N^{\epsilon}$ and $s \leq 2^{N^{\epsilon}}$, and $K \geq L$, there exists a partial matching $\rho$ with $K$ unset pigeons, so that $\phi_i \upharpoonright \rho$ has a pigeonhole decision tree of depth $\leq d$ for all $i \leq s$.
% \end{theorem}



% \section{[OLD] Pigeonhole stuff}

% \todo{(Toni) Notation/defns: Notation for mapping disjunctions $\phi$ or $D$?  Generalized hole disjunction? Define $\phi|_\rho$ where $\phi$ is a hole disjunction, and also $T|_\rho$ where $T$ is a pigeonhole decision tree. Define when two $T$'s are (functionally) equivalent and also define the function ${\cal T}$ that maps a decision tree to a hole disjunction.}
% In the following we will be concerned with a function variable $f: [N] \to [M]$, where $M = N+L$, $L > 0$; we will only be concerned with assignments to $f$ which are 1-1 functions. \todo{(Toni) I think it would be good to explain all of this in terms of our underlying functional encoding of a 1-1 assignment and also a partial 1-1 assignment. i.e.,  we encode with variables $p_{i,x}$, $ i \in [N], x \in \B^{\log M}$ and vars $q_j$ and intended meaning. And then what a partial assignment looks like. This will help to understand the next defns. Also can we add to defns 1,2 how they connect to variables in our encoding? Ie. a partial matching corresponds to what kind of restriction/conjunction and same for a hole restriction.}
 
% \begin{definition}[Partial Matchings and hole restrictions]
% A partial matching $\pi$ is a partially defined 1-1 function from $[N]$ to $[M]$. We use $\dom(\pi), \rng(\pi)$ to refer to the domain and range of $\pi$, and $\nodes(\pi)= \dom(\pi) \sqcup \rng(\pi)$. More generally a ``hole restriction'' $\tau = \langle \pi, E\rangle$ consists of a partial matching $\pi$ and a set $E \subseteq [M]$ satisfying $\rng(\pi) \cap E = \emptyset$. Let $\nodes(\tau) = \nodes(\pi) \cup E$. We use $|\tau|$ to refer to the value $|\pi| + |E|$. \todo{(Toni) Explain that $E$ is a subset of $[M]$, which we view as restricting to 1-1 fxns f such that the range of f does not intersect $E$. Also later you use terminology $\rho$ and $U$ so let's unify notation.} 

% For two hole restrictions $\tau_1=\langle \pi_1,E_1\rangle, \tau_2=\langle \pi_2,E_2\rangle$, we say that the maps agree, written $\tau_1 \Vert \tau_2$, if $\langle \pi_1 \cup \pi_2, E_1\cup E_2\rangle$ is a hole restriction. We say $\tau_1$ extends $\tau_2$, written $\tau_1\supseteq \tau_2$, if $\pi_1 \supseteq \pi_2$ and $E_1 \supseteq E_2$. \todo{(Toni) I would prefer to define when two hole restrictions are consistent instead of defining when one extends another.} We will think of a partial matching $\pi$ as a special kind of hole restriction of the form $\langle \pi, \emptyset\rangle$, and a total 1-1 assignment $f$ as a special hole restriction of the form $\langle f, \corng(f)\rangle$, where $\corng(f)= [M] \setminus \rng(f)$. In this way we will use the above terminology to define relations between matchings, maps and total 1-1 assignments, e.g. $f \Vert \pi$ and $\pi \subseteq \tau$. 
% \end{definition}

% \begin{definition}[Matching and hole disjunctions]
% A hole disjunction $\phi = \lor_i \tau_i$ is defined by a collection of hole restrictions, which we refer to as the ``terms'' of $D$. We say $\phi(f)=1$, or ``$f$ satisfies $\phi$,'' if $f \Vert \tau_i$ for some $i$. If $|\tau_i|\leq w$ for all $i$ we say that $\phi$ has width $w$. We call $\phi$ a matching disjunction if all $\tau_i$ are partial matchings.

% For a hole restriction $\tau = \langle \pi, E \rangle$ and a partial matching $\sigma$, we denote by $\tau \upharpoonright \sigma$ the \todo{simplified}. Similarly for a hole disjunction $\phi$ we use $\phi \upharpoonright \sigma$ \todo{simplified disjunction}.
% \end{definition}
% \subsection{Depth 3 Lower Bound for Strong Avoid}

% \begin{definition}
% A pigeonhole decision tree $T$ over the function variable $f$ is defined by a rooted tree with unbounded fanout \todo{(Toni) define in beginning what function variables corresponding to $f: N \rightarrow M$ is. Also maybe add unbounded fanout $\leq M$?} . Each internal node $v \in T$ is labeled by a query $q(v) \in [N] \sqcup [M]$. If $q(v) \in [N]$ then $v$ has $M$ outgoing edges each labeled by an element of $M$, and if $q(v) \in [M]$ then $v$ has $N+1$ outgoing edges each labeled by an element of $[N] \sqcup \{\bot\}$. The leaf nodes are then labeled by a binary value. The depth of $T$ is the length of the longest root to leaf path.
% \todo{(Toni) I think to avoid confusion, we need a brief discussion about how the queries are related to our variables in encoding of $f$, and why they look like they do.}

% Given a pigeonhole decision tree $T$, each 1-1 assignment $f$ uniquely defines a root to leaf path in $T$ as follows:
% \begin{enumerate}
%     \item At node $v$, if $q(v) \in [N]$ traverse the edge labeled by $f(q(v))$.
%     \item At node $v$, if $q(v) \in [M]$ and $q(v) \in \rng(f)$, then traverse the edge labeled by the (unique) $u \in [N]$ with $f(u)=q(v)$; otherwise traverse the edge labeled by $\bot$.
% \end{enumerate}
% In this way we can identify $T$ with a boolean function $T: \mathcal{F}^{\leftrightarrow}_{N,M} \to \B$, where $\mathcal{F}^{\leftrightarrow}_{N,M}$ is the set of 1-1 functions from $[N]$ to $[M]$, given by 
% \[
% T(f) \mapsto \text{ value at leaf reached by } f \text{ in }T
% \]
% \todo{(Toni) Can we remove the "if-and-only-if" superscript since everywhere $f$ has this form?}
% For a pair of general boolean functions $F,G: \mathcal{F}^{\leftrightarrow}_{N,M} \to \B$ and some partial matching $\rho: [N] \to [M]$, we say that $F \equiv_{\rho} G$ if $F(f) = G(f)$ for all 1-1 assignments $f$ extending $\rho$. \todo{(Toni) Change $F$, $G$ to $T$, $T'$?}
% \end{definition}
% We make the following observation:
% \begin{observation}
% Let $F: \mathcal{F}^{\leftrightarrow}_{N,M} \to \B$ and $\rho$ a partial matching. If $F \equiv_{\rho} T$ for some pigeonhole decision tree of depth $w$, then there exist hole disjunctions $D_1,D_2$ of width $\leq w$ so that $F \equiv_{\rho} D_1$ and $\lnot F \equiv_{\rho} D_2$. \todo{(Toni) May be easier to just state this for empty restriction and then also point out that if two trees are equivalent then they are still equivalent under a partial matching restriction.}
% \end{observation}

% We will crucially use the following lemma whose proof is deferred to a later section:
% \todo{(Toni) Formally define space of restrictions we are sampling from. Also I would prefer to stick with standard way of stating the switching lemma, which is to define a function that maps any $\phi$ to a matching decision tree, and then state switching lemma similar to how it is stated in Paul's primer. }
% \begin{lemma}[Switching Lemma for Avoid]
% Let $M,N$. Let $D$ be a size $\leq \ell$ matching disjunction over the variable $f: [N] \to [M]$. Say that $\rho$ is sampled uniformly from the set of partial 1-1 functions $[N] \to [M]$ using \todo{certain number of edges}. Then with probability \todo{something}, there exists a pigeonhole decision tree $T$ of depth \todo{something} so that $D \equiv_{\rho} T$.
% \end{lemma}

% We can now prove the main result of this section:
% \begin{theorem}
% Let $N, M$ \todo{satisfy some conditions} and let $\Phi_1,\ldots,\Phi_k$ be depth 3, size $\leq s$ unbounded fanout formulas over the boolean variables $\{f_{x, b}\}_{x \in N, b \in \lceil \log M \rceil}$; each such assignment defines a function $f: [N] \to [M]$, where $f(x)$ is given in binary by $f_{x,1},\ldots,f_{x,\lceil \log M \rceil}$. Then provided \todo{quantitative conditions}, there exists a string $z \in \B^k$ so that for all $y \in M$, there exists a 1-1 function $f: N \rightarrow M$ such that:
% \begin{enumerate}
%     \item $\Phi_i(f) = z_i$ for all $i$
%     \item $y \in \mathrm{range}(f)$
% \end{enumerate}
% In particular it is note possible to determine an empty pigeonhole of $f$ by reading the values $\Phi_1(f),\ldots,\Phi_k(f)$.
% \end{theorem}
% \begin{proof}
% Assume towards a contradiction that $\Phi_1,\ldots,\Phi_k$ are given so that for all assignments $f$ we may determine an empty pigeonhole for $f$ by reading the values $\Phi_1(f),\ldots,\Phi_k(f)$. Without loss of generality we may assume:
% \[
% \Phi_i = \bigwedge_{j} \lambda_{i,j}
% \]
% where $\lambda_{i,j}$ is a DNF over the variables $\{f_{x, b}\}_{x \in N, b \in \lceil \log M \rceil}$. Obviously if $\Phi_1(f),\ldots,\Phi_k(f)$ could determine an empty pigeonhole on all assignments $f$ then it would do so on all 1-1 assignments, so it suffices to restrict attention to 1-1 assignments. In light of this we claim that we may further assume 
% \[
% \Phi_i = \bigwedge_{j} \phi_{i,j}
% \]
% where $\phi_{i,j}$ is a matching disjunction, at the cost of increasing the overall size of $\Phi$ from $s$ to $Ms$. In particular, say that $\lambda = \lor_u \land_v \ell_{u,v}$ is a DNF where each $\ell_{u,v}$ is a literal on the variables $\{f_{x, b}\}_{x \in N, b \in \lceil \log M \rceil}$. For each term $t_u = \land_v \ell_{u,v}$ we may replace it by the matching disjunction:
% \[
% \theta_u = \bigvee_{\pi \in \mathrm{match}(t_u)} \pi
% \]
% where $\mathrm{match}(t_a)$ is the set of minimal partial matchings $\pi$ so that for each literall $\ell_{u,v} = (\lnot)^{\xi} f_{x,b}$ in $t_a$ (with $\xi \in \B$), there is an edge $(x,y) \in\pi$ with $y_b = \xi$. Clearly we have that $\theta_u(f) = t_u(f)$ for all 1-1 assignments $f$, and thus we may replace $\lambda$ with the matching disjunction $\lor_u \theta_u$ without affecting its behavior on 1-1 assignments. Observe that $|\mathrm{match}(t_u)| \leq M |t_u|$.

% At this point we have each $\Phi_i$ in the form
% \[
% \Phi_i = \bigwedge_{j} \phi_{i,j}
% \]
% with $\phi_{i,j}$ a matching disjunction. We now apply Lemma \todo{ref}, and select a random 1-1 assignment $\rho$ of size \todo{something}. By a union bound, with nonzero probability there exist for each $i,j$ a pigeonhole decision tree $T_{i,j}$ of depth \todo{depth} so that $\phi_{i,j} \equiv_{\rho} T_{i,j}$. Applying the \todo{obvervation} we can find hole disjunctions $\psi_{i,j}$ of width \todo{width} so that $\lnot \phi_{i,j} \equiv_{\rho} \psi_{i,j}$, and therefore $\lnot \Phi_i \equiv_{\rho} \bigvee_j \psi_{i,j}$ which is a width \todo{width} hole disjunction. By the conditions on the original $\Phi_1,\ldots,\Phi_k$, for all 1-1 assignments $f$ extending $\rho$ we must be able to determine an empty pigeonhole of $f$ by reading the values $\lnot \Phi_1(f),\ldots,\lnot \Phi_k(f)$, which contradicts \todo{1.5 lower bound theorem}.
% \end{proof}
% It reamins to prove the following lower bound \todo{give a better segue}

% % \begin{definition}[Generalized Proposition]
% % With $f$ as above we say that $\tau(f)$ is a ``generalized proposition'' if either:
% % \begin{enumerate}
% %     \item $\tau(f) = \mathbbm{1}\{f(x)=y\}$ for some fixed $x\in N, y \in M$
% %     \item $\tau(f) = \mathbbm{1}\{y \notin \mathrm{range}(f)\}$ for some fixed $y \in M$
% % \end{enumerate}

% % We say that $T$ is a generalized term if $T = \land_{i \leq w} \tau_i$ where $\tau_i$ are generalized propositions; in this case the width of the term $T$ is $w$. Finally we say that $\psi$ is a generalized hole disjunction if $\psi = \lor_i T_i$ where $T_i$ are generalized terms; in the case that all $T_i$ have width $w$ we say that $\psi$ overall has width $w$.
% % \end{definition}
% \todo{switch $s$ to $k$ here}
% \begin{theorem}
% Let $\psi_1,\ldots,\psi_s$ be hole disjunctions of width $w$ over the variable $f: N \rightarrow M$. Then provided \todo{quantitative condition}, there exists a string $z \in \B^s$ so that for all $y \in M$, there exists a 1-1 function $f: N \rightarrow M$ such that:
% \begin{enumerate}
%     \item $\phi_i(f) = z_i$ for all $i$
%     \item $y \in \mathrm{range}(f)$
% \end{enumerate}
% \end{theorem}

% To prove this we need the following definition:
% \begin{definition}[Hitting Sets]
% If $\phi$ is a hole disjunction and $A \subseteq [N] \sqcup [M]$, we say that $A$ is a hitting set for $\phi$ if $A \cap \nodes(\tau)  \neq \emptyset$ for all terms $\tau \in \phi$. 
% \end{definition}
% We can observe the following property of hitting sets:
% \begin{observation}\label{obs:width-killing}
% Let $\pi: N \to M$ be a partial matching, $\phi$ a nonempty hole disjunction, and $A\subseteq [N] \sqcup [M]$ a non-empty hitting set for $\phi\upharpoonright\pi$.  Say that $\pi'$ is an extension of $\pi$ such that $A \subseteq \nodes(\pi')$. Then for each term $\tau \in \phi\upharpoonright\pi$, $\tau$ has strictly smaller width in $\phi\upharpoonright\pi'$ then it originally did in $\phi\upharpoonright\pi$. 
% \end{observation}
% We can now prove the main claim:
% \begin{proof}
% To prove this we will construct partial 1-1 function $\pi: N \to M$, a set $I \subseteq [s]$, and values $\{z_i \in \B \mid i \in I\}$ so that:
% \begin{enumerate}
%     \item For all total 1-1 assignments $f$ extending $\pi$, $\phi_i(f) = z_i$ for all $i \in I$.
%     \item For all $y_0 \in M$, there exists a total 1-1 assignment $f$ extending $\pi$ so that $\phi_i(f)=1$ for all $i \notin I$, and $y_0 \in \mathrm{range}(f)$.
% \end{enumerate}
% We will initialize $\pi = \emptyset$ and expand it in stages by the following procedure:
% \begin{enumerate}
%     \item If there exists a nonempty $A \subseteq [N] \sqcup [M], |A| \leq 2sw + 2$ and some $i\in s$ so that $A$ is a hitting set for $\phi_i \upharpoonright \pi$, then extend $\pi$ to a 1-1 map $\pi'$ so that $A \subseteq \nodes(\pi')$. If none exist then halt the procedure and output the current $\pi$.
%     \item Set $\pi := \pi'$ and return to the previous step.
% \end{enumerate}
% We claim that we will always be able to extend $\pi$ in the appropriate way until no further hitting sets can be found, and that at the end we will have $|\pi| \leq sw(2sw+2)=2s^2w^2 +2sw$. To see this, observe that if $\phi$ is a generalized hole disjunction of width $t$ \todo{(Toni) width $w$?} and $\nodes(\pi)$ is a hitting set for $\phi$, then every term in $\phi$ will either be killed in $\phi\upharpoonright\pi$ or else have its width decreased by 1 by Observation~\ref{obs:width-killing}. Therefore the above procedure can only repeat $sw$ times before all the $\phi_i$ have been killed. In each step at most $2sw+2$ edges need to be added to $\pi$ in order to cover $A$, so the total size of $\pi$ at any step is at most $sw(2sw+2)=2s^2w^2 +2sw$, and therefore there are always enough available pigeons/holes to extend $\pi$ in the appropriate way at each step since \todo{quantitative condition}.

% Now we can choose $I \subseteq [s]$, which will consist of those indices $i$ so that $\phi_i\upharpoonright\pi$ is forced to a constant (i.e. one of its terms is already satisfied by $\pi$ in which case it is forced to one, or else all of its terms are killed and it is forced to 0). For $i \in I$ we denote by $z_i \in \B$ the value it is forced to. It remains to show that for any given $y_0 \in [M]$, we can find a total 1-1 assignment $f$ extending $\pi$ so that $y_0 \in \mathrm{range}(f)$ and $\phi_i(f)=1$ for all $i \notin I$. To do this we will construct a hole restriction $\langle \rho, U\rangle $, where $\rho$ is another partial matching disjoint from and compatible with \todo{(Toni) replace compatable with consistent? Also doesn't really have to be disjoint? } $\pi$, and $U \subseteq [M]$ satisfying:
% \begin{enumerate}
%     \item $y_0 \in \mathrm{range}(\pi \cup \rho)$
%     \item $|\rho| \leq sw$, $|U| \leq sw$
%     \item $U \cap \mathrm{range}(\pi \cup \rho) = \emptyset$
%     \item For each $i \in [s] \setminus I$ there is a term $\langle \sigma, E\rangle \in \phi_i$, so that for all $x \in \dom(\sigma)$ we have $\pi \cup \rho(x) = \sigma(x)$, and $E \subseteq U$.
% \end{enumerate}

% If we can accomplish this then the proof is complete. We simply extend $\pi \cup \rho$ to a total 1-1 assignment $f$ which leaves $U$ out of its range; by the assumption that $L = M-N \geq sw$ and \todo{another quantitative condition} such a total assignment exists, and by construction we have $\phi_i(f)=1$ for all $i \notin I$ and $y \in \mathrm{range}(f)$.

% Observe that by construction of $\pi$, for each $i \notin I$ we have that $\phi_i\upharpoonright\pi$ has no hitting set of size $\leq 2sw+2$. We will construct $\rho, U$ in stages. Initially we check if $y_0 \in \mathrm{range}(\pi)$, if not we initialize $\rho = \{(x_0,y_0)\}$ where $x_0$ is an arbitrary element unmapped in $\pi$, otherwise we initialize $\rho = \emptyset$. 
% In addition we initialize $U = \emptyset$. Now we go through each $i \in [s] \setminus I$ in order and do the following. We search for some term $\tau = \langle \sigma, E\rangle \in \phi_i \upharpoonright (\pi \cup \rho)$ so that $\nodes(\tau) \cap \nodes(\langle \rho, U\rangle) = \emptyset$. If found, then we add the nodes in $E$ to $U$ and the edges in $\sigma$ to $\rho$. We claim that it is always possible to find such a term while maintaining that $\pi \cup \rho$ is 1-1 and $U \cap \mathrm{range}(\pi \cup \rho) = \emptyset$. In particular say that we have gotten to some step $i \in [s] \setminus I$ where this is not possible. Let $A = \nodes(\langle \rho, U\rangle)$. Observe that by construction $|A|\leq 2sw+2$. Recall that by construction of $\pi$ we have that $A$ cannot be a hitting set for $\phi_i\upharpoonright\pi$. By definition of a hitting set this implies the existence of the required term $\tau \in \phi_i$.
% \end{proof}

% \subsection{\texorpdfstring{$\mathsf{BPP^{NP}}$}{BPP\^NP} Lower Bound for 1-1 Strong Avoid}
% Recalling the definition of $\mathsf{BPP^{NP}}$ (in the decision tree model), if 1-1 Strong Avoid has $\mathsf{BPP^{NP}}$ decision trees of complexity $w$, then there is a distribution $\mathcal{T}$ over $\mathsf{P^{NP}}$ decision trees of complexity $r$, so that for all $f: [N] \to [N+1]$ $\mathcal{T}$ outputs a 1-1 Strong Avoid Solution with probability $\geq 2/3$. By Yao's minimax principle, this implies the existence of a fixed tree $T \in \mathrm{support}(\mathcal{T})$, so that $T(f)$ succeeds in finding an empty pigeonhole of $f$ with probability $\geq 2/3$ when $f$ is sampled uniformly from the space of 1-1 functions $[N] \to [N+1]$. It thus suffices to rule out the existence of such a tree $T$. In the following let $\mathcal{F}_N$ be the set of all 1-1 functions $[N] \to [N+1]$, and $f \sim \mathcal{F}_N$ denote a uniform sample.
% \begin{theorem}\label{thm:bpp-np-lb}
% There is an absolute constant $\epsilon > 0$ so the following holds: if $T$ is a $\mathsf{P^{NP}}$ decision tree of complexity $N^{\epsilon}$ over the bit variables of $f: [N] \to [N+1]$ and leaves labeled by elements of $[N+1]$, then:
% \[
% \Pr_{f \sim \mathcal{F}_N} [T(f) \notin \rng(f)] = N^{-\Omega(1)}
% \]

% \end{theorem}
% \begin{proof}
% By definition of a $\mathsf{P^{NP}}$ decision tree, $T$ is a binary tree of depth $r$ with each internal node branching on the value of some DNF of width $r$ over the bit variables of $f$. Let $D_1,\ldots,D_s$ be the set of all DNF associated to the nodes of $T$; so $s \leq 2^r$. As in the proof of Theorem\todo{ref} we may assume all $D_i$ are matching disjunctions since we only care about their behavior on assignments in $\mathcal{F}_N$. Fix some parameter $K$ to be specified later. We will consider sampling $f \sim \mathcal{F}_N$ in the following way: first choose a uniform partial matching $\rho$ with $N-K$ edges, then a uniform extension $f \supseteq \rho$ of $\rho$ to a total assignment. We define $\mathcal{B}$ as the event that $D_i \upharpoonright \rho$ does not have a pigeonhole decision tree of depth $\leq r$ for some $i \leq s$. In the case $\rho \notin \mathcal{B}$, we may represent $T\upharpoonright \rho$ by a pigeonhole decision tree $\mathcal{T}$ of depth $r^2$. Let $\mathbb{P}_d(J)$ be the maximum, over all pigeonhole decision trees $Q$ of depth $d$ on $g: [J] \to [J+1]$, of $\Pr_{g \sim \mathcal{F}_J}[Q(g) \notin \rng(g)]$. Then we have: 

% \begin{gather*}
% \Pr_{f \sim \mathcal{F}_N} [T(f) \notin \rng(f)] = \Pr_{\rho, f} [T(f) \notin \rng(f)] \\
% \leq \Pr_{\rho}[\rho \in \mathcal{B}] + \max_{\rho \notin \mathcal{B}}\Pr_{f \supseteq \rho}[T(f) \notin \rng(f)]\\
% \leq \Pr_{\rho}[\rho \in \mathcal{B}] + \mathbb{P}_{r^2}(K)\\
% \leq 2^r \max_i \Pr_{\rho}[D_i\upharpoonright \rho \text{ has depth } > r] + \mathbb{P}_{r^2}(K)
% \end{gather*}
% We will prove in the next lemma that $\mathbb{P}_{r^2}(K) \leq \frac{r^2}{K+1}$. Setting $\epsilon$ sufficiently small and $K = r^3$, we can apply Lemma~\ref{lem:switching-main} and complete the proof.
% \end{proof}

% It remains to bound $\mathbb{P}_d(K)$. At the cost of increasing the depth of the tree by 1, we assume that a pigeonhole decision tree always queries a hole before outputting it as a solution.

% \begin{lemma}\label{lem:pdt-lb}
% $\mathbb{P}_d(N) = \frac{d}{N+1}$
% \end{lemma}
% \begin{proof}
% Let $T$ be a tree witnessing $\mathbb{P}_d(N)$. 

% Say that $T$ first queries a hole $y \in N$, so that 
% \[
% T(f) = \begin{cases}
% T_1(f) \text{\;\;if } f(1) = y\\
% \cdots \\
% T_{N}(f) \text{ if } f(N) = y\\
% y \text{\;\;\;\;\;\;\;\; if } y \notin \rng(f)
% \end{cases}
% \]
% then we have:
% \begin{gather*}
% \mathbb{P}_d(N) = \Pr_{f}[T(f) \notin \rng(f)] \\
% = \Pr_f[y \notin \rng(f)] + \sum\limits_{x}\Pr_f[f(x)=y] \cdot \Pr_f[T_x(f) \notin \rng(f) \mid f(x)=y] \\
% \leq \frac{1}{N+1} + \frac{N}{N+1} \mathbb{P}_{d-1}(N-1)
% \end{gather*}

% Here we are using that fact that a uniform $f$ conditioned on $f(x)=y$ is equivalent, up to relabeling, to a uniform member of $\mathcal{F}_{N-1}$, and the labeling of pigeons and holes has no effect on the value of $\mathbb{P}_{d-1}(N-1)$. By the same reasoning if $T$ first queries a pigeon $x$ then we have the simpler bound:

% \begin{gather*}
% \Pr_{f}[T(f) \notin \rng(f)] = \sum\limits_{y} \Pr_f[f(x)=y] \cdot \Pr_f[T_y(f) \notin \rng(f) \mid f(x)=y]
% \leq  \mathbb{P}_{d-1}(N-1)
% \end{gather*}

% so overall
% \[
% \mathbb{P}_d(N) \leq \max\{\mathbb{P}_{d-1}(N-1),\frac{1}{N+1} + \frac{N}{N+1} \mathbb{P}_{d-1}(N-1)\}
% \]
% Therefore by induction on $d$ we conclude that the optimal pigeonhole decision tree achieving $\mathbb{P}_d(N)$ only queries holes, and without loss of generality queries them in order $1,\ldots, d$ (in the base case we use the assumption that $T$ must query a hole before using it as an answer), then outputs the index of the first empty hole that was queried (if any). Therefore we have:
% \[
% \mathbb{P}_d(N) = \Pr_f[\{1,\ldots,d\} \nsubseteq \rng(f)] = \frac{d}{N+1}
% \]
% \end{proof}

% \begin{corollary}
% Relative to a generic oracle, $\mathsf{\Sigma_2^P} \cap \mathsf{\Pi_2^P} \not\subseteq \mathsf{BPP^{NP}}$
% \end{corollary}
% \begin{proof}
% Consider the type-2 decision problem: given $f: [N] \to [N+1]$ and $i \in \lceil \log (N + 1) \rceil$, accept if $f$ is a 1-1 function whose unique empty pigeonhole $y$ has $y_i=1$. By definition this is a decision problem in $\mathsf{\Sigma_2^P} \cap \mathsf{\Pi_2^P}$, since we can witness $f$ is a yes-instance by supplying $y$, and that it is a no-instance by supplying either $y$ or a collision. By the above Theorem, it cannot be solved by \todo{randomized decision lists/ np oracle decision trees, however we decide to phrase this}.
% \end{proof}
% \subsection{Switching Lemma}
% \begin{definition}[Canonical Decision Tree]
% For a partial matching $\pi$ and a hole restriction $\tau$, we say that $\tau$ covers $\pi$ if $\nodes(\tau) \supseteq \nodes(\pi)$. If $\tau = \langle \sigma, E\rangle$, we say that $\tau$ is a minimal cover of $\pi$ if it fails to be a cover after removing any edge from $\sigma$ or any element from $E$.

% For a set $A \subseteq [N] \sqcup [M]$, we denote by $\mathcal{C}_A$, the complete unlabeled decision tree over $A$, the pigeonhole decision tree (without leaf labels) which queries each element of $A$ in a fixed order. Observe that the leaves of $\mathcal{C}_A$ are in one to one correspondence with minimal covers of $A$. 

% For a matching disjunction $D = \pi_1 \lor \ldots \lor \pi_{s}$ with some fixed ordering on terms, we associate a pigeonhole decision tree $\mathcal{T}_D$, referred to as the ``canonical decision tree for $D$,'' as follows: construct $\mathcal{C}_{\nodes(\pi_1)}$; recall that each leaf is associated to a minimal cover $\tau$ of $\pi_1$. At the unique such leaf associated to $\tau = \pi_1$ we label the leaf by the value 1. At any leaf $\tau$ such that $\tau \nparallel \pi_1$ we label it with the value $0$. Otherwise we construct recursively $\mathcal{T}_{D \upharpoonright \tau}$ and attach a copy of this subtree to the current leaf associated with $\tau$.
% \end{definition}
% % \todo{modified def:}
% % \begin{definition}[Canonical Decision Tree]
% % Let $\pi = (x_1,y_1), \ldots, (x_k,y_k)$ be a partial matching with some fixed ordering on its constituent edges. We define $\mathcal{T}_{\pi}$, the canonical decision tree for $\pi$, to be an unlabeled pigeonhole decision tree obtained as follows:
% % \begin{enumerate}
% %     \item The root makes the query $x_1$. At the child labeled by the response $y_1$ we place recursively a copy of $\mathcal{T}_{\pi \setminus (x_1, y_1)}$
% % \end{enumerate}
% % \end{definition}

% It follows from the definition that $\mathcal{T}_D \equiv D$. We will now prove Lemma \todo{ref} according to the usual switching-lemma strategy: we first prove a special case of the theorem when the width of the matching disjunctions are small, and then a second (and much easier) restriction argument to reduce the case of size-$s$ disjunctions to width-$O(\log s)$ disjunctions.
% \begin{lemma}
% Let $M,N,d \in \mathbb{N}$, $L = M - N$. Let $D$ be a width-$w$ matching disjunction over the variable $f: [N] \to [M]$. Say that $\rho$ is sampled uniformly from the set of partial 1-1 functions $[N] \to [M]$ using $N-K$ edges. If $L \leq K \leq \frac{N}{4}$ and $N,K,d,w$ sufficiently large, then
% \[
% \Pr_{\rho} [\mathcal{T}_{D\upharpoonright \rho} \text{ has depth } \geq d] \leq \exp\Bigl(d (\log wK^5 -\log N^{1/2} + O(1))\Bigr)
% \]
% \end{lemma}
% We first need the following definition:
% \begin{definition}
% Let $\mathrm{stars}(w,t)$ denote the set of all sequences $\beta_1\cdots\beta_{\ell}$ such that each $\beta$ is a string in $\{\star,\bot\}^w \setminus \{\bot\}^w$ and the total number of $\star$ symbols across all $\beta_i$ is at most $t$.
% \end{definition}
% \begin{lemma}\todo{ref beame}
% $|\mathrm{stars}(w,t)| \leq (w/\ln 2)^t$
% \end{lemma}
% \begin{proof}[Proof of Lemma~\ref{lem:switching-main}]
% Let $D = \lambda_1 \lor \ldots \lor \lambda_s$ be a matching disjunction of width $w$. Let $\mathcal{B}$ be the set of partial matchings $\rho$ with $N - K$ edges, such that $\mathcal{T}_{D \upharpoonright \rho}$ has depth $\geq d$. 
% If $\rho \in \mathcal{B}$, there exists a total assignment $f \supseteq \rho$ yielding a path $P$ of length $d$ from the root in $\mathcal{T}_{D \upharpoonright \rho}$. We define a sequence of hole restrictions $\tau_1,\dots,\tau_{\ell}$ and a sequence of partial matchings $\sigma_1,\ldots,\sigma_{\ell}$, and a sequence of strings $\beta_1,\ldots,\beta_{\ell} \in \{\star, \bot\}^w$ each with at least one $\star$ index as follows. Define $\tau_0 = \emptyset$. Once $\tau_i$ has been defined, let $a_{i+1}$ be the index of the first term in $D$ so that $\lambda_{a_{i+1}} \upharpoonright \rho \cup \tau_1 \cup \cdots \cup \tau_i$ is unkilled. Then let $\tau_{i+1}$ be the minimal map $\tau_{i+1} \subseteq f$ which covers $\lambda_{a_{i+1}} \upharpoonright \rho \cup \tau_1 \cup \cdots \cup \tau_i$, and let $\sigma_i = \lambda_{a_{i+1}} \upharpoonright \rho \cup \tau_1 \cup \cdots \cup \tau_i$. Finally let $\beta$ be the string with a star corresponding to each edge of $\lambda_{a_i}$ which is contained in $\sigma_i$. We claim that such a sequence can be constructed so that at the end $\ell \leq d$, and if we take $\sigma = \cup_i \sigma_i$, then $\sigma$ is a partial matching with $\nodes(\sigma) \cap \nodes(\rho) = \emptyset$ and $\frac{d}{2} \leq |\sigma| \leq dw$. This follows from the definition of $\mathcal{T}_{D \upharpoonright \rho}$ and \todo{the fact that $f$ achieves a deep decision tree, explain... the fact that $\sigma$ is a matching follows from...}.

% Let $\mathcal{M}_R$ be the set of partial matchings $[N] \to [M]$ with exactly $N - R$ edges. So $\rho \in \mathcal{M}_K$, while $\rho \cup \sigma \in \mathcal{M}_{K-t}$ for some $\frac{d}{2} \leq t \leq dw$. We will define a map 
% \[
% \mathrm{Dec}: \bigcup\limits_{t \geq \frac{d}{2}}\mathcal{M}_{K-t}  \times \mathrm{stars}(w,t) \times \Delta_t \to \mathcal{M}_K
% \]
% for some finite sets $\{\Delta_t\}_{t \leq N}$ so that whenever $\rho \in \mathcal{B}$ and $\sigma, \beta_1,\ldots,\beta_{\ell}$ are constructed from $\rho$ as above, there exists some ``advice'' $\delta \in \Delta_t$ so that $\mathrm{Dec}(\rho \cup \sigma, \delta, \beta_1\cdots\beta_{\ell}) = \rho$. This will imply that
% \[
% \Pr_{\rho \sim \mathcal{M}_{K}}[\rho \in \mathcal{B}] \leq \frac{1}{|\mathcal{M}_K|} \sum\limits_{t \geq \frac{d}{2}} |\mathcal{M}_{K-t}|  |\mathrm{stars}(w,t)| |\Delta_t|
% \]
% since it follows by definition that $\beta_1\cdots\beta_{\ell} \in \mathrm{stars}(w,t)$. This will yield the theorem provided we can choose the sets $\Delta_t$ sufficiently small.

% Let $\rho \in \mathcal{B}$, $\tau_i,a_i,\sigma_i, \beta_i$ be as above. Define $\gamma_{i}:=\rho \cup \tau_1 \cup \ldots \cup \tau_{i} \cup \sigma_{i+1} \cup \ldots \cup \sigma_{\ell}$, where $\gamma_0 = \rho \cup \sigma$. Let $\sigma=\cup_i\sigma_i$ and $t = |\sigma_i|$. First, we claim that if we are given $\tau_1,\ldots,\tau_{i-1}$, $\sigma_1,\ldots,\sigma_{i-1}$, and  and $\gamma_{i-1}$ then we may decode from these the index $a_i$. To do this, we simply search for the first satisfied term in $D \upharpoonright \gamma_{i-1}$ \todo{justify this}. Next, observe that if we have knowledge of $a_i$, then using $\beta_i$ we may reconstruct $\sigma_i$ \todo{this is direct from the definition}. Once we have $\gamma_{i-1}, \sigma_i, \tau_i$, we may construct $\gamma_i$ directly, by removing $\sigma_i$ in $\gamma_{i-1}$ and replacing it with $\tau_i$. So it remains to show how to specify $\tau$ using the extra information $\Delta_t$, so that we may recover $\tau_i$ from $\sigma_i$. Obverse that $\nodes(\tau)$ is contained in $\nodes(\sigma) \cup ([N] \sqcup [M] \setminus \nodes(\rho \cup \sigma))$. Therefore if we construct a map $\delta: [K] \to [K + L]$ with $|\delta| \leq 2d$ which specifies, for each $u \in \nodes(\sigma)$ the preimage/image of $u$ in $\tau$ (or the lack of preimage if $u$ is marked empty), then from $\sigma_i$ we may reconstruct $\tau_i$ by searching for the images/preimages of the nodes in $[K] \sqcup [K+L]$ corresponding to $\nodes(\sigma_i)$. Thus for every $t$ we can take $\Delta_t = \Delta$ to be the set of all hole restrictions $[K] \to [K+L]$ of size at most $d$, so that $|\Delta| \leq \binom{K}{2d}\binom{K+L}{2d}(2d)!$.

% It remains to estimate the value
% \begin{gather*}
% \Pr_{\rho \sim \mathcal{M}_{K}}[\rho \in \mathcal{B}] \leq \frac{1}{|\mathcal{M}_K|} \sum\limits_{t \geq \frac{d}{2}} |\mathcal{M}_{K-t}| |\mathrm{stars}(w,t)| |\Delta_t| \\
% =
% \sum\limits_{t \geq \frac{d}{2}} \frac{|\mathcal{M}_{K-d}|}{|\mathcal{M}_K|}  |\mathrm{stars}(w,t)| |\Delta_t|
% \end{gather*}

% Observe that for each $\kappa \in \mathcal{M}_{K-t}$ there are at least $\binom{N-K+t}{t}$ elements $\alpha \in \mathcal{M}_{K}$ such that $\alpha \subseteq \kappa$, and on the other hand for any $\alpha \in \mathcal{M}_{K}$ there are at most $\binom{K}{t}\binom{K+L}{t}t!$ elements $\kappa \in \mathcal{M}_{K-t}$ so that $\alpha \subseteq \kappa$. Thus
% \[
% \frac{ |\mathcal{M}_{K-t}|}{|\mathcal{M}_K|} \leq \frac{\binom{K}{t}\binom{K+L}{t}t!}{\binom{N-K+t}{t}} 
% \]

% So overall we have:

% \begin{gather*}
% \Pr_{\rho \sim \mathcal{M}_{K}}[\rho \in \mathcal{B}] \leq \sum\limits_{t \geq \frac{d}{2}}
% \frac{\binom{K}{t}\binom{K+L}{t}t!}{\binom{N-K+t}{t}} (w/\ln 2)^t |\Delta_t|\\
% \leq |\Delta|\sum\limits_{t \geq \frac{d}{2}}
% \frac{\binom{K}{t}\binom{2K}{t}t!}{\binom{N-K+t}{t}} (w/\ln 2)^t \\
% \leq 
% \binom{K}{2d}\binom{K+L}{2d}(2d)!\sum\limits_{t \geq \frac{d}{2}}
% \frac{\binom{K}{t}\binom{2K}{t}t!}{\binom{N-K+t}{t}} (w/\ln 2)^t \\
% \leq (\frac{2K}{\sqrt{2d}})^{4d}\sum\limits_{t \geq \frac{d}{2}} \left(\frac{4wK^2}{N}\right)^t
% \end{gather*}
% Where we apply the assumptions $K \geq L$ and $K\leq \frac{N}{4}$. We may safely assume $\frac{wK^2}{N} = o(1)$, otherwise the conclusion of the theorem holds trivially. Thus by geometric decay of terms in the above sum we have, for some absolute $C \in \mathbb{N}$, the bounds:

% \begin{gather*}
% \Pr_{\rho \sim \mathcal{M}_{K}}[\rho \in \mathcal{B}] \leq O (\frac{2K}{\sqrt{2d}})^{4d}\left(\frac{4wK^2}{N}\right)^{\frac{d}{2}} \\
% \leq (2K)^{4d} \left(\frac{4wK^2}{N}\right)^{\frac{d}{2}}
% \end{gather*}
% provided $K, d$ are sufficiently large with respect to $C$. So overall
% \begin{gather*}
% \Pr_{\rho \sim \mathcal{M}_{K}}[\rho \in \mathcal{B}] \leq \exp(4d \log(2K) + \frac{d}{2}\log(4wK^2) - \frac{d}{2} \log N) \\
% \leq \exp(d(\log (16K^4) + \log (4wK) + \log N^{1/2}))
% = \exp\Bigl(d (\log wK^5 -\log N^{1/2} + O(1))\Bigr)
% \end{gather*}

% \end{proof}

% We now need the following simpler lemma:
% \begin{lemma}
% Let $\epsilon > 0$ be a sufficiently small constant. Let $D$ be a matching disjunction with $s$ terms on the function variable $f: [N] \to [M]$. Say a partial matching $\rho$ is sampled as follows: sample $A \subseteq [N]$ by including each element independently with probability $1-\epsilon$; now choose a uniform 1-1 function $A \to [M]$.
% \[
% \Pr_{\rho}[D\upharpoonright \rho \text{ has width } \geq w] < \exp(\log s - \Omega(w))
% \]
% \end{lemma}
% \begin{proof}
% Let $D = \lambda_1,\ldots,\lambda_s$. We will ignore each term whose width is already $<w$; its width cannot increase under the restriction. Let $\lambda \in D$ be a term of width $\geq w$. We bound the probability that $\rho$ does not kill the term $\lambda$ as:
% \begin{gather*}
% \Pr_{\rho}[\lambda \text{ survives }\rho] = \sum\limits_{U \subseteq \dom(\lambda)} \Pr[\dom(\lambda) \cap \dom(\rho) = U] \frac{1}{|U|! \binom{M}{|U|}} \\
% \leq \sum_{t=1}^w\Pr[|\dom(\lambda) \cap \dom(\rho)| = t] \frac{1}{t! \binom{M}{t}}\\
% \leq \frac{1}{(w/2)! \binom{M}{w/2}} +\Pr[|\dom(\lambda) \cap \dom(\rho)| < \frac{w}{2}] \\
% \leq \frac{1}{(w/2)! \binom{M}{w/2}} + \Pr_{\xi_1,\ldots,\xi_w \sim \mathrm{Bern}(1-\epsilon)} [\sum_i \xi_i \leq \frac{w}{2}] \leq \exp(- \Omega(w))
% \end{gather*}
% Thus by a union bound 
% \[
% \Pr_{\rho}[D\upharpoonright \rho \text{ has width } \geq w] \leq \exp(\log s - \Omega(w))
% \]
% \end{proof}
% Taking another union bound we get:
% \begin{corollary}
% There are absolute constants $\epsilon > 0$, $C \in \mathbb{N}$ so that the following holds. Say $D_1,\ldots,D_s$ are a collection of matching disjunctions, each of size at most $s$, over the function variable $f: [N] \to [M]$. Then there exists a partial assignment $\rho$ of size $N-K$, $K\geq \epsilon N$ so that for all $D_i$, $D_i \upharpoonright \rho$ has width at most $C\log S$.
% \end{corollary}

% \begin{theorem}
% Let $D_1,\ldots,D_s$ be a collection of matching disjunctions of size $\leq s$ over $f: [N] \to [M]$ and let $L = M-N$. Provided $(Kd\log s)^{10} < N$ \todo{check this} and $K \geq L$, there exists a partial matching $\rho$ with $K$ unset pigeons, so that $D_i \upharpoonright \rho$ has a pigeonhole decision tree of depth $\leq d$ for all $i \leq s$.
% \end{theorem}



\section{Linear Ordering Principle}\label{sec:lop}
In this section we investigate the complexity of the Linear Ordering Principle. We restate the definition in a slightly different terminology:
\begin{definition*}[Linear Ordering Principle, Restatement of \ref{def:lop}]
The input to LOP is a binary relation $\prec$ on $[N] \times [N]$ specified by a circuit. The following solutions are sought:
\begin{enumerate}
    \item A witness that $\prec$ does not define a total order on $[N]$ (trivial solution)
    \item An element $a_0 \in [N]$ so that $a_0 \prec a$ for all $a \neq a_0 \in [N]$ (nontrivial solution)
\end{enumerate}
We say an instance $\prec$ is nontrivial if it has no trivial solutions. The trivial solutions are enumerated formally as follows:
\begin{enumerate}[leftmargin=8.5\parindent]
    \item[\textbf{Antireflexivity Violation:}] $x$ such that $x \prec x$
    \item[\textbf{Transitivity Violation:}] $x,y,z$ such that $x \prec y \prec z$ and $x \not\prec z$
    \item[\textbf{Totality Violation:}] $x\neq y$ such that $x \not\prec y$ and $y \not\prec x$
\end{enumerate}
\end{definition*}
Note that in the introduction we defined $\prec$ to be a relation on $\B^n \times \B^n$. We state it in this more general form here, since the space we define a linear ordering on in our main reduction will not have size being a power of two. It is straightforward to pad an instance of LOP which is defined on some subset of $\B^n$ to the whole space without affecting the value of the solution.

Recall from Section ~\ref{sec:totalsearch}
that LOP has ``essentially unique solutions:'' the trivial solutions are checkable in polynomial time, and any instance without trivial solutions has a unique nontrivial solution. Our primary interest in this problem stems from the following reduction (stated in the Introduction):
\medskip

%\begin{theorem*}\ref{thm:avoid-to-lop-intro}
\noindent {\bf Theorem~\ref{thm:avoid-to-lop-intro}}
{\it Weak Avoid is polynomial-time many-one reducible to LOP.}
%\end{theorem*}

\medskip
We will describe the reduction in two parts, using an intermediate search problem called Forest Termination:
\begin{definition}[Forest Termination]\label{def:forest}
The input consists of:
\begin{enumerate}
    \item A function $\mathrm{Pred}: [M] \times [N] \times [N] \to [M]$
    \item A function $\mathrm{Col}:[M] \times [N] \to \B$
\end{enumerate}
specified by boolean circuits. We think of the input as representing a layered rooted forest $F$ of depth $N$ with nodes partitioned into sets $L_1, \ldots, L_{N}$, each of size $M$, with nodes in $L_1$ being roots and all other nodes in $\{L_j\}_{j> 1}$ having a unique parent in $L_{j-1}$. We say that the function $\mathrm{Pred}$ ``validly represents a forest'' if the following holds: for each $u \in [M]$, $k \leq j \leq i \in [N]$, we have that $\mathrm{Pred}(\mathrm{Pred}(u,i,j),j,k) = \mathrm{Pred}(u,i,k)$. If this holds then we may think of $\mathrm{Pred}$ as representing some forest $F$ in the following strong sense: for any node $u \in L_i$ with $i >1$, the nodes $\mathrm{Pred}(u, i, 1),\ldots, \mathrm{Pred}(u,i,i)$ form the unique path from $L_1$ to $u$ in $F$. Finally for each each $u \in F$ with think of $\mathrm{Col}(u, i)$ as coloring the node $u \in L_i$ red or blue.

We say that $u \in F \cup \{\bot\}$ is a termination point of $F$ if the following holds:
\begin{enumerate}
    \item $u = \{\bot\}$ and $L_1$ contains no blue node.
    \item $u$ is a node in $F$, the path from $L_1$ to $u$ in $F$ uses only blue nodes, and $u$ does not have a blue child in $F$.
\end{enumerate}

Now the search problem is: given $\mathrm{Pred}, \mathrm{Col}$, find a witness that $\mathrm{Pred}$ does not validly represent a forest, or else find a termination point in the forest it represents.
\end{definition}
The totality of this problem follows by starting at at a blue node in $L_1$ (if one exists) and finding a maximal blue path through its descendants in $F$. We start by giving a reduction from Weak Avoid to Forest Termination. This is the aspect of the proof which is very similar to arguments by Li and by Paris Wilkie and Woods \cite{Li,Paris-Wilkie-Woods}.

\begin{lemma}\label{lem:avoid-forest}
Weak Avoid with $2n$ stretch is polynomial-time reducible to Forest Termination.
\end{lemma}
\begin{proof}
Let $f: \B^n \to \B^{2n}$ be given; we denote by $f_0,f_1: \B^n \to \B^n$ the functions obtained by restricting to the first and last $n$ bits of output respectively. Let $\mathcal{S}_n$ denote the set of all binary strings of length at most $n$, including the empty string $\epsilon$. Naturally $\mathcal{S}_n$ is associated with the nodes of a full binary tree of depth $n$; however since the Forest Termination instance we construct involves a forest whose structure and depth differ from that of $\mathcal{S}_n$, we will use the terminology of binary strings to refer to elements of $\mathcal{S}_n$ in order to avoid confusion. We say $s \sqsubseteq s'$ if $s$ is a prefix of $s'$, and $s\sqsubset s'$ if it is a proper prefix. We use $s \cdot s'$ for concatenation. If $s \sqsubseteq s'$, with $s' = s \cdot p$, we use $s' - s$ to denote $p$. We define the ``preorder'' $<$ on $\mathcal{S}_n$ recursively with respect to prefixes: for a prefix $p$ and two distinct extensions $p \cdot s, p \cdot s'$ with $s < s'$ (with the relative order of $s,s'$ defined recursively), we put $p \cdot s < p \cdot s' < p$. In the view of $\mathcal{S}_n$ as a depth $n$ perfect binary tree, this corresponds to the recursive subtree ordering ``left subtree $< $ right subtree $< $ root.'' Note that this is a total ordering with $\epsilon$ being the greatest element and $0^n$ the least. Finally, for an element $s \in \mathcal{S}_n$, we say that $s'$ is a ``left outlet'' of $s$ if either $s'=s$ or $s'=p0$ for some $p \sqsubseteq s$. We use $\mathrm{LO}(s)$ to denote its set of left outlets; note that $|\mathrm{LO}(s)| \leq n$. Observe that for any $s' \leq s$, there exists a unique $p \in \mathrm{LO}(s)$ such that $p \sqsubseteq s'$.

For some $s \in \mathcal{S}_n$ and some value $x \in \B^n$, define $f_s(x)$ recursively as follows: if $s = \epsilon$ then $f_s(v)=v$. If $s = bs'$ for some bit $b$ and substring $s'$, $f_{s} = f_{s'}(f_b(x))$.

We define a ``transcript'' as an element $\pi \in (\B^n)^n$. We say the transcript $\pi$ is valid for some $s \in \mathcal{S}_n$ if $\pi_i = 0^n$ for all $i > |\mathrm{LO}(s)|$. In this way the valid transcripts for $s$ are in one-to-one correspondence with functions $\mathrm{LO}(s) \to \B^n$. Thus for a valid transcript $\pi$ for $s$ and some $p \in \mathrm{LO}(s)$, we use $\pi(p) \in \B^n$ to denote the value the transcript associates to $p$. 

For a transcript $\pi$, and strings $s' < s \in \mathcal{S}_n$, we define the transcript $\pi' = \mathrm{Pred}(\pi, s, s')$ as follows. We first check if $\pi$ is valid for $s$; if not we set $\pi'$ to be some fixed canonical choice of invalid transcript for $s'$. Otherwise, for each $p \in \mathrm{LO}(s')$, we set $\pi'(p) = f_{p-q}(\pi(q))$ where $q \in \mathrm{LO}(s)$ is the unique left outlet of $s$ which is an ancestor of $p$. 

We now construct an instance of Forest Termination with $[N] = [2^{n+1}-1] = \mathcal{S}_n$ and $M = [2^{n^2}] = (\B^n)^n$. $[N]$ is associated naturally to $\mathcal{S}_n$, so that each layer corresponds to an element of $\mathcal{S}_n$ and the layers are ordered according to $<$, i.e. $L_1$ is associated to $0^n \in \mathcal{S}_n$ and $L_N$ to $\epsilon \in \mathcal{S}_n$. $[M]$ is associated to the set of possible transcripts, so that $L_s$ will be the set of all possible transcripts for the string $s$ (some valid and others invalid). We use $\mathrm{Pred}$ defined above to give the $\mathrm{Pred}$ function in the instance of Forest Termination. It should be noted that the depth of the forest defined here is $N = 2^{n+1}-1$, which is much larger then the depth of tree naturally associated to $\mathcal{S}_n$ (whose depth is $n$); this is why we have used the notation of binary strings rather then nodes in a tree to refer to elements of $\mathcal{S}_n$. It remains to define the node coloring function $\mathrm{Col}$.

If $(\pi,s) \in [M] \times [N]$ are such that $\pi$ is not a valid transcript for $s$, then $\mathrm{Col}(\pi,s)$ is red; if $\pi$ is a valid transcript for $s$ and $|s| < n$ then it is colored blue ($|s|$ denotes its length as a binary string). Otherwise if $\pi$ is a valid transcript for $s$ and $|s|=n$, we check that $\pi(s)= \lnot f_s(s)$, where $\lnot$ denotes the string obtained by flipping all bits. If so we color it blue, otherwise we color it red. Note that since $|s|=n$, $f(s)$ is well defined, and thus so is $f_s(s)$.

It follows by construction that $\mathrm{Pred}$ validly defines a forest in the sense of Definition~\ref{def:forest}. It remains to show that given any termination point we may determine a solution to the original instance of range avoidance. We first claim that there exists a blue node in the first layer $L_{0^n}$, so that the solution to Forest Termination cannot be $\bot$. Observe that $\mathrm{LO}(0^n) = \{0^n\}$, thus the transcript $\pi$ which sets $\pi(0^n) = \lnot f_{0^n}(0^n)$ will be blue. Now, say that $\pi \in L_s$ is a termination node. We claim that $L_s$ cannot be the last layer, i.e. $s \neq \epsilon$. Say towards a contradiction that $\pi$ is a termination node in layer $L_{\epsilon}$. So $\pi$ is a valid transcript for $\epsilon$, which stores a single value $x = \pi(\epsilon)$. By the assumption $\pi$ is a termination point all of its predecessors are blue. In particular if we view $x$ as a length-$n$ element of $\mathcal{S}_n$, then we have that $\mathrm{Pred}(\pi,\epsilon, x)$ is blue, which by construction means that $f_{x}(x) = \lnot f_{x}(x)$, which is a direct contradiction.

Now say that $\pi$ is a termination node in layer $L_s$ where $s \neq \epsilon$; let $s'$ be the successor of $s$ in the ordering $<$. First say that $|s'|=n$; in this case $(\pi, s)$ cannot be a termination point, since we may construct a blue transcript $\pi'$ for $s'$ such that $\mathrm{Pred}(\pi',s',s)=\pi$ by simply appending the value $\lnot f_{u'}(u')$ as $\pi'(u')$ to the original transcript $\pi$ for $s$. Otherwise we have $s=s'1$. Let $p=s'0$; so $p$ is a left-outlet of $s$, and therefore $\pi$ stores some values $x_0 = \pi(p), x_1 = \pi(s)$. We claim that $x_0x_1$ must be a solution to the Avoid instance $f$. Say this is not true and there is some $x \in \B^n$ such that $f(x)=x_0x_1$; i.e. $f_0(x)=x_0$ and $f_1(x)=x_1$. Then we can generate a blue transcript $\pi'$ for $s'$ as follows. Observe that $\mathrm{LO}(s') = \mathrm{LO}(s) \setminus \{p\} \cup \{s'\}$. Thus if we remove $p$ from the domain of $\pi$ and set $\pi'(s') = x$ to obtain $\pi'$, then $\pi'$ will be a blue transcript for $s'$ and we will have $\mathrm{Pred}(\pi',s',s)=\pi$, which means $(\pi, s)$ cannot be a termination point. To see this, note that in the computation of $\mathrm{Pred}(\pi',s',s)$, the only modification made to $\pi'$ to generate $\pi$ will be to compute $z_0 =f_0(\pi'(s')), z_1 = f_1(\pi(s'))$, and add $z_0$ as the value $\pi$ associates to $p$ and $z_1$ as the value it associates to $s$. By construction $z_0 = x_0$ and $z_1=x_1$ so we end up with the original transcript $\pi$.

\end{proof}

Next we reduce Forest Termination to LOP, which is simpler:
\begin{lemma}\label{lem:forest-lop}
Forest Termination is polynomial-time many-one reducible to LOP.
\end{lemma}
We will prove shortly that any search problem which is $\mathsf{P^{NP}}$-reducible to LOP is also polynomial-time many-one reducible to LOP (Theorem~\ref{thm:l2p-closure-intro}); therefore it suffices to give a $\mathsf{P^{NP}}$ reduction.
\begin{proof}
Let $\mathrm{Pred},\mathrm{Col}$ an instance of Forest Termination. We start by using an $\mathsf{NP}$ oracle to search for a witness that $\mathrm{Pred}$ fails to validly represent a forest; if found we output this as our solution, otherwise we know that the instance indeed represents a forest $F$ on the nodes $L_1 \sqcup \cdots \sqcup L_N$. For conceptual simplicity, we modify $F$ into a tree $T$ by adding a virtual blue root node $\bot$ whose children consist of all nodes in $L_1$. We use $<$ for the lexicographical ordering on each $L_i$. We now define a new ordering $\prec$ on the nodes of $T$. Its definition is given recursively with respect to subtrees of $T$ as follows: say that $u$ is a node with blue-rooted subtrees $B_1 < \ldots < B_K$ and red-rooted subtrees $R_1<\ldots<R_{K'}$ ordered lexicographically by their respective roots. Define $\prec$ recursively inside each $B_i, R_j$; now between them use the ordering $B_1 \prec \ldots \prec B_K \prec \{u\} \prec R_1 \prec \ldots \prec R_{K'}$. By induction on the depth of $T$ we see that $\prec$ is a total ordering. The relation $\prec$ will be our instance of LOP output by the reduction; it remains to show that $\prec$ is efficiently computable and that its minimal element yeilds a Forest Termination solution.

We first show how to compute $\prec$. Given $u_{0} \neq u_{1} \in T$:
\begin{enumerate}
    \item If $u_0$ is an ancestor of $u_1$ then let $v$ be the child of $u_0$ from which $u_0$ can be reached. If $v$ is blue then set $u_1 \prec u_0$, else if it is red set $u_0 \prec u_1$.
    \item Otherwise let $v$ be the least common ancestor of $u_0$ and $u_1$, and let $w_0,w_1$ be the children of $w$ from which $u_0$ and $u_1$ can be reached respectively. We define the relative ordering of $w_0,w_1$ in $\prec$, which is then inherited by the pair $u_0,u_1$: if $w_b$ is blue and $w_{\lnot b}$ is red for some $b \in \B$, then set $w_b \prec w_{\lnot b}$. Otherwise use the lexicographical order $<$ to order the pair $w_0,w_1$.
\end{enumerate}
The only nontrivial step here is to compute a least common ancestor of $u_1,u_2$, which can be accomplished with binary search. Say $u_1,u_2$ lie in layers $L_i, L_{i'}$ respectively, with $i' > i$. The using $\mathrm{Pred}$ we find the ancestor of $u_2$ in $L_{i}$; call it $\hat{u}_2$. If $\hat{u}_2 = u_1$ then $u_1$ is the ancestor of $u_2$. Otherwise let $j = \lceil i/2 \rceil$, and compute the ancestors $v_1,v_2$ of $u_1, \hat{u}_2$ in layer $L_j$. If they are distinct then we repeat the same procedure starting from $v_1,v_2$. Otherwise if $v_1 = v_2$ we prune all the layers $L_1,\ldots,L_{j-1}$ from consideration and repeat the procedure from $u_1,\hat{u}_2$, taking $L_j$ to be the root layer with $v_1=v_2$ as its root, and renumbering the layers accordingly. Overall the number of steps required is $O(\log N)$.

Finally it remains to show that any solution to LOP on the instance $\prec$ yields a solution to the given Tree Termination instance, i.e. if $u \in\{\bot\} \sqcup L_1 \sqcup \cdots \sqcup L_N$ is minimal in the ordering $\prec$ then it must correspond to a termination point. Say that $u$ is not a termination point, then either:
\begin{enumerate}
    \item There is an edge $(v,w)$ in the path from the root to $u$ so that $w$ is colored red. In this case we must have $v \prec u$.
    \item $u$ has a blue child $v$. In this case we have $v \prec u$.
\end{enumerate}
In each case we see that $u$ cannot be the minimal element. An example of a Forest Termination instance and its corresponding order $\prec$ is given in Figure~\ref{fig:lop-tree}.
\end{proof}
\begin{figure}
    \centering
    \includegraphics[width=0.4\textwidth]{Lop_tree.pdf}
    % \captionsetup{width=1\textwidth}
    \caption{Example instance of Forest Termination with $M = 4$ and $N=3$, with virtual root added at the top. Blue nodes are circles and red nodes are diamonds. We have sorted each layer left to right so that blue nodes come first, and within each color class the nodes are ordered left to right lexicographically. The number on each node represents its position in the ordering $\prec$. The nodes numbered $1$ and $4$ are the termination points of this instance.}
    \label{fig:lop-tree}
\end{figure}

We now observe some special properties held by the search problem LOP, the first of which is its strong closure properties under a wide class of reductions:
\begin{theorem*}[Theorem~\ref{thm:l2p-closure-intro}]
Let $R$ be any search problem which has a polynomial-time, $\mathsf{NP}$-oracle Turing reduction to LOP. Then $R$ is polynomial-time many-one reducible to LOP.
\end{theorem*}
\begin{proof}
Let $A$ be a polynomial time $\mathsf{P^{NP}}$ Turing reduction from some search problem $R$ to LOP. Let $x$ be an instance of the search problem of length $n$. We first claim that we can modify $A$ so that:
\begin{enumerate}
    \item $A(x)$ only calls the LOP oracle on instances $\prec$ which are nontrivial, i.e. they define a true total order.
    \item $A(x)$ does not use its $\mathsf{NP}$ oracle.
\end{enumerate}
To achieve the first condition, before each call to the LOP oracle on an instance $\prec$, we first use the $\mathsf{NP}$ oracle to check if $\prec$ defines a total order, and to find a violation if not. If a violation is found we no longer need to use the LOP oracle on this instance; otherwise we know that $\prec$ is a nontrivial instance and we use the oracle to solve it. Now to achieve the second condition, if $\phi$ is an instance of $\mathsf{SAT}$ with $m$ variables, we define the ordering $\prec_{\phi}$ on $\B^m$ by: $y \prec_{\phi} z$ if $\phi(y) = 1 \land \phi(z)=0$ or vice versa, otherwise $\prec_{\phi}$ orders them lexicographically. Clearly $\prec_{\phi}$ is a nontrivial LOP instance, and the minimal element of $\prec_{\phi}$ will tell us if $\phi$ is satisfiable.

At this point we have a polynomial time reduction $A$ which makes adaptive oracle calls to LOP on nontrivial instances, each of which has a unique solution. We next modify $A(x)$ to make exactly $m$ calls on all computation paths, and have each call of the form $\prec: \B^m \times \B^m \to \B$ for some fixed value $m = \mathrm{poly}(n)$. We can accomplish this easily by padding the algorithm with dummy queries and padding any instance of LOP to a larger bit-length. We now define a new instance $\prec_{x}: \B^{m^2} \times \B^{m^2} \to \B$ as follows. An element of $\B^{m^2}$ is given by a sequence $(u_1,\ldots,u_m)$ with $u_i \in \B^m$. To compare two elements $\bar{u} = (u_1,\ldots,u_m), \bar{v}= (v_1,\ldots v_m)$, we find the first index $i \in [m]$ so that $u_i \neq v_i$. We simulate $A(x)$ through the first $i-1$ queries, plugging in $u_j = v_j$ as the answer to the $j^{th}$ LOP oracle query. Once we get to the $i^{th}$ query, we look at the instance $\prec_i$ of LOP defining the next query, and compare $u_i, v_i$ using $\prec_i$; we then order $\bar{u},\bar{v}$ accordingly, e.g. $\bar{u} \prec_x \bar{v}$ if $u_i \prec_i v_i$. It follows by construction that $\prec_x$ is a total order, and its least element  corresponds to the unique sequence of correct oracle responses in the computation $A(x)$. Therefore by the correctness of $A$ we may read off in polynomial time from this sequence the solution to the search problem $R$ on input $x$.
\end{proof}

At this point we can prove Theorem~\ref{thm:avoid-to-lop-intro} that Weak Avoid is polynomial time many-one reducible to LOP:
\begin{proof}[Proof of Theorem~\ref{thm:avoid-to-lop-intro}]
By \cite{Korten} there is a $\mathsf{P^{NP}}$ reduction from the general Weak Avoid problem to the special case when the stretch is $2n$. Composing this with the reductions in Lemmas \ref{lem:avoid-forest} and \ref{lem:forest-lop} we obtain a $\mathsf{P^{NP}}$ reduction from Weak Avoid to LOP. Applying Theorem~\ref{thm:l2p-closure-intro} this in turn yields a polynomial-time many-one reduction.
\end{proof}

Recall that LOP enjoys the special property of having essentially unique solutions. Together with its closure under $\mathsf{P^{NP}}$ reductions, this endows LOP with the unusual property of being equivalent to a decision problem. We define now a decisional complexity class $\mathsf{L_2^P}$, possessing a few equivalent definitions, whose complete problem is equivalent to the search problem LOP: 
\begin{definition*}[Reminder of Definition~\ref{def:l2p}]
A language $L$ is in the complexity class $\mathsf{L_2^P}$ if there is a polynomial time relation $R: (\B^*)^3 \to \B$ and a polynomial $p$, so that for all $x$, $R(x, \cdot, \cdot)$ defines a total order on $\B^{p(n)}$ whose minimal element $a$ has $a_1 = L(x)$. 
\end{definition*}
% \todo{note that for black-box definable TFNP search problems, the only search problems equivalent to a language are in $\mathsf{P}$ so this is a situation that can only occur at the second level of the hiearchy.}
Recalling Lemma~\ref{lem:lop-in-s2p}, we have the following upper bound for this class:
\begin{observation}
$\mathsf{L_2^P} \subseteq \mathsf{S_2^P}$.
\end{observation}
At first sight $\mathsf{L_2^P}$ seems to crucially involve a promise, namely that $\prec_x$ defines a total order for all $x$, just like the promise defining the class $\mathsf{S_2^P}$. However it turns out that this promise can be eliminated and $\mathrm{LOP}$ can be given a purely syntactic characterization. This is summarized in Theorem~\ref{thm:l2p-complete-intro} which we restate here:

\medskip

%\begin{theorem*}[Theorem~\ref{thm:l2p-complete-intro}]
\noindent {\bf Theorem~\ref{thm:l2p-complete-intro}}
{\it The following are equivalent for a language $L$:
\begin{enumerate}
    \item $L \in \mathsf{L_2^P}$
    \item $L$ is $\mathsf{P^{NP}}
$-Turing reducible to LOP.
    \item $L$ is polynomial time many-one reducible to LOP.
\end{enumerate}
Conversely, the search problem LOP is polynomial time truth table reducible to a language in $\mathsf{L_2^P}$.}
%\end{theorem*}

\begin{proof}
If $L \in \mathsf{L_2^P}$ then by definition it is many-one reducible to LOP which gives $(1)\rightarrow(2)$. Theorem~\ref{thm:l2p-closure-intro} gives $(2)\rightarrow(3)$; inspecting the proof, for each $x$ we may define $\prec_x$ so that the value $L(x)$ is the first bit of the LOP solution to $\prec_x$, which actually gives $(2)\rightarrow(1)$.
\end{proof}

% \begin{lemma}
% In the decision tree model, $\mathsf{L_2^P} \not\subseteq \mathsf{P^{NP}} $; consequently  $\mathsf{S_2^P} \not\subseteq \mathsf{P^{NP}} $ and $\mathsf{ZPP^{NP}} \not\subseteq \mathsf{P^{NP}}$ hold in the decision tree model as well.
% \end{lemma}
% \todo{Can we use the depth 2.5 lower bound to show that Strong Avoid is not equivalent to any language in this sense? We know that if a language $L$ reduces to any $\mathsf{TF\Sigma_2}$ search problem then it lies in $\mathsf{\Sigma_2 \cap \Pi_2}$ (same holds under $\mathsf{P^{NP}}$ reductions). So then if Strong Avoid is equivalent to a language, then it must lie in $\mathsf{FP^{\Sigma_2 \cap \Pi_2}}$. But then we should have $\mathsf{FP^{\Sigma_2 \cap \Pi_2}} \subseteq \mathsf{FP}_{\Vert}^{\Sigma_2}$ and our lower bound says Strong Avoid is not contained in the latter class.}
Finally we prove Theorem~\ref{thm:l2p-structural-intro} from the introduction, which tells us that $\mathsf{L_2^P}$ satisfies two of the three most interesting properties of the larger class $\mathsf{S_2^P}$:

%\begin{theorem*}[Theorem~\ref{thm:l2p-structural-intro}]

\medskip

\noindent{\bf Theorem~\ref{thm:l2p-structural-intro}}
{\it \phantom\\
\begin{enumerate}
    \item $\mathsf{P^{NP}} \subseteq \mathsf{L_2^P}$ and $\mathsf{BPP} \subseteq \mathsf{MA} \subseteq \mathsf{L_2^P}$
    \item $\mathsf{L_2^E}$\footnote{$\mathsf{L_2^E}$ is the exponential-time analogue of $\mathsf{L_2^P}$, where we replace ``polynomial time'' with $2^{O(n)}$ time'' in its definition} contains a language of circuit complexity $2^n/n$.
\end{enumerate}
}
%\end{theorem*}

%\medskip

\begin{proof}
The inclusion $\mathsf{P^{NP}} \subseteq \mathsf{L_2^P}$ follows directly from the closure of $\mathsf{L_2^P}$ under $\mathsf{P^{NP}}$ reductions. For the second inclusion, we know that every language in $\mathsf{MA}$ is $\mathsf{P^{NP}}$-reducible to the construction of a $O(\log n)$-seed length PRG for $O(n)$-size circuits \cite{NW}. The explicit construction of such PRGs is in turn reducible to Weak Avoid \cite{Korten}, and Weak Avoid is reducible to $\mathsf{LOP}$ by Theorem~\ref{thm:avoid-to-lop-intro}. Applying Theorem~\ref{thm:l2p-closure-intro} the result follows.

For the second part, the proof is identical to the case of $\mathsf{S_2^E}$ shown in \cite{Li} given our reduction from Weak Avoid to LOP. Given $x \in \B^n$ we may produce in $2^{O(n)}$ time an instance of Weak Avoid $C_n:\B^{\ell} \to \B^{2^n}$ so that any solution is a truth table of a function $f: \B^n \to \B$ of circuit complexity $2^n/n$. By Theorem~\ref{thm:avoid-to-lop-intro} we may then construct in $2^{O(n)}$ time a nontrivial instance of LOP $\prec$ so that from the unique minimal element of $\prec$ we may read off a uniquely defined Weak Avoid solution $f_n$ for $C_n$. We then accept/reject $x$ based on $f_n(x)$. Using the closure properties of $\mathsf{L_2^E}$ the language $\{f_n\}_{n \in \mathbb{N}}$ thus defined lies in $\mathsf{L_2^E}$.
\end{proof}

The only interesting property of $\mathsf{S_2^P}$ which we cannot prove is inherited by its subclass $\mathsf{L_2^P}$ is the Karp-Lipton theorem:
\begin{theorem*}[\cite{cai}, credited to Sengupta]\label{thm:s2p-kl}
If $\mathsf{NP} \in \mathsf{P/poly}$ then $\mathsf{PH} = \mathsf{S_2^P}$.    
\end{theorem*}
The proof of this result does not seem to generalize to $\mathsf{L_2^P}$.
\section{Proof of the Pigeonhole Switching Lemma}
\label{sec:php-sl}
As in the standard proof of the switching Lemma, we %need the following definition which associates a particular pigeonhole decision tree to any matching disjunction.
we need to define a procedure for representing a matching disjunction by a {\it canonical} complete pigeonhole decision tree. We start with some preliminary definitions.

\begin{definition}[Complete Pigeonhole Trees on a subset $A$] 
For a partial matching $\pi$ and a hole restriction $\tau$, we say that $\tau$ covers $\pi$ if $\nodes(\tau) \supseteq \nodes(\pi)$. If $\tau = \langle \sigma, E\rangle$, we say that $\tau$ is a minimal cover of $\pi$ if it fails to be a cover after removing any edge from $\sigma$ or any element from $E$.

For every tuple $(A,N,M)$ where $A \subseteq [N] \sqcup [M]$, 
we define a canonical complete pigeonhole decision tree, denoted denoted by $\mathcal{C}_A(N,M)$, as follows. This tree is unlabelled, and is defined with respect to some fixed ordering of the elements of $A$. So lets assume the ordering $a_1,\ldots,a_r$ of $A$. 
We build $\mathcal{C}_A(N,M)$ in $r$ steps:
\begin{enumerate}
    \item In step 1, if $a_1 \in [N]$, we query pigeon $a_1$ at the root, with all possible outedges: $(a_i,y)$, $y \in [M]$. Otherwise if $a_1 \in [M]$, we query hole $a_1$ with all possible outedges: $(x,a_i), x \in [N] \sqcup \{\bot\}$. 
    \item In step $i>1$, for each leaf node $l$ in $T$, let $\tau_l = \langle \sigma_l,E_l \rangle$ denote the hole restriction associated with the path to $l$.  If $a_i$ is a pigeon or hole in $nodes(\tau_l)$ (that is, if $a_i$ has already been "determined" by the hole restriction thus far), then we do nothing.  Otherwise,  we query $a_i$ at the leaf $l$, with all possible outedges that are consistent with $\tau_l$. 
    \end{enumerate}
\end{definition}
Observe that the leaves of $\mathcal{C}_A(N,M)$ are in one-to-one correspondence with the set of all hole restrictions that minimally cover $A$, and thus ${\cal C}_A(N,M)$ is a complete pigeonhole decision tree.

\begin{definition}[Canonical Pigeonhole Trees]
Let $\phi = \pi_1 \lor \ldots \lor \pi_{s}$ be a matching disjunction over $[N] \sqcup [M]$  where we have fixed this ordering on the terms. 
The canonical pigeonhole tree for $\phi$, 
$\mathcal{T}_{\phi}(N,M)$, is defined recursively below. When $N,M$ is clear from context, we will abbreviate $\mathcal{T}_{\phi}(N,M)$ by $\mathcal{T}_{\phi}$.

\begin{enumerate}
    \item If $s=0$, then $\phi$ is a constant function. If $\phi=1$ then $\mathcal{T}_{\phi}(N,M)$ consists of a single node labelled by 1, and similarly if $\phi$ is 0, then the decision tree is a single node labelled by 0.
    \item Otherwise, let $\pi_1$ be the first term in $\phi$ according to the ordering of terms in $\phi$. First, we create the complete pigeonhole tree ${\cal C}_{ nodes(\pi_1)}(N,M)$. For each leaf $l$ in the tree constructed so far, let $\tau_l = \langle \sigma_l,E_l \rangle$ be the hole restriction associated with the path to $l$. Then for each leaf $l$: 
    \begin{enumerate}
        \item Let $P \subseteq [N]$ be the subset of pigeons that have not been matched  by $\pi_l$, and let $H \subseteq [M]$ be the subset of holes not matched by $\pi_l$ and not in $E_l$. $P$ and $H$ define the pigeons and holes that are still undetermined after applying $\tau_l$.  If the size of $P$ is $N'$ and the size of $H$ is $M'$, then we can rename the vertices so that we identify $P$ with $[N']$ and $H$ with $[M']$.
        \item Recursively construct ${\cal T}_{\phi\upharpoonright \tau_l}(N',M')$. and attach a copy of this subtree to the leaf $l$.
    \end{enumerate}
    \end{enumerate}
\end{definition}

%\begin{definition}[Canonical Decision Tree]
%For a partial matching $\pi$ and a hole restriction $\tau$, we say that $\tau$ covers $\pi$ if $\nodes(\tau) \supseteq \nodes(\pi)$. If $\tau = \langle \sigma, E\rangle$, we say that $\tau$ is a minimal cover of $\pi$ if it fails to be a cover after removing any edge from $\sigma$ or any element from $E$.

%For a set $A \subseteq [N] \sqcup [M]$, we denote by $\mathcal{C}_A$, the complete unlabeled decision tree over $A$, the pigeonhole decision tree (without leaf labels) which queries each element of $A$ in a fixed order. Observe that the leaves of $\mathcal{C}_A$ are in one to one correspondence with minimal covers of $A$. \todo{Change so that we delete paths not consistent with a 1-1 function.}

%For a matching disjunction $\phi = \pi_1 \lor \ldots \lor \pi_{s}$ with some fixed ordering on terms, we associate a pigeonhole decision tree $\mathcal{T}_{\phi}$, referred to as the ``canonical decision tree for $\phi$,'' as follows: 
%First, if $\phi$ is 1 then the decision tree consists of a single node labelled by 1, and similarly if $\phi$ is 0, then the decision tree is a single node labelled by 0. Otherwise, we construct $\mathcal{C}_{\nodes(\pi_1)}$; recall that each leaf is associated to a minimal cover $\tau$ of $\pi_1$. At the unique such leaf associated to $\tau = \pi_1$ we label the leaf by the value 1. At any leaf $\tau$ such that $\tau \nparallel \pi_1$ we label it with the value $0$. Otherwise we construct recursively $\mathcal{T}_{\phi \upharpoonright \tau}$ and attach a copy of this subtree to the current leaf associated with $\tau$.
%\todo{I would replace everything above starting with "recall that each leaf" with : Now for each leaf of  ${\mathcal C}_{\nodes(\pi_1)}$, we apply the same procedure recursively on the restricted matching disjuncdtion $\phi|_{\tau}$, where $\tau$ is the partial matching associated with this leaf.} 
%\end{definition}
% \todo{modified def:}
% \begin{definition}[Canonical Decision Tree]
% Let $\pi = (x_1,y_1), \ldots, (x_k,y_k)$ be a partial matching with some fixed ordering on its constituent edges. We define $\mathcal{T}_{\pi}$, the canonical decision tree for $\pi$, to be an unlabeled pigeonhole decision tree obtained as follows:
% \begin{enumerate}
%     \item The root makes the query $x_1$. At the child labeled by the response $y_1$ we place recursively a copy of $\mathcal{T}_{\pi \setminus (x_1, y_1)}$
% \end{enumerate}
% \end{definition}

\medskip
It follows from the definition that $\mathcal{T}_{\phi} \equiv \phi$. 
%\begin{definition}[Distribution of Partial Restrictions]
Next, we restate the Pigeonhole Switching Lemma.

\medskip

%\begin{lemma*}[Lemma \ref{lem:switching-main}]
%Let $\mathcal{M}_K$ be the uniform distribution over the set  of all partial matchings $[N] \to [M]$ with exactly $N - K$ edges. 
%\begin{lemma}\label{lem:switching-main}
\noindent{\bf Lemma~\ref{lem:switching-main}} 
{\it Let $M,N,d \in \mathbb{N}$. 
%$L = M - N$. 
Let $\phi$ be a width-$w$ matching disjunction over $[M]^{[N]}$. 
%Say that $\rho$ is sampled uniformly from the set of partial matchings $[N] \to [M]$ using $N-K$ edges. 
If $M-N \leq K \leq \frac{N}{4}$ and $N,K,d,w$ sufficiently large, then:
\[
\Pr_{\rho \sim \mathcal{M}_K} [\mathcal{T}_{\phi\upharpoonright \rho} \text{ has depth } \geq d] \leq \exp\Bigl(d (\log wK^5 -\log N^{1/2} + O(1))\Bigr)
\]
}
%\end{lemma*}
% We first need the following definition:
% \begin{definition}
% Let $\mathrm{stars}(w,t)$ denote the set of all sequences $\beta_1\cdots\beta_{\ell}$ such that each $\beta$ is a string in $\{\star,\bot\}^w \setminus \{\bot\}^w$ and the total number of $\star$ symbols across all $\beta_i$ is at most $t$.
% \end{definition}
% \begin{lemma}\todo{ref beame}
% $|\mathrm{stars}(w,t)| \leq (w/\ln 2)^t$
% \end{lemma}
%\medskip

\begin{proof}
Let $\phi= \lambda_1 \lor \ldots \lor \lambda_s$ be a matching disjunction of width $w$. Let $\mathcal{B}$ be the set of partial matchings $\rho$ with $N - K$ edges, such that $\mathcal{T}_{\phi \upharpoonright \rho}$ has depth $\geq d$. 
If $\rho \in \mathcal{B}$, there exists a total assignment $f \supseteq \rho$ yielding a path $P$ of length $d$ from the root in $\mathcal{T}_{\phi \upharpoonright \rho}$. We define a sequence of hole restrictions $\tau_1,\dots,\tau_{\ell}$ and a sequence of partial matchings $\sigma_1,\ldots,\sigma_{\ell}$, and set $S \subseteq [\ell] \times [w]$ as follows. Define $\tau_0 = \emptyset$. Once $\tau_i$ has been defined, let $a_{i+1}$ be the index of the first term in $\phi$ so that $\lambda_{a_{i+1}} \upharpoonright \rho \cup \tau_1 \cup \cdots \cup \tau_i$ is unkilled. Then let $\tau_{i+1}$ be the minimal map $\tau_{i+1} \subseteq f$ which covers $\lambda_{a_{i+1}} \upharpoonright \rho \cup \tau_1 \cup \cdots \cup \tau_i$, and let $\sigma_i = \lambda_{a_{i+1}} \upharpoonright \rho \cup \tau_1 \cup \cdots \cup \tau_i$. Finally let $S \subseteq [\ell] \times [w]$ contain each index $(i,j) \in [\ell] \times [w]$ such that the $j^{th}$ edge of $\lambda_i$ occurs in $\sigma_i$. We claim that such a sequence can be constructed so that at the end $\ell \leq d$, and if we take $\sigma = \cup_i \sigma_i$, then $\sigma$ is a partial matching with $\nodes(\sigma) \cap \nodes(\rho) = \emptyset$ and $|\sigma|=t$ for some $\frac{d}{2} \leq t \leq dw$. This follows from the definition of $\mathcal{T}_{\phi \upharpoonright \rho}$ and and the fact that $f$ achieves a path of length $\geq d$ in this decision tree.

%Let $\mathcal{M}_R$ be the uniform distribution over the set  of all partial matchings $[N] \to [M]$ with exactly $N - R$ edges. 
Thus, $\rho \in \mathcal{M}_K$, while $\rho \cup \sigma \in \mathcal{M}_{K-t}$ for some $\frac{d}{2} \leq t \leq dw$. We will define a map 
\[
\mathrm{Dec}: \bigcup\limits_{t \geq \frac{d}{2}}\mathcal{M}_{K-t}  \times \Gamma_t \times \Delta_t \to \mathcal{M}_K
\]
for some finite sets $\{\Delta_t\}_{t \leq N}$ so that whenever $\rho \in \mathcal{B}$ and $\sigma, S$ are constructed from $\rho$ as above, there exists some ``advice'' $\delta \in \Delta_t$ so that $\mathrm{Dec}(\rho \cup \sigma, S,\delta) = \rho$. This will imply that
\[
\Pr_{\rho \sim \mathcal{M}_{K}}[\rho \in \mathcal{B}] \leq \frac{1}{|\mathcal{M}_K|} \sum\limits_{t \geq \frac{d}{2}} |\mathcal{M}_{K-t}|  |\Gamma_t| |\Delta_t|
\]
which will yield the theorem provided we can choose the sets $\Gamma_t$, $\Delta_t$ sufficiently small.

Let $\rho \in \mathcal{B}$, $\tau_i,a_i,\sigma_i, S$ be as above. Define $\gamma_{i}:=\rho \cup \tau_1 \cup \ldots \cup \tau_{i} \cup \sigma_{i+1} \cup \ldots \cup \sigma_{\ell}$, where $\gamma_0 = \rho \cup \sigma$. Let $\sigma=\cup_i\sigma_i$ and $t = |\sigma_i|$. First, we claim that if we are given $\tau_1,\ldots,\tau_{i-1}$, $\sigma_1,\ldots,\sigma_{i-1}$, and  and $\gamma_{i-1}$ then we may decode from these the index $a_i$. To do this, we simply search for the first satisfied term in $\phi \upharpoonright \gamma_{i-1}$. Next, observe that if we have knowledge of $a_i$, then using $S$ we may reconstruct $\sigma_i$. Once we have $\gamma_{i-1}, \sigma_i, \tau_i$, we may construct $\gamma_i$ directly, by removing $\sigma_i$ in $\gamma_{i-1}$ and replacing it with $\tau_i$. So it remains to show how to specify $\tau$ using the extra information $\Delta_t$, so that we may recover $\tau_i$ from $\sigma_i$. Obverse that $\nodes(\tau)$ is contained in $\nodes(\sigma) \cup ([N] \sqcup [M] \setminus \nodes(\rho \cup \sigma))$. Therefore if we construct a map $\delta: [K] \to [K + L]$, where $M-N=L$ and $|\delta| \leq 2d$ which specifies, for each $u \in \nodes(\sigma)$ the preimage/image of $u$ in $\tau$ (or the lack of preimage if $u$ is marked empty), then from $\sigma_i$ we may reconstruct $\tau_i$ by searching for the images/preimages of the nodes in $[K] \sqcup [K+L]$ corresponding to $\nodes(\sigma_i)$. Thus for every $t$ we can take $\Delta_t = \Delta$ to be the set of all hole restrictions $[K] \to [K+L]$ of size at most $d$, so that $|\Delta| \leq \binom{K}{2d}\binom{K+L}{2d}(2d)!$. Finally it remains to show that the sets $S$ lie in $\Gamma_t$ for some suitably small set $\Gamma_t$. Observe that the set $S \subseteq [\ell] \times [w]$ constructed in this argument has the following form: $|S \cap \{i\} \times [w]| \geq 1$ for all $i \in \ell$, and $|S| \leq t$. There are at most $(2w)^t$ sets of this form (ranging over all possible values of $\ell$), thus we can take $|\Gamma_t|= (2w)^t$.

It remains to estimate the value
\begin{gather*}
\Pr_{\rho \sim \mathcal{M}_{K}}[\rho \in \mathcal{B}] \leq \frac{1}{|\mathcal{M}_K|} \sum\limits_{t \geq \frac{d}{2}} |\mathcal{M}_{K-t}| |\Gamma_t| |\Delta_t| \\
=
\sum\limits_{t \geq \frac{d}{2}} \frac{|\mathcal{M}_{K-d}|}{|\mathcal{M}_K|}  |\Gamma_t| |\Delta_t|
\end{gather*}

Observe that for each $\kappa \in \mathcal{M}_{K-t}$ there are at least $\binom{N-K+t}{t}$ elements $\alpha \in \mathcal{M}_{K}$ such that $\alpha \subseteq \kappa$, and on the other hand for any $\alpha \in \mathcal{M}_{K}$ there are at most $\binom{K}{t}\binom{K+L}{t}t!$ elements $\kappa \in \mathcal{M}_{K-t}$ so that $\alpha \subseteq \kappa$. Thus
\[
\frac{ |\mathcal{M}_{K-t}|}{|\mathcal{M}_K|} \leq \frac{\binom{K}{t}\binom{K+L}{t}t!}{\binom{N-K+t}{t}} 
\]

So overall we have:
\begin{gather*}
\Pr_{\rho \sim \mathcal{M}_{K}}[\rho \in \mathcal{B}] \leq \sum\limits_{t \geq \frac{d}{2}}
\frac{\binom{K}{t}\binom{K+L}{t}t!}{\binom{N-K+t}{t}} |\Gamma_t| |\Delta_t|\\
\leq |\Delta|\sum\limits_{t \geq \frac{d}{2}}
\frac{\binom{K}{t}\binom{2K}{t}t!}{\binom{N-K+t}{t}} (2w)^t \\
\leq 
\binom{K}{2d}\binom{K+L}{2d}(2d)!\sum\limits_{t \geq \frac{d}{2}}
\frac{\binom{K}{t}\binom{2K}{t}t!}{\binom{N-K+t}{t}} (2w)^t \\
\leq (\frac{2K}{\sqrt{2d}})^{4d}\sum\limits_{t \geq \frac{d}{2}} \left(\frac{8wK^2}{N}\right)^t
\end{gather*}
Where we apply the assumptions $K \geq L$ and $K\leq \frac{N}{4}$. We may safely assume $\frac{wK^2}{N} = o(1)$, otherwise the conclusion of the theorem holds trivially. Thus by geometric decay of terms in the above sum we have, for some absolute $C \in \mathbb{N}$, the bounds:

\begin{gather*}
\Pr_{\rho \sim \mathcal{M}_{K}}[\rho \in \mathcal{B}] \leq C (\frac{2K}{\sqrt{2d}})^{4d}\left(\frac{8wK^2}{N}\right)^{\frac{d}{2}} \\
\leq (2K)^{4d} \left(\frac{8wK^2}{N}\right)^{\frac{d}{2}}
\end{gather*}
provided $K, d$ are sufficiently large with respect to $C$. So overall
\begin{gather*}
\Pr_{\rho \sim \mathcal{M}_{K}}[\rho \in \mathcal{B}] \leq \exp(4d \log(2K) + \frac{d}{2}\log(8wK^2) - \frac{d}{2} \log N) \\
\leq \exp(d(\log (16K^4) + \log (8wK) + \log N^{1/2}))
= \exp\Bigl(d (\log wK^5 -\log N^{1/2} + O(1))\Bigr)
\end{gather*}



\end{proof}

% \todo{should pose the two questions: (1) is LOP reducible to AVOID, (2) is AVOID $P^{NP}$ equivalent to a language? A positive answer to the first implies a positive answer to the second.}
% \todo{can we prove the following:}
% \begin{conjecture}
% If $\mathsf{NP} \subseteq \mathsf{P/poly}$ then every language in the polynomi_al hierarchy is reducible to LOP.
% \end{conjecture}
% \todo{I think we can cast Li's result in this way... should probably figure out connection to witnessing theorem, it would be nice to extend li's result to say that any $TF\Sigma_2^P$ search problem which is provably total in a certain fragment of buss hierarchy has this sort of non-trivial witnessing theorem. }

% \todo{if we could prove such a witnessing theorem it would be interesting to see if we could get a separation for some TFNP problem which is provably total in a higher fragment of bounded depth frege, i.e.}
% \begin{definition}[Generalized Linear Ordering Principle]
% The search problem GLOP is defined as follows. We are given a binary relation $\trianglelefteq$ on $[N] \times [N]$ and a set $A \subseteq [N]$, and must find one of the following solutions:
% \begin{enumerate}
%     \item A witness that $\trianglelefteq$ does not define a partial order on $A$.
%     \item An element $a_0 \in A$ so that $a_0 \trianglelefteq a$ for all $a \in A$.
% \end{enumerate}
% \end{definition}
% Observe that in this case the solution is not gaurunteed to be unique.

% Could we show anything like the following?
% \begin{conjecture}
% Let $R \in TF\Sigma_2^P$ (a first order principle) and let $R' \in TFNP$ be its skolemized form. If $R'$ is provable in $res(\log)$ then $R$ is reducible (in some sense) to LOP.
% \end{conjecture}
